\documentclass[10pt,a4paper]{amsart}
\usepackage[margin=19mm]{geometry}
\usepackage{amsmath,amssymb,mathtools}
\usepackage[scr=rsfs,cal=euler]{mathalfa}
\usepackage{xfrac}
\usepackage[unicode,hidelinks,bookmarks=false]{hyperref}
\newcommand{\ie}{i.e.\ }
\newcommand{\cf}{cf.\ }
\newcommand{\resp}{respectively\ }
\newcommand{\Subsection}{\subsection}
\providecommand{\nobreakdash}{\nobreak\hbox{-}}
\setlength{\emergencystretch}{2em}
\multlinegap0pt
\usepackage{amssymb}
\newtheorem{assumption}{Assumption}
\newtheorem{theorem}{Theorem}
\newtheorem{lemma}{Lemma}
\newtheorem{corollary}{Corollary}
\newtheorem{prop}{Proposition}
\newcommand\y{{{\langle y \rangle}}}
\title{Well-posedness of the two-dimensional unsteady Prandtl system in Sobolev space with degenerate critical points}
\author{Shi-Yong Zhu, Ya-Guang Wang}
\date{Source: arXiv:2609.10936v1 (CC BY 4.0)}
\begin{document}
\maketitle

\noindent\emph{Source and licence.} Well-posedness of the two-dimensional unsteady Prandtl system in Sobolev space with degenerate critical points. Version arXiv:2609.10936v1 (CC BY 4.0), distributed by arXiv under the Creative Commons Attribution 4.0 International licence, http://creativecommons.org/licenses/by/4.0/, as recorded in the arXiv licence metadata for this exact version and shown on the abstract page https://arxiv.org/abs/2609.10936v1. Full licence notice: \texttt{licenses/arxiv-2609.10936-CC-BY-4.0.txt}.\par\smallskip

\noindent\textit{Editorial scope.} Counted unit: the article's theorem Phe (source0.tex lines 5164--5184). The article's complete original proof of it is delivered with it (source0.tex lines 5185--6816, one \texttt{proof} environment of 1632 lines). No mathematics is written, repaired or re-derived: the statement and the proof are the article's own lines, in the article's own notation, and no retained line is substituted or re-flowed.  Retained uncounted as the source-local context the proof needs: lines 81--88; lines 90--92; lines 145--167; lines 172--178; lines 209--235; lines 238--244; lines 342--348; lines 483--495; lines 512--515; lines 581--583; lines 604--619; lines 613--616; lines 831--878; lines 884--903; lines 964--981; lines 1048--1055; lines 1083--1117; lines 1177--1195; lines 2808--2815; lines 2820--2826; lines 3617--3629; lines 3633--3642; lines 3644--3646; lines 3678--3680; lines 3690--3695; lines 3698--3707; lines 4046--4055; lines 4156--4163; lines 4165--4172; lines 4210--4244; lines 4671--4693; lines 4904--4918; lines 5128--5160.  Nothing was omitted from any retained span, so the package carries no omission record.  No retained line contains a citation, so the wrapper carries no bibliography and no bibliography entry.  No retained line contains a figure environment, an image-inclusion command or drawing code, so no image is delivered and the package declares no asset dependency.  Editorial formatting only, in addition: an AMS \texttt{amsart} preamble that restates the article's own declarations and macros used by the retained lines, namely the environments \texttt{assumption}, \texttt{theorem}, \texttt{lemma}, \texttt{corollary}, \texttt{prop} on separate counters, together with the standard packages those lines need.  The retained environments print in this document as Assumption 1, Theorem 1, Lemma 1, Corollary 1, Proposition 1 in order of appearance, whereas the article prints the same items with the numbering of its own full document.  The complete native source of the article as distributed in the arXiv e-print for this exact version is bundled unchanged as \texttt{sources/209968/source0.tex}.
\begin{equation}\label{Prandtl}
\begin{cases}
\partial_{t}u+u\partial_{x}u+v\partial_{y}u=\partial_{y}^{2}u-\partial_{x}P,\\
\partial_{x}u+\partial_{y}v=0,\\
u|_{y=0}=v|_{y=0}=0,~\lim\limits_{y\to+\infty} u=U(t,x),\\
u|_{t=0}=u_{in}(x,y)
\end{cases}
\end{equation}

\begin{equation}\label{B}
\partial_{t}U+U\partial_{x}U=-\partial_{x}P.
\end{equation}

\begin{assumption}\label{inA'}
Assume that there exist $y_\ast>0$, and $0<\check{y}<y_{*}<\hat{y}<+\infty$ such that
the initial data $u_{in}(x,y)$ satisfies
\begin{align}\label{inA1}
\partial_{y}u_{in}(x,y_{*})=\partial_{y}^{2}u_{in}(x,y_{*})=0,~\forall x\in\mathbb{T},
\end{align}
\begin{align}\label{inA2}
\partial_{y}^{3}u_{in}\geq\mathcal{C}>0~\text{in}~[\check{y},\hat{y}],
\end{align}
\begin{align}\label{inA3}
\partial_{y}u_{in}\geq c_{in}\varpi (y) e^{-\delta\y}, \quad \forall x\in  \mathbb{T}, ~y\ge 0,
\end{align}
where $\varpi(y)\in C^\infty[0,+\infty)$ is a nonnegative function satisfying
\begin{align}\label{pi}
\varpi(y)=
\begin{cases}
\mathcal{C}(y-y_{*})^{2},~\text{in}~[\check{y},\hat{y}],\\
1,~\text{in}~[0,+\infty)\setminus[\check{y},\hat{y}].
\end{cases}
\end{align}
with $\delta, c_{in}$ and $ \mathcal{C}$ positive constants.

\end{assumption}

\begin{align}\label{u0}
u_{0}(t,x,y)=
\begin{cases}
c_{0}\phi(y)+c_{0}ty,~&\text{in}~[0,T]\times\mathbb{T}\times[0,\hat{y}],\\
\phi(y)U(t,x),  ~&\text{in}~[0,T]\times\mathbb{T}\times[\hat{y}+1,+\infty),
\end{cases}
\end{align}

\begin{assumption}\label{inA}
For the initial data $u_{in}(x,y)$, let $\tilde{u}_{in}(x,y)=u_{in}(x,y)-u_{0}(0,x,y)$, with $u_0$ given in \eqref{u0}.
For some $s\in\mathbb{N}_{+}$, assume there exists positive constants $\varepsilon_{0}\in\left(0,\frac{1}{4}\right)$ and $\mathcal{C}_{in}$ such that
\begin{align}\label{inA4}
&\|\tilde{u}_{in}\|_{\hat{H}_{\psi_{in}}^{s}(\Omega)}
+\|\omega_{\varepsilon_{0}}^{in}\mathcal{P}_{in}\partial_{y}\partial_{x}^{s}\tilde{u}_{in}\|_{L_{\hat{\psi}_{in}}^{2}(\Omega)}
+\|\omega_{\varepsilon_{0}}^{in}\mathcal{P}_{in}^{2}\partial_{x}^{s}\tilde{u}_{in}\|_{L_{\hat{\psi}_{in}}^{2}(\Omega)}\\
&+\|\omega_{\varepsilon_{0}}^{in}\mathcal{P}_{in}\partial_{y}\partial_{x}^{s+1}\tilde{u}_{in}\|_{L_{\check{\psi}_{in}}^{2}(\Omega)}
+\|\omega_{\varepsilon_{0}}^{in}\mathcal{P}_{in}^{2}\partial_{x}^{s+1}\tilde{u}_{in}\|_{L_{\check{\psi}_{in}}^{2}(\Omega)}
\leq \mathcal{C}_{in},\nonumber
\end{align}
with $\psi_{in}=\psi(0,y)=e^{\frac{3}{4}\delta\y}$, $\hat{\psi}_{in}=\y^{-\frac{1}{2}} \psi_{in}$, $\check{\psi}_{in}=\y^{-1} \psi_{in}$,
and
\begin{align}\label{inA5}
|D^{\gamma}\tilde{u}_{in}|\leq \mathcal{C}_{in}\y^{-1}e^{-\delta\y}~\text{in}~\Omega,
\end{align}
for any $\gamma\in\Gamma_{9}$, with $\Gamma_{m}=\{(\gamma_{1},\gamma_{2})\in\mathbb{N}^{2}| \gamma_{1}+\gamma_{2}\leq m\}$  denoting a set of multiindex. Moreover, assume that
\begin{align}\label{inA6}
\begin{cases}
|\partial_{y}\rho_{in}|\leq \mathcal{C}_{in}\y^{-1}, |\partial_{x}\rho_{in}|\leq \mathcal{C}_{in},~\text{in}~\mathbb{T}\times[\hat{y},+\infty),\\
\rho_{in}\leq -\delta,~\text{in}~\mathbb{T}\times[Y,+\infty),
\end{cases}
\end{align}
for fixed constants $Y>\hat{y}$.


\end{assumption}

\begin{theorem}\label{MR}
Assume that the initial data $u_{in}$ of the problem \eqref{Prandtl} satisfy  Assumptions \ref{inA'} and \ref{inA} for  a fixed odd integer $s\geq 11$, and the compatibility conditions of the problem \eqref{Prandtl} up to order $s-1$. Moreover, the outflow
$U\in C([0,T];H^{s+2}(\mathbb{T}))\bigcap C^{1}([0,T];H^{s+1}(\mathbb{T}))\bigcap\limits_{i=2}^{\frac{s-3}{2}}C^{i}([0,T];H^{s-2i}(\mathbb{T}))$ is positive.
Then, there exists $t_{*}\in(0,T]$ such that the problem \eqref{Prandtl} admits a unique classical solution $u$ within $[0,t_{*}]$ satisfying that
$$w:=u-u_{0}\in L^{\infty}((0, t_{*}), \hat{H}_{\psi}^{s}(\Omega))\cap L^{2}((0, t_{*}), H_{\hat{\psi}}^{s}\cap H_{\check{\psi}}^{s+1}(\Omega)),$$
where $ \check{\psi}(t,y)=\y^{-1}\psi(t,y)$, and the function $u_0(t,x,y)$ is given in \eqref{u0}.
\end{theorem}

\begin{align}\label{theta}
\begin{cases}
\vartheta_{o}(y)=1,~ \forall y\in[0,Y],\\
\vartheta_{o}'\leq 0~\text{and}~|\vartheta_{o}^{(m)}|\leq c_{0}\vartheta_{c},~ \forall y\in [0,+\infty),\\
\vartheta_{o}\leq c_{1}e^{-\delta\y},~ \forall y\in [0,+\infty),
\end{cases}
\end{align}

\begin{equation}\label{app}
\begin{cases}
\partial_{t}u
+\mathfrak{U}\partial_{x}u+u\partial_{x}\mathfrak{U}
-\partial_{y}^{-1}[\partial_{x}u](\partial_{y}\mathfrak{U}+\varsigma)
-\partial_{y}^{-1}[\partial_{x}\mathfrak{U}]\partial_{y}u
-h\partial_{y}^{-1}[\partial_{x}u]\partial_{y}u\\
\quad
=\partial_{y}^{2}u+\eta\partial_{x}^{2}u+\kappa\partial_{x}^{2}\mathbb{U}+F,\\
u|_{y=0}=0,~\lim\limits_{y\to+\infty} u=0,\\
u|_{t=0}=\check{u}(x,y).
\end{cases}
\end{equation}

\begin{align}\label{rhoPd}
\rho(t,x,y)=\frac{\partial_{y}\dot{\mathfrak{U}}+\varsigma'}{\dot{\mathfrak{U}}+\varsigma}, \quad
\mathcal{P}(t,x,y)=\frac{\chi(y)}{\sqrt{\dot{\mathfrak{U}}+\varsigma}}+1-\chi(y),
\end{align}

\begin{equation}\label{varphi}
\varphi(t,y)=\y^{-\frac{1}{2}}\psi(t,y)\sqrt{\zeta_{\lambda}(t,y)},
\end{equation}

\begin{lemma}\label{weighte}
Let $\varphi$ be defined as in \eqref{varphi}, with $\lambda$ a even positive integer. If $\lambda$ is large enough, there exists positive constants $\Lambda_{\lambda}$ and $\iota_{\lambda}$ depending only on $\lambda$ such that for any
$\Lambda>\Lambda_{\lambda}, \iota\in(0,\iota_{\lambda})$ one can find a constant $T_{\lambda,\Lambda}\in(0,T)$ depending only on $\lambda$ and $\Lambda$ to obtain the following inequalities in $[0,T_{\lambda,\Lambda}]\times[0,+\infty)$,
\begin{align}\label{We}
\partial_{t}\varphi^{2}+\partial_{y}^{2}\varphi^{2}
\leq -\frac{\lambda}{32}\omega_{\lambda}\varphi^{2}
     -\frac{3}{2}\Lambda\delta\y\varphi^{2},
\end{align}
moreover,
\begin{align}\label{varphiy}
|\partial_{y}\varphi^{2}|
\leq C_{l}\varphi^{2}+C_{l}\lambda\sqrt{\omega_{\lambda}}\varphi^{2},
\end{align}
where $C_{l}$ is a positive constant independent of $\lambda,\iota$ and $\Lambda$, and
$$\omega_{\lambda}(t,y)=\frac{1}{\frac{\lambda^{-2}}{16}(y-y_{*})^{2}+\bar{\varsigma}+t}.$$
\end{lemma}

\begin{align}
|\partial_{y}\varphi^{2}|
\leq C_{l}\varphi^{2}+C_{l}\lambda\sqrt{\omega_{\lambda}}\varphi^{2},
\end{align}

\begin{assumption}\label{MA}
Let $\rho$ and $\mathcal{P}$ be given in \eqref{rhoPd} and $\mathbb{U}=\mathfrak{U}-u_{0}$. Assume that there exist $\hat{T}_{*}\in(0,T]$  such that the following inequalities hold,
$$\dot{\mathfrak{U}}\geq C\varpi(U-\mathfrak{U})\geq C\varpi e^{-(1+\theta t)\delta\y}~\text{in}~\Omega_{\hat{T}_{*}},$$
$$\dot{\mathfrak{U}}\geq C(\varpi+t)~\text{in}~\Omega_{\hat{T}_{*}}^{*},$$
$$|\partial_{x}\dot{\mathfrak{U}}|\leq C(1+t\y)\dot{\mathfrak{U}},
  |\partial_{x}^{2}\dot{\mathfrak{U}}|\leq C(1+t\y)\dot{\mathfrak{U}},
  |\partial_{y}\partial_{x}^{2}\dot{\mathfrak{U}}|\leq C(1+t\y)\dot{\mathfrak{U}}^{\frac{1}{2}}~\text{in}~\Omega_{\hat{T}_{*}},$$
$$|\vartheta_{c}^{2}\partial_{x}^{2}\mathbb{U}|\leq C_{*}\y^{-1}\dot{\mathfrak{U}}~\text{in}~\check{\Omega}_{\hat{T}_{*}}^{\hat{y}},$$
\begin{align*}
& \|\mathbb{U}\|_{L_{t}^{\infty}\mathcal{H}_{\psi,\varphi}^{s}(\Omega_{\hat{T}_{*}})}
 +\|\partial_{y}\mathbb{U}\|_{L_{t}^{2}\hat{H}_{\psi}^{s}(\Omega_{\hat{T}_{*}})}
 +\|\partial_{y}\mathbf{U}_{s}\|_{L_{t}^{2}L_{\varphi}^{2}(\Omega_{\hat{T}_{*}})}
 +\|\mathcal{P}\mathbf{U}_{s}\|_{L_{t}^{2}L_{\varphi}^{2}(\Omega_{\hat{T}_{*}})}\\
&+\|\mathbf{U}_{s+1}\|_{L_{t}^{2}L_{\hat{\varphi}}^{2}(\Omega_{\hat{T}_{*}})}
 +\|\mathcal{P}\mathbf{U}_{s+1}\|_{L_{t}^{2}L_{\hat{\varphi}}^{2}(\Omega_{\hat{T}_{*}})}
 +\epsilon^{-2}\|\sqrt{\eta}\partial_{x}\mathbf{U}_{s+1}\|_{L_{t}^{2}L_{\hat{\varphi}}^{2}(\Omega_{\hat{T}_{*}})}
\leq \mathcal{Z},
\end{align*}
and
$$\rho\leq -\frac{15}{16}\delta,~\text{in}~\check{\Omega}_{T_{*}}^{Y},$$
\begin{align*}
|\rho|\leq C\mathcal{P},~|\partial_{x}\rho|\leq C\mathcal{P},~
|\partial_{y}\rho|\leq C\mathcal{P}^{2},~
\left|\partial_{x}(\mathcal{P}\partial_{x}\rho)\right|\leq C(1+t\y)^{2}\mathcal{P}^{2},~\text{in}~\Omega_{T_{*}},
\end{align*}
where $\varpi$ and $\vartheta_{c}$ are given in \eqref{pi} and Section 3 respectively.
Moreover, the following constraints on $F$ hold,
\begin{align*}
|\partial_{y}F|\leq C_{*}\iota^{-1}(\dot{\mathfrak{U}}+\bar{\varsigma})~\text{in}~\Omega_{\hat{T}_{*}}~\text{and}~
\|\mathcal{P}(\partial_{y}^{2}F-\rho\partial_{y}F)\|_{L^{\infty}(\Omega_{\hat{T}_{*}})}\leq C_{*}\iota^{-1},
\end{align*}
$$
\|\partial_{t}F\|_{L_{t}^{\infty}H_{\psi}^{s-4}(\Omega_{\hat{T}_{*}})}+\|F\|_{L_{t}^{\infty}H_{\psi}^{s-2}(\Omega_{\hat{T}_{*}})}
\leq C_{*}\bar{\eta},
$$
$$
|D^{\gamma}F|\leq C_{*}e^{-\frac{4}{3}\delta\y}~\text{in}~\Omega_{\hat{T}_{*}}, \forall \gamma\in\Gamma_{2},
$$
and
\begin{align*}
\|\mathbf{F}_{s}\|_{L_{t}^{2}L_{\hat{\psi}}^{2}(\Omega_{\hat{T}_{*}})}^{2}
+\|F\|_{L_{t}^{2}\hat{H}_{\psi}^{s}(\Omega_{\hat{T}_{*}})}^{2}
+\sum_{i=0}^{\frac{s-1}{2}}\|\partial_{t}^{i}F\|_{L_{t}^{\infty}\mathring{H}^{s-1-2i}([0,\hat{T}_{*}]\times\mathbb{T})}^{2}
+\|\partial_{x}^{s}F\|_{L_{t}^{\infty}\mathring{H}^{0}([0,\hat{T}_{*}]\times\mathbb{T})}^{2}
\leq C_{*}\bar{\eta},
\end{align*}
where $\mathbf{F}_{s}=\mathcal{P}\left(\partial_{y}\partial_{x}^{s}F-\rho \partial_{x}^{s}F\right)$.
\end{assumption}

\begin{theorem}\label{wpae}
There exists a positive constant $\lambda_{*}$ such that for any $\lambda>\lambda_{*}$, there are positive constants $\Lambda_{*}, \iota^{*}$ and $\mathcal{Z}_{*}$ depending only on $\lambda$ for which, when $\Lambda\geq\Lambda_{*}$, $\theta>0$, $\iota\in(0,\iota^{*})$ and $\mathcal{Z}\geq\mathcal{Z}_{*}$, there is $t^{*}\in(0,\min\{T_{\lambda,\Lambda},\frac{1}{12\Lambda},T\}]$ independent of $\bar{\eta}$ such that if Assumption \ref{MA} holds for any $t_{v}\in(0,t^{*}]$, moreover, the initial date $\check{u}$ satisfies
$$\|\check{u}\|_{\hat{H}_{\psi_{in}}^{s}(\Omega)}
 +\|\mathcal{P}(0,x,y)(\partial_{y}\partial_{x}^{s}\check{u}-\rho(0,x,y)\partial_{x}^{s}\check{u})\|_{L_{\varphi_{in}}^{2}(\Omega)}
\leq \frac{(1-\epsilon_{0})\sqrt{\bar{\eta}}\mathcal{Z}}{8\epsilon_{0}},~\text{in}~\Omega,$$
and $\check{u}\mathbb{U}=0$ in $\Omega_{t_{v}}$, where $\varphi_{in}=\varphi(0,y)$, then the problem \eqref{app} admits a unique solution
$u\in L_{t}^{\infty}\mathcal{H}_{\psi,\varphi}^{s}(\Omega_{t_{v}})$ satisfying
\begin{align}\label{priore}
& \|u\|_{L_{t}^{\infty}\mathcal{H}_{\psi,\varphi}^{s}(\Omega_{t_{v}})}
 +\|u\|_{L_{t}^{2}\hat{H}_{\psi,\y}^{s}(\Omega_{t_{v}})}
 +\|\partial_{y}u\|_{L_{t}^{2}\hat{H}_{\psi}^{s}(\Omega_{t_{v}})}
 +\|\partial_{x}u\|_{L_{t}^{2}\hat{H}_{\psi,\eta}^{s}(\Omega_{t_{v}})}\\
&+\|\mathcal{P}\mathcal{U}_{s}\|_{L_{t}^{2}L_{\varphi}^{2}(\Omega_{t_{v}})}
 +\|\sqrt{\y}\mathcal{U}_{s}\|_{L_{t}^{2}L_{\varphi}^{2}(\Omega_{t_{v}})}
 +\|\partial_{y}\mathcal{U}_{s}\|_{L_{t}^{2}L_{\varphi}^{2}(\Omega_{t_{v}})}
 +\|\sqrt{\eta}\partial_{x}\mathcal{U}_{s}\|_{L_{t}^{2}L_{\varphi}^{2}(\Omega_{t_{v}})}
\leq \frac{1-\epsilon_{0}}{\epsilon_{0}}\sqrt{\bar{\eta}}\mathcal{Z}.\nonumber
\end{align}
Here $T_{\lambda,\Lambda}$ is given in Lemma \ref{weighte}.
\end{theorem}

\begin{lemma}\label{nc}
Assume $\psi_{c}(y)\in C^{\infty}(\mathbb{R}_{+})$ to be a positive nondecreasing function,
and $\mathfrak{g}(y)\in C^{1}(\mathbb{R}_{+})\cap L_{\psi_{c}}^{2}(\mathbb{R}_{+})$ to be a positive function.
Let $\Phi(y)=\psi_{c}(y)\omega_{c}(y)\left[\frac{\chi(y)}{\sqrt{\mathfrak{g}(y)}}+1-\chi(y)\right]^{2}$, with
$$\omega_{c}(y)=\frac{\tilde{\chi}(y)}{((y-y_{*})^{2}+\mathfrak{c})^{\frac{\varepsilon_{0}}{2}}}+1-\tilde{\chi}(y),$$
where $\epsilon_{0}$ is given as in Assumption \ref{inA}, $\mathfrak{c}$ is a positive constant, $\chi(y)$ is given in Section 2, and $\tilde{\chi}(y)\in C^{\infty}(\mathbb{R}_{+})$ is a function such that
$0\leq\tilde{\chi}(y)\leq 1$.
For any $f(y)\in L_{\psi_{c}}^{2}(\mathbb{R}_{+})$ satisfying $f(0)=0$, define $$\dot{f}(y)=f'(y)-\varrho(y)f(y),$$
with $\varrho(y)=\frac{\mathfrak{g}'(y)}{\mathfrak{g}(y)}$.
Assume there exist a constant $\ell\in(0,1)$ and negative constants $\check{\mathfrak{c}}$ and $\hat{\mathfrak{c}}$ such that
 $\check{\mathfrak{c}}\leq\varrho\leq -\ell^{-1}\frac{\psi_{c}'(y)}{\psi_{c}(y)}\leq\hat{\mathfrak{c}}$ for any $y\in[Y,+\infty)$,
and $\dot{f}(y)\in L_{\Phi}^{2}(\mathbb{R}_{+})$,
then for any nonnegative function $\xi(y)\in C^{1}[0,+\infty)$ such that $\xi'(y)\leq0$,
there exist a positive constant $\mathcal{N}$ depending on the values of $\mathfrak{g}$ on $[0,Y]\setminus[\check{y}_{*},\hat{y}_{*}]$, $\psi_{c}$ and $\ell$ but independent of $f$ and $\mathfrak{c}$ such that
\begin{align}\label{nce0}
\|\xi f\|_{L_{\Phi}^{2}(\mathbb{R}_{+})}\leq \mathcal{N}\|\xi\dot{f}\|_{L_{\Phi}^{2}(\mathbb{R}_{+})}.
\end{align}
\end{lemma}

\begin{corollary}\label{ncC1}
Assume $\tilde{\xi}(y)\in C^{1}[0,+\infty)$ to be a nonnegative bounded function. Let $\mathcal{N}$ be given in Lemma \ref{nc}.
Under the assumption of Lemma \ref{nc}, if $|\tilde{\xi}'(y)|\leq \frac{1}{2\mathcal{N}}\tilde{\xi}(y)$ for any $y\in[0,+\infty)$, then
\begin{align*}
\|\tilde{\xi} f\|_{L_{\Phi}^{2}(\mathbb{R}_{+})}\leq 2\mathcal{N}\|\tilde{\xi}\dot{f}\|_{L_{\Phi}^{2}(\mathbb{R}_{+})}.
\end{align*}

\end{corollary}

\begin{corollary}\label{ncC2}
Assume there exist $T_{c}\in(0,T)$ such that
$$\dot{\mathfrak{U}}(t,x,\cdot)\in C^{1}(\mathbb{R}_{+})\cap L_{\psi}^{2}(\mathbb{R}_{+}),~\forall (t,x)\in[0,T_{c}]\times\mathbb{T},$$
and
$$\mathcal{P}^{-1}\leq C,~|\rho|\leq C\mathcal{P},~|\partial_{x}\rho|\leq C\mathcal{P},~|\partial_{y}\rho|\leq C\mathcal{P}^{2},~\text{in}~ \Omega_{T_{c}},$$
then for any positive integer $r$ such that $r\leq s+1$ and any nonnegative function $\xi(y)\in C^{1}[0,+\infty)$ such that $\xi'(y)\leq0$, the following norm inequalities hold for any $t\in[0,T_{c}]$,
\begin{align}\label{ncc21}
\|\xi\partial_{x}^{r}f\|_{L_{\Psi}^{2}(\Omega)}
\leq C\|\xi\mathcal{P}\mathcal{F}_{r}\|_{L_{\varphi_{o}}^{2}(\Omega)},~
\|\xi\mathcal{P}\partial_{y}\partial_{x}^{r}f\|_{L_{\varphi_{o}}^{2}(\Omega)}
\leq C\|\xi\mathcal{F}_{r}\|_{L_{\varphi_{o}}^{2}(\Omega)}
    +C\|\xi\mathcal{P}\mathcal{F}_{r}\|_{L_{\varphi_{o}}^{2}(\Omega)},
\end{align}
and
\begin{align}\label{ncc22}
\|\xi\partial_{y}^{2}\partial_{x}^{r}f\|_{L_{\varphi_{o}}^{2}(\Omega)}
\leq C\|\xi\mathcal{F}_{r}\|_{L_{\varphi_{o}}^{2}(\Omega)}
    +C\|\xi\partial_{y}\mathcal{F}_{r}\|_{L_{\varphi_{o}}^{2}(\Omega)}
    +C\|\xi\mathcal{P}\mathcal{F}_{r}\|_{L_{\varphi_{o}}^{2}(\Omega)},
\end{align}
moreover,
\begin{align}\label{ncc23}
\|\xi\partial_{x}^{r}f\|_{L_{\Psi_{c}}^{2}(\Omega)}
\leq C\|\xi\mathcal{P}\mathcal{F}_{r}\|_{L_{\psi}^{2}(\Omega)},~
\|\xi\mathcal{P}\partial_{y}\partial_{x}^{r}f\|_{L_{\Psi_{c}}^{2}(\Omega)}
\leq C\|\xi\mathcal{F}_{r}\|_{L_{\psi}^{2}(\Omega)}
    +C\|\xi\mathcal{P}\mathcal{F}_{r}\|_{L_{\psi}^{2}(\Omega)},
\end{align}
\begin{align}\label{ncc24}
\|\xi\partial_{y}^{2}\partial_{x}^{r}f\|_{L_{\Psi_{c}}^{2}(\Omega)}
\leq C\|\xi\mathcal{F}_{r}\|_{L_{\psi}^{2}(\Omega)}
    +C\|\xi\partial_{y}\mathcal{F}_{r}\|_{L_{\psi}^{2}(\Omega)}
    +C\|\xi\mathcal{P}\mathcal{F}_{r}\|_{L_{\psi}^{2}(\Omega)}.
\end{align}
\end{corollary}

\begin{corollary}\label{ncC3}
Assume there exist $T_{c}\in(0,T)$ such that
$$\dot{\mathfrak{U}}(t,x,\cdot)\in C^{1}(\mathbb{R}_{+})\cap L_{\psi}^{2}(\mathbb{R}_{+}),~\forall (t,x)\in[0,T_{c}]\times\mathbb{T},$$
and
$$\mathcal{P}^{-1}\leq C,~\partial_{x}\mathcal{P}^{-1}\leq C,~|\rho|\leq C\mathcal{P},~|\partial_{x}\rho|\leq C\mathcal{P},~\text{in}~ \Omega_{T_{c}}.$$
Assume that $\xi(y)\in C^{1}[0,+\infty)$ is a nonnegative function such that $\xi'(y)\leq0$, then one has the following norm inequalities for any $t\in[0,T_{c}]$,
\begin{align}\label{ncc3}
\|\xi\partial_{x}^{r+1}f\|_{L_{\Psi}^{2}(\Omega)}
\leq C\|\mathcal{P}\|_{L^{\infty}(\Omega)}\|\xi\partial_{x}\mathcal{F}_{r}\|_{L_{\varphi_{o}}^{2}(\Omega)}
    +C\|\mathcal{P}\|_{L^{\infty}(\Omega)}\|\xi\mathcal{P}\mathcal{F}_{r}\|_{L_{\varphi_{o}}^{2}(\Omega)},
\end{align}
and
\begin{align}\label{ncc3'}
\|\xi\partial_{x}^{r+1}f\|_{L_{\Psi_{c}}^{2}(\Omega)}
\leq C\|\mathcal{P}\omega_{o}^{-1}\|_{L^{\infty}(\Omega)}\|\xi\partial_{x}\mathcal{F}_{r}\|_{L_{\varphi_{o}}^{2}(\Omega)}
    +C\|\mathcal{P}\omega_{o}^{-1}\|_{L^{\infty}(\Omega)}\|\xi\mathcal{P}\mathcal{F}_{r}\|_{L_{\varphi_{o}}^{2}(\Omega)},
\end{align}

\end{corollary}

\begin{align}\label{I32}
&2\left(\frac{\partial_{y}\mathcal{P}}{\mathcal{P}}\right)^{2}
+2\eta\left(\frac{\partial_{x}\mathcal{P}}{\mathcal{P}}\right)^{2}
-\frac{3\chi}{4\mathcal{P}\sqrt{\dot{\mathfrak{U}}+\varsigma}}\left[\rho^{2}+\eta\left(\frac{\partial_{x}\dot{\mathfrak{U}}}{\dot{\mathfrak{U}}+\varsigma}\right)^{2}\right]\\
\leq&-\frac{\chi}{4\mathcal{P}\sqrt{\dot{\mathfrak{U}}+\varsigma}}\left[\rho^{2}+\eta\left(\frac{\partial_{x}\dot{\mathfrak{U}}}{\dot{\mathfrak{U}}+\varsigma}\right)^{2}\right]
     +C_{*}
\leq C_{*}.\nonumber
\end{align}

\begin{align}\label{I33}
&\frac{\partial_{y}^{-1}\left[\partial_{x}\dot{\mathfrak{U}}\right]\varsigma'+\varsigma''-\partial_{y}\eta\partial_{x}^{2}\mathbb{U}+\eta\partial_{x}^{2}\partial_{y}u_{0}+\partial_{y}F+\partial_{y}[\kappa\partial_{x}^{2}\mathbb{U}]}{2(\dot{\mathfrak{U}}+\varsigma)}\\
&-(1-\chi)\frac{\partial_{y}^{-1}\left[\partial_{x}\mathfrak{U}\right]\varsigma'+\varsigma''-\partial_{y}\eta\partial_{x}^{2}\mathbb{U}+\eta\partial_{x}^{2}\partial_{y}u_{0}+\partial_{y}F+\partial_{y}[\kappa\partial_{x}^{2}\mathbb{U}]}{2\mathcal{P}(\dot{\mathfrak{U}}+\varsigma)}\nonumber\\
&\frac{-\partial_{y}^{-1}[\partial_{x}\mathfrak{U}]\chi'+\chi'\rho-\chi''}{\mathcal{P}\sqrt{\dot{\mathfrak{U}}+\varsigma}}
+\frac{\partial_{y}^{-1}\left[\partial_{x}\mathfrak{U}\right]\chi'+\chi''}{\mathcal{P}}
\leq C_{*}\iota^{-1}.\nonumber
\end{align}

\begin{equation}\label{app'}
\begin{cases}
\partial_{t}u
+\tilde{\mathfrak{U}}\partial_{x}u+u\partial_{x}\tilde{\mathfrak{U}}
-\partial_{y}^{-1}[\partial_{x}u](\partial_{y}\tilde{\mathfrak{U}}+\tilde{\varsigma})
-\partial_{y}^{-1}[\partial_{x}\tilde{\mathfrak{U}}]\partial_{y}u
-h\partial_{y}^{-1}[\partial_{x}u]\partial_{y}u\\
\quad
=\partial_{y}^{2}u+\eta\partial_{x}^{2}u+\kappa\partial_{x}^{2}\tilde{\mathbb{U}}+F,\\
u|_{y=0}=0,~\lim\limits_{y\to+\infty} u=0,\\
u|_{t=0}=\check{u}(x,y).
\end{cases}
\end{equation}

\begin{equation}\label{App+1}
\begin{cases}
\partial_{t}\mathfrak{U}
+\mathfrak{U}\partial_{x}\mathfrak{U}
-\partial_{y}^{-1}[\partial_{x}\mathfrak{U}]\partial_{y}\mathfrak{U}
=\partial_{y}^{2}\mathfrak{U}+\eta\partial_{x}^{2}\mathbb{U}-\partial_{x}P-\hat{F},\\
\mathfrak{U}|_{y=0}=0,~\lim\limits_{y\to+\infty}\mathfrak{U}=U(t,x),\\
\mathfrak{U}|_{t=0}=u_{in}(x,y).
\end{cases}
\end{equation}

\begin{align}\label{hatF}
\hat{F}=-\partial_{y}^{-1}[\partial_{x}u]\tilde{\varsigma}-u\partial_{x}u+(1-h)\partial_{y}^{-1}[\partial_{x}u]\partial_{y}u,
\end{align}

\begin{align}\label{us-3}
\|u\|_{H_{\psi}^{s-3}(\Omega)}\leq C\check{c}+C_{*}\sqrt{\bar{\eta}}t,~\forall t\in[0,t^{*}].
\end{align}

\begin{assumption}\label{FA}
Let $\hat{F}$ be given  in \eqref{hatF}. For some $\hat{t}_{*}\in(0,T]$ and any multiindex $\gamma=(\gamma_{1},\gamma_{2})\in\Gamma_{9}$, assume $\hat{F}$ satisfies
\begin{align}\label{Fd}
|D^{\gamma}\hat{F}|\leq C_{*}e^{-\frac{4}{3}\delta\y},~\text{in}~\Omega_{\hat{t}^{*}}.
\end{align}
\end{assumption}

\begin{lemma}\label{DI}
Assume the initial date $\tilde{u}_{in}$ in the problem \eqref{App+1} satisfying Assumption \ref{inA}.
There exists a positive constant $\theta_{*}$ such that for any $\theta\geq\theta_{*}$, there corresponds a constant $\bar{t}_{0}\in(0,t^{*}]$ such that for any multiindex $\gamma\in\Gamma_{6}$ and any $t_{v}\in(0,\bar{t}_{0}]$,
if $\mathfrak{U}$ is a solution to the problem \eqref{App+1} on the time interval $[0,t_{v}]$ such that
$\|\mathbb{U}\|_{L_{t}^{\infty}\hat{H}_{\psi}^{s}(\Omega_{t_{v}})}\leq C_{*}$,
and Assumption \ref{FA} holds for $t_{v}$, then
\begin{align}\label{Di}
|D^{\gamma}\mathbb{U}|\leq C_{*}\y^{-1}e^{-(1-\theta t)\delta\y}, ~\text{in}~\Omega_{t_{v}}.
\end{align}
\end{lemma}

\begin{lemma}\label{wdotu}
Let $\theta_{*}$ and $\bar{t}_{0}$ be obtained in Lemma \ref{DI}.
Assume the initial date $\tilde{u}_{in}$ in the problem \eqref{App+1} satisfies Assumption \ref{inA'} and \ref{inA}.
There exists a positive constant $\theta^{*}\geq \theta_{*}$ and positive constants $K_{1},K_{2}$ independent of $\mathcal{Z}$ such that for any $\theta\geq\theta^{*}$, there is a corresponding $t_{0}\in(0,\bar{t}_{0}]$ such that for any $t_{v}\in[0,t_{0}]$, if $\mathfrak{U}$ is a solution to the problem \eqref{App+1} on the time interval $[0,t_{v}]$ such that
$\|\mathbb{U}\|_{L_{t}^{\infty}\hat{H}_{\psi}^{s}(\Omega_{t_{v}})}\leq C_{*}$,
and Assumption \ref{FA} holds for $t_{v}$, then the following inequality holds in $\Omega_{t_{v}}$,
\begin{align}\label{dotu}
K_{1}\varpi e^{-(1+\theta t)\delta\y}\leq \dot{\mathfrak{U}}\leq K_{2}e^{-(1-\theta t)\delta\y}.
\end{align}
\end{lemma}

\begin{corollary}\label{wdotuC'}
Under the assumption of Lemma \ref{wdotu}, for any multiindex $\gamma\in\Gamma_{6}$ and $t_{v}\in[0,t_{0}]$, one has
\begin{align*}
|\vartheta_{c}^{2}D^{\gamma}\mathbb{U}|
\leq C_{*}\y^{-1}e^{-(1+\theta t)\delta\y}
\leq C_{*}\y^{-1}\dot{\mathfrak{U}},~\text{in}~\check{\Omega}_{t_{v}}^{\hat{y}}.
\end{align*}
\end{corollary}

\begin{corollary}\label{wdotuC}
Under the assumption of Lemma \ref{wdotu}, for any $t_{v}\in[0,t_{0}]$ one has
\begin{align}\label{Umu}
\frac{K_{1}}{(1+\theta t)\delta}e^{-(1+\theta t)\delta\hat{y}}e^{-(1+\theta t)\delta\y}
\leq U-\mathfrak{U}
\leq \frac{K_{2}}{(1-\theta t)\delta}e^{-(1-\theta t)\delta\y}, ~\text{in}~\Omega_{t_{v}}.
\end{align}
\end{corollary}

\begin{lemma}\label{sdotu}
Let $\Gamma_{*}=\{(m,k)|m,k\in\mathbb{N}, 0\leq m\leq 2, 0\leq m+k\leq3\}$.
Let $\theta^{*}$ and $t_{0}$ be obtained in Lemma \ref{wdotu}.
Assume the initial date $\tilde{u}_{in}$ in the problem \eqref{App+1} satisfies Assumption \ref{inA'} and \ref{inA}.
There exists positive constants $K_{mk}$ independent of $\mathcal{Z}$ such that for any $\theta\geq\theta^{*}$, there is a corresponding $t_{1}\in(0,t_{0}]$ such that for any $t_{v}\in[0,t_{1}]$, if $\mathfrak{U}$ is a solution to the problem \eqref{App+1} on the time interval $[0,t_{v}]$ such that
$\|\mathbb{U}\|_{L_{t}^{\infty}\hat{H}_{\psi}^{s}(\Omega_{t_{v}})}\leq C_{*}$,
and Assumption \ref{FA} holds for $t_{v}$, then the following inequalities hold in $\Omega_{t_{v}}$,
\begin{align}\label{dotul}
|\dot{\mathfrak{U}}|\leq K_{00}(U-\mathfrak{U}),~
|\partial_{y}\dot{\mathfrak{U}}|\leq K_{01}(U-\mathfrak{U}),~
|\partial_{y}^{2}\dot{\mathfrak{U}}|\leq K_{02}(U-\mathfrak{U}),~
|\partial_{y}^{3}\dot{\mathfrak{U}}|\leq K_{03}(U-\mathfrak{U}),
\end{align}
\begin{align}\label{dotux1}
|\partial_{x}\dot{\mathfrak{U}}|\leq K_{10}(1+\y t)(U-\mathfrak{U}),~
|\partial_{x}\partial_{y}\dot{\mathfrak{U}}|\leq K_{11}(1+\y t)(U-\mathfrak{U}),~
|\partial_{x}\partial_{y}^{2}\dot{\mathfrak{U}}|\leq K_{12}(1+\y t)(U-\mathfrak{U})
\end{align}
\begin{align}\label{dotux2}
|\partial_{x}^{2}\dot{\mathfrak{U}}|\leq K_{20}(1+\y t)(U-\mathfrak{U}),~
|\partial_{x}^{2}\partial_{y}\dot{\mathfrak{U}}|\leq K_{21}(1+\y t)(U-\mathfrak{U}),
\end{align}
in addition, there exists a positive constant $\bar{\iota}$ such that if
\begin{align}\label{Fe}
|\partial_{y}\hat{F}|\leq \bar{\iota}+C|y-y_{*}|+C_{*}t~\text{in}~\Omega_{t_{v}}^{*},
\end{align}
then there are positive constants $k_{0},k_{1},k_{2},k_{3}$ and $k_{4}$ independent of $\mathcal{Z}$ such that
\begin{align}\label{dotug}
\dot{\mathfrak{U}}\geq k_{0}\varpi(U-\mathfrak{U}),~\text{in}~\Omega_{t_{v}},
\end{align}
and
\begin{align}\label{dotub}
k_{1}\varpi+k_{2}t\leq \dot{\mathfrak{U}}\leq k_{3}\varpi+k_{4}t,~\text{in}~\hat{\Omega}_{t_{v}}\cap\check{\Omega}_{t_{v}}.
\end{align}
\end{lemma}

\begin{lemma}\label{rhoP}
Let $\theta^{*}$ be obtained in Lemma \ref{wdotu} and $t_{1}$ be obtained in Lemma \ref{sdotu}.
Assume the initial date $\tilde{u}_{in}$ in the problem \eqref{App+1} satisfies Assumption \ref{inA'} and \ref{inA}.
There exists positive constant $K$ independent of $\mathcal{Z}$ such that for any $\theta\geq\theta^{*}$, there is a corresponding $t_{2}\in(0,t_{1}]$ such that for any $t_{v}\in[0,t_{2}]$, if $\mathfrak{U}$ is a solution to the problem \eqref{App+1} on the time interval $[0,t_{v}]$ such that
$\|\mathbb{U}\|_{L_{t}^{\infty}\hat{H}_{\psi}^{s}(\Omega_{t_{v}})}\leq C_{*}$,
and Assumption \ref{FA} holds for $t_{v}$, then the following inequalities hold in $\Omega_{t_{v}}$,
\begin{align}\label{rhop1}
|\rho|\leq K\mathcal{P},~
|\partial_{y}\rho|\leq K\mathcal{P}^{2},~
|\partial_{x}\rho|\leq K(1+\y t)\mathcal{P},~
|\partial_{x}^{2}\rho|\leq K(1+\y t)^{2}\mathcal{P},~
~\text{in}~\Omega_{t_{v}},
\end{align}
\begin{align}\label{rhop2}
|\partial_{x}\dot{\mathfrak{U}}|\leq K(1+\y t)\dot{\mathfrak{U}},~
|\partial_{x}^{2}\dot{\mathfrak{U}}|\leq C(1+\y t)\dot{\mathfrak{U}},~
|\partial_{x}^{2}\partial_{y}\dot{\mathfrak{U}}|\leq C(1+\y t)\dot{\mathfrak{U}}^{\frac{1}{2}},~\text{in}~\Omega_{t_{v}}.
\end{align}
\begin{align}\label{rhop3}
|\partial_{x}\mathcal{P}|\leq K\mathcal{P},~
\left|\partial_{x}(\mathcal{P}\partial_{x}\rho)\right|\leq K(1+\y t)^{2}\mathcal{P}^{2}~\text{in}~\Omega_{t_{v}}.
\end{align}
\end{lemma}

\begin{lemma}\label{rhoP'}
Let $\theta^{*}$ be obtained in Lemma \ref{wdotu} and $t_{2}$ be obtained in Lemma \ref{rhoP}.
Assume the initial date $\tilde{u}_{in}$ in the problem \eqref{App+1} satisfies \eqref{inA1}--\eqref{inA3}, and Assumption \ref{inA}.
There exists positive constant $K$ independent of $C$ such that for any $\theta\geq\theta^{*}$, there is a corresponding $t_{3}\in(0,t_{2}]$ such that for any $t_{v}\in[0,t_{3}]$, if $\mathfrak{U}$ is a solution to the problem \eqref{App+1} on the time interval $[0,t_{v}]$ such that
$\|\mathbb{U}\|_{L_{t}^{\infty}\hat{H}_{\psi}^{s}(\Omega_{t_{v}})}\leq C_{*}$,
and Assumption \ref{FA} holds for $t_{v}$, then the following inequalities hold in $\Omega_{t_{v}}$,
\begin{align}\label{rhox}
|\partial_{x}\rho|\leq C\mathcal{P},~\text{in}~\Omega_{t_{v}},
\end{align}
and
\begin{align}\label{rhol}
\rho\leq -\frac{15}{16}\delta,~\text{in}~\check{\Omega}_{t_{v}}^{Y}.
\end{align}

\end{lemma}

\begin{assumption}\label{peA}
For some $\hat{t}_{*}\in(0,T]$, assume that
Assumption \ref{FA} holds for $\hat{t}_{*}$ such that  Lemma \ref{DI}--Lemma \ref{rhoP'} are valid on $[0,\hat{t}_{*}]$. Assume $u$ is a solution to the problem \eqref{app'} on the time interval $[0,\hat{t}_{*}]$ such that
\begin{align*}
 \|u\|_{L_{t}^{\infty}\mathcal{H}_{\psi,\tilde{\varphi}}^{s}(\Omega_{\hat{t}_{*}})}
+\|\partial_{y}\mathcal{U}_{s}\|_{L_{t}^{2}L_{\tilde{\varphi}}^{2}(\Omega_{\hat{t}_{*}})}
+\|\sqrt{\eta}\partial_{x}\mathcal{U}_{s}\|_{L_{t}^{2}L_{\tilde{\varphi}}^{2}(\Omega_{\hat{t}_{*}})}
+\|\tilde{\mathcal{P}}\mathcal{U}_{s}\|_{L_{t}^{2}L_{\tilde{\varphi}}^{2}(\Omega_{\hat{t}_{*}})}
\leq C_{*}\sqrt{\bar{\eta}},
\end{align*}
and $\mathfrak{U}$ is a solution to the problem \eqref{App+1} on the time interval $[0,\hat{t}_{*}]$ with $\hat{F}$ satisfying
$$|\partial_{y}\hat{F}|\leq C_{*}\iota^{-1}(\dot{\mathfrak{U}}+\bar{\varsigma})~\text{in}~\Omega_{\hat{t}_{*}},$$
and
$$\|\mathcal{P}(\partial_{y}^{2}\hat{F}-\rho\partial_{y}\hat{F})\|_{L^{\infty}(\Omega_{\hat{t}_{*}})}\leq C_{*},$$
such that $\mathbb{U}=\mathfrak{U}-u_{0}$ satisfies
$$\|\mathbb{U}\|_{L_{t}^{\infty}\hat{H}_{\psi}^{s}(\Omega_{\hat{t}_{*}})}\leq \mathcal{Z},
~\|\mathbb{U}\|_{L_{t}^{2}\hat{H}_{\psi,\y}^{s}(\Omega_{\hat{t}_{*}})}\leq \mathcal{Z}.$$
Moreover, assume that
\begin{align*}
 \|\sqrt{\tilde{\eta}}\partial_{x}\tilde{\mathbf{U}}_{s}\|_{L_{t}^{2}L_{\tilde{\varphi}}^{2}(\Omega_{\hat{t}_{*}})}
+\|\tilde{\mathcal{P}}\tilde{\mathbf{U}}_{s}\|_{L_{t}^{2}L_{\tilde{\varphi}}^{2}(\Omega_{\hat{t}_{*}})}
+\|\tilde{\mathbf{U}}_{s}\|_{L_{t}^{\infty}L_{\tilde{\varphi}}^{2}(\Omega_{\hat{t}_{*}})}
\leq \mathcal{Z},
\end{align*}
and
\begin{align*}
 \|\sqrt{\tilde{\eta}}\partial_{x}\tilde{\mathbf{U}}_{s+1}\|_{L_{t}^{2}L_{\hat{\tilde{\varphi}}}^{2}(\Omega_{\hat{t}_{*}})}
+\|\tilde{\mathcal{P}}\tilde{\mathbf{U}}_{s+1}\|_{L_{t}^{2}L_{\hat{\tilde{\varphi}}}^{2}(\Omega_{\hat{t}_{*}})}
+\|\tilde{\mathbf{U}}_{s+1}\|_{L_{t}^{\infty}L_{\hat{\tilde{\varphi}}}^{2}(\Omega_{\hat{t}_{*}})}
\leq \mathcal{Z},
\end{align*}
where $\hat{\tilde{\varphi}}=\y^{-\frac{1}{2}}\tilde{\varphi}$.
\end{assumption}

\begin{theorem}\label{Phe}
Let $\theta^{*}$ be given in Lemma \ref{wdotu} and $t_{3}$ be given in Lemma \ref{rhoP'}. Assume that the assumption of Theorem \ref{MR} holds.
There exists a positive constant $\lambda^{*}$ such that for any $\lambda>\lambda^{*}$, there are positive constants $\Lambda^{*}, \iota_{*}$ and $\mathcal{Z}^{*}$ depending only on $\lambda$ for which, when $\Lambda\geq\Lambda^{*}$, $\theta\geq\theta^{*}$, $\iota\in(0,\iota_{*})$ and $\mathcal{Z}\geq\mathcal{Z}^{*}$, there exists $t_{*}\in(0,t_{3}]$ such that if Assumption \ref{peA} holds for $t_{*}$, then we have
\begin{align}\label{Hoe1}
&\|\mathbf{U}_{s}\|_{L_{t}^{\infty}L_{\varphi}^{2}(\Omega_{t_{*}})}
 +\Lambda\delta\|\sqrt{\y}\mathbf{U}_{s}\|_{L_{t}^{2}L_{\varphi}^{2}(\Omega_{t_{*}})}
 +\|\mathcal{P}\mathbf{U}_{s}\|_{L_{t}^{2}L_{\varphi}^{2}(\Omega_{t_{*}})}\\
&+\|\partial_{y}\mathbf{U}_{s}\|_{L_{t}^{2}L_{\varphi}^{2}(\Omega_{t_{*}})}
 +\|\sqrt{\eta}\partial_{x}\mathbf{U}_{s}\|_{L_{t}^{2}L_{\varphi}^{2}(\Omega_{t_{*}})}
\leq \mathcal{Z},\nonumber
\end{align}
and
\begin{align}\label{Hoe2}
&\|\mathbf{U}_{s+1}\|_{L_{t}^{\infty}L_{\hat{\varphi}}^{2}(\Omega_{t_{*}})}
 +\Lambda\delta\|\sqrt{\y}\mathbf{U}_{s+1}\|_{L_{t}^{2}L_{\hat{\varphi}}^{2}(\Omega_{t_{*}})}
 +\|\mathcal{P}\mathbf{U}_{s+1}\|_{L_{t}^{2}L_{\hat{\varphi}}^{2}(\Omega_{t_{*}})}\\
&+\|\partial_{y}\mathbf{U}_{s+1}\|_{L_{t}^{2}L_{\hat{\varphi}}^{2}(\Omega_{t_{*}})}
 +\|\sqrt{\eta}\partial_{x}\mathbf{U}_{s+1}\|_{L_{t}^{2}L_{\hat{\varphi}}^{2}(\Omega_{t_{*}})}
\leq \mathcal{Z}.\nonumber
\end{align}
\end{theorem}

\begin{proof}[Proof of Theorem \ref{Phe}]

We shall prove that the inequalities \eqref{Hoe1} and \eqref{Hoe2} holds for some $t_{*}\in (0,t_{3}]$, under the assumption that Assumption \ref{peA} holds for $t_{3}$, with $\lambda,\Lambda, \theta, \mathcal{Z}$ sufficiently large and $\iota$ sufficiently small. As can be seen from our proof, the inequalities \eqref{Hoe1} and \eqref{Hoe2} still hold for $t_{*}$ when Assumption \ref{peA} holds for $t_{*}$.

Let $r=s$ or $s+1$. From \eqref{App+1} one knows that $\mathbf{U}_{r}(t,x,y)$ obeys the following problem
\begin{equation}\label{Un}
\begin{cases}
\partial_{t}\mathbf{U}_{r}
+\mathbf{L}_{1}\partial_{x}\mathbf{U}_{r}
+\mathbf{L}_{2}\partial_{y}\mathbf{U}_{r}
+\mathbf{L}_{0}\mathbf{U}_{r}
-\partial_{y}^{2}\mathbf{U}_{r}
-\eta\partial_{x}^{2}\mathbf{U}_{r}
=\sum\limits_{i=1}^{6}\mathcal{T}_{i},\\
\left(\partial_{y}\mathbf{U}_{r}+\rho\mathbf{U}_{r}\right)|_{y=0}=\partial_{x}^{r+1}P,~
\lim\limits_{y\to+\infty} \mathbf{U}_{r}=0,\\
\mathbf{U}_{r}|_{t=0}=\mathcal{P}(0,x,y)(\partial_{y}\partial_{x}^{s}\tilde{u}_{in}-\rho(0,x,y)\partial_{x}^{s}\tilde{u}_{in}),
\end{cases}
\end{equation}
where
\begin{align*}
\mathbf{L}_{0}
=&~r\partial_{x}\mathfrak{U}-2\partial_{y}\rho
  -2\left(\frac{\partial_{y}\mathcal{P}}{\mathcal{P}}\right)^{2}-2\eta\left(\frac{\partial_{x}\mathcal{P}}{\mathcal{P}}\right)^{2}
  +\frac{3\chi}{4\mathcal{P}\sqrt{\dot{\mathfrak{U}}+\varsigma}}\left[\rho^{2}+\eta\left(\frac{\partial_{x}\dot{\mathfrak{U}}}{\dot{\mathfrak{U}}+\varsigma}\right)^{2}\right]\\
 &-\frac{\partial_{y}^{-1}\left[\partial_{x}\mathfrak{U}\right]\varsigma'+\varsigma''-\partial_{y}\eta\partial_{x}^{2}\mathbb{U}+\eta\partial_{x}^{2}\partial_{y}u_{0}+\partial_{y}\hat{F}}{2(\dot{\mathfrak{U}}+\varsigma)}\\
 &+(1-\chi)\frac{\partial_{y}^{-1}\left[\partial_{x}\mathfrak{U}\right]\varsigma'+\varsigma''-\partial_{y}\eta\partial_{x}^{2}\mathbb{U}+\eta\partial_{x}^{2}\partial_{y}u_{0}+\partial_{y}\hat{F}}{2\mathcal{P}(\dot{\mathfrak{U}}+\varsigma)}\\
 &-\frac{-\partial_{y}^{-1}[\partial_{x}\mathfrak{U}]\chi'+\chi'\rho-\chi''}{\mathcal{P}\sqrt{\dot{\mathfrak{U}}+\varsigma}}
  -\frac{\partial_{y}^{-1}\left[\partial_{x}\mathfrak{U}\right]\chi'+\chi''}{\mathcal{P}},
\end{align*}
$$\mathbf{L}_{1}=\mathfrak{U}+2\eta\frac{\partial_{x}\mathcal{P}}{\mathcal{P}},
~\mathbf{L}_{2}=2\frac{\partial_{y}\mathcal{P}}{\mathcal{P}}-\partial_{y}^{-1}\left[\partial_{x}\mathfrak{U}\right],$$
and $\mathcal{T}_{1}=\mathcal{P}(-\partial_{y}\bar{\mathcal{T}}_{1}+\rho\bar{\mathcal{T}}_{1})$, with
$$
\bar{\mathcal{T}}_{1}(t,x,y)=
 \sum_{i=2}^{r-1}\binom{r}{i}\partial_{x}^{r-i+1}\mathbb{U}\partial_{x}^{i}\mathbb{U}
-\sum_{i=1}^{r-2}\binom{r}{i}\partial_{y}^{-1}[\partial_{x}^{i+1}\mathbb{U}]\partial_{x}^{r-i}\partial_{y}\mathbb{U},
$$
and
$$
\mathcal{T}_{2}(t,x,y)
= \mathcal{P}\left(-\partial_{y}^{-1}[\partial_{x}^{r+1}\tilde{\mathfrak{U}}](\varsigma'-\rho\varsigma)
 +2\eta\partial_{x}\rho\partial_{x}^{r+1}\mathbb{U}+\eta'\partial_{x}^{r+2}\mathbb{U}\right),
$$
$$
\mathcal{T}_{3}(t,x,y)
=r\mathcal{P}\left(\partial_{x}\partial_{y}^{2}\mathfrak{U}-\rho\partial_{x}\partial_{y}\mathfrak{U}\right)\partial_{y}^{-1}[\partial_{x}^{r}\mathfrak{U}],
$$
and $\mathcal{T}_{4}(t,x,y)=\mathcal{P}\bar{\mathcal{T}}_{4}\frac{\partial_{x}^{r}\mathbb{U}}{\dot{\mathfrak{U}}+\varsigma}$, with
\begin{align*}
\bar{\mathcal{T}}_{4}(t,x,y)
=&\partial_{y}^{2}\hat{F}+\varsigma'\partial_{x}\mathfrak{U}+\varsigma'''+\partial_{y}^{-1}[\partial_{x}\mathfrak{U}]\varsigma''+\eta\partial_{x}^{2}\partial_{y}^{2}u_{0}-\varsigma\partial_{x}\partial_{y}\mathfrak{U}\\
 &-\rho\left(\partial_{y}\hat{F}+\partial_{y}^{-1}[\partial_{x}\mathfrak{U}]\varsigma'+\varsigma''+\eta\partial_{x}^{2}\partial_{y}u_{0}\right)
  -2\eta\partial_{x}\partial_{y}\mathfrak{U}\partial_{x}\rho,
\end{align*}
and
$$
\mathcal{T}_{5}=\mathcal{P}(\varsigma\partial_{x}^{r+1}u-\partial_{y}\partial_{x}^{r}\mathbb{F}+\rho\partial_{x}^{r}\mathbb{F}),
$$
with $\mathbb{F}=-u\partial_{x}u+(1-h)\partial_{y}^{-1}[\partial_{x}u]\partial_{y}u$, and finally, $\mathcal{T}_{6}=\mathcal{P}\bar{\mathcal{T}}_{6}$, with
\begin{align*}
\bar{\mathcal{T}}_{6}
=& \rho\sum_{i=2}^{r}\binom{r}{i}\partial_{x}^{r-i+1}\mathbb{U}\partial_{x}^{i}u_{0}
  +\rho\sum_{i=0}^{r-1}\binom{r}{i}\partial_{x}^{r-i+1}u_{0}\partial_{x}^{i}\mathbb{U}
  -\sum_{i=0}^{r}\binom{r}{i}\partial_{y}\partial_{x}^{r-i+1}u_{0}\partial_{x}^{i}u_{0}\\
 &-\sum_{i=0}^{r-1}\binom{r}{i}\partial_{y}\partial_{x}^{r-i+1}u_{0}\partial_{x}^{i}\mathbb{U}
  -\sum_{i=2}^{r}\binom{r}{i}\partial_{y}\partial_{x}^{r-i+1}\mathbb{U}\partial_{x}^{i}u_{0}\\
 &+\sum_{i=0}^{r-2}\binom{r}{i}\partial_{y}^{-1}[\partial_{x}^{i+1}\mathbb{U}]\partial_{x}^{r-i}\partial_{y}^{2}u_{0}
  +\sum_{i=1}^{r-2}\binom{r}{i}\partial_{y}^{-1}[\partial_{x}^{i+1}u_{0}]\partial_{x}^{r-i}\partial_{y}^{2}\mathbb{U}
  +\sum_{i=0}^{r-2}\binom{r}{i}\partial_{y}^{-1}[\partial_{x}^{i+1}u_{0}]\partial_{x}^{r-i}\partial_{y}^{2}u_{0}\\
 &-\rho\sum_{i=0}^{r-2}\binom{r}{i}\partial_{y}^{-1}[\partial_{x}^{i+1}\mathbb{U}]\partial_{x}^{r-i}\partial_{y}u_{0}
  -\rho\sum_{i=1}^{r-2}\binom{r}{i}\partial_{y}^{-1}[\partial_{x}^{i+1}u_{0}]\partial_{x}^{r-i}\partial_{y}\mathbb{U}
  -\rho\sum_{i=0}^{r-2}\binom{r}{i}\partial_{y}^{-1}[\partial_{x}^{i+1}u_{0}]\partial_{x}^{r-i}\partial_{y}u_{0}\\
 &+\rho\partial_{x}^{r}(\partial_{t}u_{0}+u_{0}\partial_{x}u_{0}+\partial_{x}P)
  +\partial_{y}^{3}\partial_{x}^{r}u_{0}
  -\rho\partial_{x}^{r}\partial_{y}^{2}u_{0}
  -\partial_{t}\partial_{y}\partial_{x}^{r}u_{0}.
\end{align*}

Let $\varphi_{r}=\varphi, \hat{\Psi}_{r}=\hat{\Psi}=\y^{-\frac{1}{2}}\Psi$ when $r=s$ and  $\varphi_{r}=\hat{\varphi},\hat{\Psi}_{r}=\y^{-1}\Psi$ when $r=s+1$, where $\Psi$ is given below Corollary \eqref{ncC1}.
Correspondingly, let $\tilde{\varphi}_{r}=\tilde{\varphi}, \tilde{\hat{\Psi}}_{r}=\y^{-\frac{1}{2}}\tilde{\mathcal{P}}^{2}(t,x,y)\varphi_{o}(t,y)$ when $r=s$ and  $\tilde{\varphi}_{r}=\y^{-\frac{1}{2}}\tilde{\varphi},\tilde{\hat{\Psi}}_{r}=\y^{-1}\tilde{\mathcal{P}}^{2}(t,x,y)\varphi_{o}(t,y)$ when $r=s+1$.
For any $t\in[0,t_{3}]$, multiplying $\varphi_{r}^{2}\mathbf{U}_{r}$ and then integrating over $\Omega$, by using the integration by parts we get
\begin{align}\label{UnE}
&\frac{1}{2}\frac{d}{dt}\|\mathbf{U}_{r}\|_{L_{\varphi_{r}}^{2}(\Omega)}^{2}
+\|\partial_{y}\mathbf{U}_{r}\|_{L_{\varphi_{r}}^{2}(\Omega)}^{2}
+\|\sqrt{\eta}\partial_{x}\mathbf{U}_{r}\|_{L_{\varphi_{r}}^{2}(\Omega)}^{2}\\
=&~\frac{1}{2}(\mathbf{U}_{r},(\partial_{t}\varphi_{r}^{2}+\partial_{y}^{2}\varphi_{r}^{2})\mathbf{U}_{r})
  -(\mathbf{L}_{1}\partial_{x}\mathbf{U}_{r},\varphi_{r}^{2}\mathbf{U}_{r})
  -(\mathbf{L}_{2}\partial_{y}\mathbf{U}_{r},\varphi_{r}^{2}\mathbf{U}_{r})
  -(\mathbf{L}_{0}\mathbf{U}_{r},\varphi_{r}^{2}\mathbf{U}_{r})\nonumber\\
 &+(\mathbf{U}_{r},\varphi_{r}\partial_{y}\varphi_{r}\mathbf{U}_{r})_{L_{x}^{2}}|_{y=0}
  -(\partial_{y}\mathbf{U}_{r},\varphi_{r}^{2}\mathbf{U}_{r})_{L_{x}^{2}}|_{y=0}
  +\sum_{i=1}^{6}(\mathcal{T}_{i},\varphi_{r}^{2}\mathbf{U}_{r})\nonumber\\
:=&\sum_{i=0}^{11}J_{i}.\nonumber
\end{align}


To get the a posterior estimates of the approximate solutions, we shall estimate $J_{1}$, $J_{2}$,...,$J_{13}$ term by term.
In what follows, we shall always take $\lambda, \Lambda, \theta, \iota, \mathcal{Z}$ to be positive constants to be determined later, satisfying $\lambda>\lambda_{*}, \Lambda>\Lambda_{*}, \theta>\theta^{*}, \iota\in(0,\iota^{*}), \mathcal{Z}>\mathcal{Z}_{*}$, where $\lambda_{*}, \Lambda_{*},\iota^{*}, \mathcal{Z}_{*}$ are given in Theorem \ref{wpae}, and $\theta^{*}$ is given in Lemma \ref{wdotu}.

\noindent{\bf \underline{Estimate of~$J_{0}$.}}
When $r=s$, parallel to the estimate of $I_{0}$ in Section 4.4, one has
\begin{align}\label{J00}
J_{0}\leq -\frac{\lambda}{64}\|\sqrt{\omega_{\lambda}}\mathbf{U}_{s}\|_{L_{\varphi}^{2}(\Omega)}^{2}
          -\frac{3}{4}\Lambda\delta\|\sqrt{\y}\mathbf{U}_{s}\|_{L_{\varphi}^{2}(\Omega)}^{2}.
\end{align}
When $r=s+1$, it follows from Lemma \ref{weighte} that
\begin{align}\label{J01}
J_{0}
=&~\frac{1}{2}(\mathbf{U}_{s+1},(\partial_{t}\varphi^{2}+\partial_{y}^{2}\varphi^{2})\y^{-1}\mathbf{U}_{s+1})
  +(\mathbf{U}_{s+1},\partial_{y}\varphi^{2}\partial_{y}\y^{-1}\mathbf{U}_{s+1})
  +\frac{1}{2}(\mathbf{U}_{s+1},\varphi^{2}\partial_{y}^{2}\y^{-1}\mathbf{U}_{s+1})\\
\leq&-\frac{\lambda}{64}\|\sqrt{\omega_{\lambda}}\mathbf{U}_{s+1}\|_{L_{\hat{\varphi}}^{2}(\Omega)}^{2}
     -\frac{3}{4}\Lambda\delta\|\sqrt{\y}\mathbf{U}_{s+1}\|_{L_{\hat{\varphi}}^{2}(\Omega)}^{2}
     +C\lambda(\mathbf{U}_{s+1},\sqrt{\omega_{\lambda}}\hat{\varphi}^{2}\mathbf{U}_{s+1})
     +C_{*}\|\mathbf{U}_{s+1}\|_{L_{\hat{\varphi}}^{2}(\Omega)}^{2}\nonumber\\
\leq&-\left(\frac{\lambda}{64}-\frac{1}{16}\right)\|\sqrt{\omega_{\lambda}}\mathbf{U}_{s+1}\|_{L_{\hat{\varphi}}^{2}(\Omega)}^{2}
     -\frac{3}{4}\Lambda\delta\|\sqrt{\y}\mathbf{U}_{s+1}\|_{L_{\hat{\varphi}}^{2}(\Omega)}^{2}
     +C_{*}\|\mathbf{U}_{s+1}\|_{L_{\hat{\varphi}}^{2}(\Omega)}^{2}.\nonumber
\end{align}
As a consequence of \eqref{J00} and \eqref{J01}, one has
\begin{align*}
J_{0}
\leq -\left(\frac{\lambda}{64}-\frac{1}{16}\right)\|\sqrt{\omega_{\lambda}}\mathbf{U}_{r}\|_{L_{\varphi_{r}}^{2}(\Omega)}^{2}
     -\frac{3}{4}\Lambda\delta\|\sqrt{\y}\mathbf{U}_{r}\|_{L_{\varphi_{r}}^{2}(\Omega)}^{2}
     +C_{*}\|\mathbf{U}_{r}\|_{L_{\varphi_{r}}^{2}(\Omega)}^{2}.
\end{align*}



\noindent{\bf \underline{Estimate of~$J_{1}$.}}
By using the integration by parts and Lemma \ref{rhoP}, one has
\begin{align*}
J_{1}
=\frac{1}{2}(\partial_{x}\mathbf{L}_{1}\mathbf{U}_{r},\varphi_{r}^{2}\mathbf{U}_{r})
=&~\frac{1}{2}(\partial_{x}\mathfrak{U}\mathbf{U}_{r},\varphi_{r}^{2}\mathbf{U}_{r})
  +\left(\eta\partial_{x}\left(\frac{\partial_{x}\mathcal{P}}{\mathcal{P}}\right)\mathbf{U}_{r},\varphi_{r}^{2}\mathbf{U}_{r}\right)\\
=&~\frac{1}{2}(\partial_{x}\mathfrak{U}\mathbf{U}_{r},\varphi_{r}^{2}\mathbf{U}_{r})
  -\frac{1}{2}\int_{\mathbb{T}}\int_{\check{y}_{*}}^{\hat{y}_{*}}\eta\varphi_{r}^{2}\left[\frac{\partial_{x}^{2}\dot{\mathfrak{U}}}{\dot{\mathfrak{U}}+\varsigma}-\left(\frac{\partial_{x}\dot{\mathfrak{U}}}{\dot{\mathfrak{U}}+\varsigma}\right)^{2}\right]\mathbf{U}_{r}^{2}dydx\\
 &+\int_{\mathbb{T}}\int_{[\check{y},\hat{y}]\setminus[\check{y}_{*},\hat{y}_{*}]}\eta\varphi_{r}^{2}\partial_{x}\left(\frac{\partial_{x}\mathcal{P}}{\mathcal{P}}\right)\mathbf{U}_{r}^{2}dydx\\
\leq&~C_{*}\|\mathbf{U}_{r}\|_{L_{\varphi_{r}}^{2}(\Omega)}^{2}.
\end{align*}



\noindent{\bf \underline{Estimate of~$J_{2}$.}}
By the definition of $\mathbf{L}_{2}$, one can decompose $J_{2}$ into two parts,
\begin{align*}
J_{2}
= (\partial_{y}^{-1}[\partial_{x}]\partial_{y}\mathfrak{U}\mathbf{U}_{r},\varphi_{r}^{2}\mathbf{U}_{r})
 -2\left(\frac{\partial_{y}\mathcal{P}}{\mathcal{P}} \partial_{y}\mathbf{U}_{r},\varphi_{r}^{2}\mathbf{U}_{r}\right)
:=J_{21}+J_{22}.
\end{align*}
The estimate of $J_{21}$ is parallel to that of $I_{21}$ in Section 4.6, and one has
\begin{align*}
J_{21}
=&~-\frac{1}{2}(\partial_{x}\mathfrak{U}\mathbf{U}_{r},\varphi_{r}^{2}\mathbf{U}_{r})
  -\frac{1}{2}(\partial_{y}^{-1}\left[\partial_{x}\mathfrak{U}\right]\mathbf{U}_{r},\partial_{y}\varphi_{r}^{2}\mathbf{U}_{r})\\
=&~-\frac{1}{2}(\partial_{x}\mathfrak{U}\mathbf{U}_{r},\varphi_{r}^{2}\mathbf{U}_{r})
  -\frac{1}{2}(\chi\partial_{y}^{-1}\left[\partial_{x}\mathfrak{U}\right]\mathbf{U}_{r},\partial_{y}\varphi_{r}^{2}\mathbf{U}_{r})
  -\frac{1}{2}((1-\chi)\partial_{y}^{-1}\left[\partial_{x}\mathfrak{U}\right]\mathbf{U}_{r},\partial_{y}\varphi_{r}^{2}\mathbf{U}_{r})\\
\leq&~C_{*}\|\mathbf{U}_{r}\|_{L_{\varphi_{r}}^{2}(\Omega)}^{2}
     +C_{*}\lambda(\mathbf{U}_{r},\sqrt{\omega_{\lambda}}\varphi_{r}^{2}\mathbf{U}_{r})
     +C_{*}(\mathbf{U}_{r},\varphi_{r}^{2}\mathbf{U}_{r})
     +C\|\y^{-1}\partial_{y}^{-1}\left[\partial_{x}\mathfrak{U}\right]\|_{L^{\infty}(\Omega)}
       \|\sqrt{\y}\mathbf{U}_{r}\|_{L_{\varphi_{r}}^{2}(\Omega)}^{2}\\
\leq&~C_{*}\|\mathbf{U}_{r}\|_{L_{\varphi_{r}}^{2}(\Omega)}^{2}
     +\frac{1}{16}\|\sqrt{\omega_{\lambda}}\mathbf{U}_{r}\|_{L_{\varphi_{r}}^{2}(\Omega)}^{2}
     +C_{*}\|\sqrt{\y}\mathbf{U}_{r}\|_{L_{\varphi_{r}}^{2}(\Omega)}^{2}.
\end{align*}
Noticing ${\bf supp}~\partial_{y}\mathcal{P}\subset(\check{y},\hat{y})$ and $\mathcal{P}^{-1}\partial_{y}\mathcal{P}=-\frac{1}{2}\rho$ on $[\check{y}_{*},\hat{y}_{*}]$, one has
\begin{align*}
J_{22}
=&~\int_{\mathbb{T}}\int_{\check{y}_{*}}^{\hat{y}_{*}}\varphi_{r}^{2}\rho\mathbf{U}_{r}\partial_{y}\mathbf{U}_{r}dydx
  -2\int_{\mathbb{T}}\int_{[\check{y},\hat{y}]\setminus[\check{y}_{*},\hat{y}_{*}]}\varphi_{r}^{2}\frac{\partial_{y}\mathcal{P}}{\mathcal{P}}\mathbf{U}_{r}\partial_{y}\mathbf{U}_{r}dydx\\
\leq&~C_{*}\|\mathbf{U}_{r}\|_{L_{\varphi_{r}}^{2}(\Omega)}^{2}
     +C\|\rho\mathbf{U}_{r}\|_{L_{\varphi_{r}}^{2}(\Omega)}^{2}
     +\frac{1}{16}\|\partial_{y}\mathbf{U}_{r}\|_{L_{\varphi_{r}}^{2}(\Omega)}^{2}\\
\leq&~C_{*}\|\mathbf{U}_{r}\|_{L_{\varphi_{r}}^{2}(\Omega)}^{2}
     +C\|\mathcal{P}\mathbf{U}_{r}\|_{L_{\varphi_{r}}^{2}(\Omega)}^{2}
     +\frac{1}{16}\|\partial_{y}\mathbf{U}_{r}\|_{L_{\varphi_{r}}^{2}(\Omega)}^{2}.
\end{align*}
To sum up, one has
\begin{align*}
J_{2}
\leq&~C_{*}\|\mathbf{U}_{r}\|_{L_{\varphi_{r}}^{2}(\Omega)}^{2}
     +\frac{1}{16}\|\sqrt{\omega_{\lambda}}\mathbf{U}_{r}\|_{L_{\varphi_{r}}^{2}(\Omega)}^{2}
     +C_{*}\|\sqrt{\y}\mathbf{U}_{r}\|_{L_{\varphi_{r}}^{2}(\Omega)}^{2}\\
    &+C\|\mathcal{P}\mathbf{U}_{r}\|_{L_{\varphi_{r}}^{2}(\Omega)}^{2}
     +\frac{1}{16}\|\partial_{y}\mathbf{U}_{r}\|_{L_{\varphi_{r}}^{2}(\Omega)}^{2}.
\end{align*}

\noindent{\bf \underline{Estimate of~$J_{3}$.}}
As is shown in \eqref{I33}, one has
\begin{align}\label{J31}
&\frac{\partial_{y}^{-1}\left[\partial_{x}\dot{\mathfrak{U}}\right]\varsigma'+\varsigma''-\partial_{y}\eta\partial_{x}^{2}\mathbb{U}+\eta\partial_{x}^{2}\partial_{y}u_{0}}{2(\dot{\mathfrak{U}}+\varsigma)}
 -(1-\chi)\frac{\partial_{y}^{-1}\left[\partial_{x}\mathfrak{U}\right]\varsigma'+\varsigma''-\partial_{y}\eta\partial_{x}^{2}\mathbb{U}+\eta\partial_{x}^{2}\partial_{y}u_{0}}{2\mathcal{P}(\dot{\mathfrak{U}}+\varsigma)}\\
&\frac{-\partial_{y}^{-1}[\partial_{x}\mathfrak{U}]\chi'+\chi'\rho-\chi''}{\mathcal{P}\sqrt{\dot{\mathfrak{U}}+\varsigma}}
+\frac{\partial_{y}^{-1}\left[\partial_{x}\mathfrak{U}\right]\chi'+\chi''}{\mathcal{P}}
\leq C_{*}\iota^{-1}.\nonumber
\end{align}
Since $|\partial_{y}\hat{F}|\leq C_{*}\iota^{-1}(\dot{\mathfrak{U}}+\bar{\varsigma})$ and $|\partial_{y}\rho|\leq C\mathcal{P}^{2}$, one has
\begin{align}\label{J32}
2\partial_{y}\rho-r\partial_{x}\mathfrak{U}
+\frac{\partial_{y}\hat{F}}{2(\dot{\mathfrak{U}}+\varsigma)}
-(1-\chi)\frac{\partial_{y}\hat{F}}{2\mathcal{P}(\dot{\mathfrak{U}}+\varsigma)}\leq C\mathcal{P}^{2}+C_{*}\iota^{-1}.
\end{align}
Combining \eqref{J31} and \eqref{J32} with \eqref{I32}, it follows that
\begin{align*}
-\mathbf{L}_{0}\leq C\mathcal{P}^{2}+C_{*}\iota^{-1},
\end{align*}
which gives
\begin{align*}
J_{3}
\leq C_{*}\iota^{-1}\|\mathbf{U}_{r}\|_{L_{\varphi_{r}}^{2}(\Omega)}^{2}
    +C\|\mathcal{P}\mathbf{U}_{r}\|_{L_{\varphi_{r}}^{2}(\Omega)}^{2}.
\end{align*}

\noindent{\bf \underline{Estimate of~$J_{4}$.}}
Using the trace theorem, we know that
\begin{align*}
J_{4}
=(\mathbf{U}_{r},\varphi_{r}\partial_{y}\varphi_{r}\mathbf{U}_{r})_{L_{x}^{2}}|_{y=0}
\leq C\|\mathbf{U}_{r}\|_{L_{\varphi_{r}}^{2}(\Omega)}^{2}
     +\frac{1}{16}\|\partial_{y}\mathbf{U}_{r}\|_{L_{\varphi_{r}}^{2}(\Omega)}^{2}.
\end{align*}


\noindent{\bf \underline{Estimate of~$J_{5}$.}}
With the help of the boundary condition in \eqref{Un}, by using the trace estimate one can deduce that
\begin{align*}
J_{5}
=-(\varphi_{r}^{2}\partial_{y}\mathbf{U}_{r},\mathbf{U}_{r})_{L_{x}^{2}}|_{y=0}
=& \varphi_{0}^{2}(\rho \mathbf{U}_{r},\mathbf{U}_{r})_{L_{x}^{2}}|_{y=0}
  -\left.\varphi_{0}^{2}(\partial_{x}^{r+1}P,\mathbf{U}_{r})_{L_{x}^{2}}\right|_{y=0}\\
\leq& C_{*}\|\mathbf{U}_{r}\|_{L_{\varphi_{r}}^{2}(\Omega)}^{2}
     +\frac{1}{16}\|\partial_{y}\mathbf{U}_{r}\|_{L_{\varphi_{r}}^{2}(\Omega)}^{2}
     +C\|\partial_{x}^{r+1}P\|_{L_{x}^{2}(\mathbb{T})}^{2},
\end{align*}
where $\varphi_{0}=\varphi(t,0)$.

\noindent{\bf \underline{Estimate of~$J_{6}$.}}
We shall estimate $J_{6}$ in different ways for $r=s$ and $r=s+1$ respectively.

If $r=s$, in view of Lemma \ref{sdotu} and Lemma \ref{rhoP} one has
\begin{align}\label{J6s1}
\left(\mathcal{P}\rho\bar{\mathcal{T}}_{1},\varphi^{2}\mathbf{U}_{s}\right)
=&~\int_{\mathbb{T}}\int_{\check{y}}^{\hat{y}}\varphi^{2}\mathcal{P}\rho\bar{\mathcal{T}}_{1}\mathbf{U}_{s}dydx
  +\int_{\mathbb{T}}\int_{\mathbb{R}_{+}\setminus[\check{y},\hat{y}]}\varphi^{2}\mathcal{P}\rho\bar{\mathcal{T}}_{1}\mathbf{U}_{s}dydx\\
\leq&~C_{*}\|\bar{\mathcal{T}}_{1}\|_{L^{2}_{x}(\mathbb{T})L^{\infty}_{y}[\check{y},\hat{y}]}
      \|\mathcal{P}\mathbf{U}_{s}\|_{L_{\varphi}^{2}(\Omega)}\|\mathcal{P}\varphi\|_{L_{y}^{2}[\check{y},\hat{y}]}
     +C_{*}\|\bar{\mathcal{T}}_{1}\|_{L_{\psi}^{2}(\Omega)}\|\mathbf{U}_{s}\|_{L_{\varphi}^{2}(\Omega)}\nonumber\\
\leq&~C_{*}\|\bar{\mathcal{T}}_{1}\|_{L^{2}_{x}(\mathbb{T})L^{\infty}_{y}[\check{y},\hat{y}]}
      \|\mathcal{P}\mathbf{U}_{s}\|_{L_{\varphi}^{2}(\Omega)}
      \left(\int_{\check{y}}^{\hat{y}}\frac{1}{((y-y_{*})^2+\bar{\varsigma}+t)^{1+\varepsilon_{0}}}dy\right)^{\frac{1}{2}}\nonumber\\
    &+C_{*}\|\bar{\mathcal{T}}_{1}\|_{L_{\psi}^{2}(\Omega)}\|\mathbf{U}_{s}\|_{L_{\varphi}^{2}(\Omega)}\nonumber\\
\leq&~C_{*}\|\mathbf{U}_{s}\|_{L_{\varphi}^{2}(\Omega)}^{2}
     +\frac{1}{2}\|\mathcal{P}\mathbf{U}_{s}\|_{L_{\varphi}^{2}(\Omega)}^{2}
     +\frac{C_{*}}{(\bar{\varsigma}+t)^{\frac{1}{2}+\varepsilon_{0}}}\|\bar{\mathcal{T}}_{1}\|_{L_{x}^{2}(\mathbb{T})L_{y}^{\infty}[\check{y},\hat{y}]}^{2}
     +C_{*}\|\bar{\mathcal{T}}_{1}\|_{L_{\psi}^{2}(\Omega)}^{2}.\nonumber
\end{align}
and
\begin{align}\label{J6s2}
\left(-\mathcal{P}\partial_{y}\bar{\mathcal{T}}_{1},\varphi^{2}\mathbf{U}_{s}\right)
\leq \frac{C_{*}}{(\bar{\varsigma}+t)^{\varepsilon_{0}}}\|\partial_{y}\bar{\mathcal{T}}_{1}\|_{L_{\psi}^{2}(\Omega)}^{2}
    +\frac{1}{2}\|\mathcal{P}\mathbf{U}_{s}\|_{L_{\varphi}^{2}(\Omega)}^{2}.
\end{align}
Recall the definition of $\bar{\mathcal{T}}_{1}$, using Sobolev imbedding inequality one has
\begin{align}\label{J6s3}
&\|\bar{\mathcal{T}}_{1}\|_{L_{x}^{2}(\mathbb{T})L_{y}^{\infty}[\check{y},\hat{y}]}^{2}
+\|\bar{\mathcal{T}}_{1}\|_{L_{\psi}^{2}(\Omega)}^{2}
+\|\partial_{y}\bar{\mathcal{T}}_{1}\|_{L_{\psi}^{2}(\Omega)}^{2}\\
\leq&~C\sum_{i=2}^{s-1}\|\partial_{x}^{s-i+1}\mathbb{U}\partial_{x}^{i}\mathbb{U}\|_{L_{x}^{2}(\mathbb{T})L_{y}^{\infty}[\check{y},\hat{y}]}^{2}
     +C\sum_{i=1}^{s-2}\|\partial_{y}^{-1}[\partial_{x}^{i+1}\mathbb{U}]\partial_{x}^{s-i}\partial_{y}\mathbb{U}\|_{L_{x}^{2}(\mathbb{T})L_{y}^{\infty}[\check{y},\hat{y}]}^{2}\nonumber\\
    &+C\sum_{i=2}^{s-1}\|\partial_{x}^{s-i+1}\mathbb{U}\partial_{x}^{i}\mathbb{U}\|_{L_{\psi}^{2}(\Omega)}^{2}
     +C\sum_{i=1}^{s-2}\|\partial_{y}^{-1}[\partial_{x}^{i+1}\mathbb{U}]\partial_{x}^{s-i}\partial_{y}\mathbb{U}\|_{L_{\psi}^{2}(\Omega)}^{2}\nonumber\\
    &+C\sum_{i=2}^{s-1}\|\partial_{y}\partial_{x}^{s-i+1}\mathbb{U}\partial_{x}^{i}\mathbb{U}\|_{L_{\psi}^{2}(\Omega)}^{2}
     +C\sum_{i=1}^{s-2}\|\partial_{y}^{-1}[\partial_{x}^{i+1}\mathbb{U}]\partial_{x}^{s-i}\partial_{y}^{2}\mathbb{U}\|_{L_{\psi}^{2}(\Omega)}^{2}\nonumber\\
\leq&~C_{*}\|\mathbb{U}\|_{\hat{H}_{\psi}^{s}(\Omega)}^{4}
\leq C_{*}.\nonumber
\end{align}
Combining \eqref{J6s1} with \eqref{J6s2}, from \eqref{J6s3} we have the following estimate of $J_{6}$ when $r=s$,
\begin{align}\label{J6s}
J_{6}
\leq  C_{*}\|\mathbf{U}_{s}\|_{L_{\varphi}^{2}(\Omega)}^{2}
     +\|\mathcal{P}\mathbf{U}_{s}\|_{L_{\varphi}^{2}(\Omega)}^{2}
     +C_{*}\left(1+\frac{1}{(\bar{\varsigma}+t)^{\frac{1}{2}+\varepsilon_{0}}}\right).
\end{align}

If $r=s+1$, the terms involving partial derivativesof order less than $s$ with respect to $x$ in $J_{6}$ can be estimated in a similar way as in the case $r=s$ above, while other terms must be estimated in different manner since they cannot be controlled by $\|\mathbb{U}\|_{\hat{H}_{\psi}^{s}(\Omega)}$. For this purpose, we divide $\mathcal{T}_{1}$ into five parts as follows,
\begin{align*}
\mathcal{T}_{1}
=&-\frac{s(s+1)}{2}\partial_{x}^{2}\mathbb{U}\mathbf{U}_{s}
  +(s+1)\mathcal{P}\left(\rho\partial_{x}^{2}\mathbb{U}-\partial_{y}\partial_{x}^{2}\mathbb{U}\right)\partial_{x}^{s}\mathbb{U}
  +\frac{s(s+1)}{2}\mathcal{P}\partial_{y}^{-1}[\partial_{x}^{s}\mathbb{U}](\partial_{x}^{2}\partial_{y}^{2}\mathbb{U}-\rho\partial_{x}^{2}\partial_{y}\mathbb{U})\\
 &+(s+1)\mathcal{P}\partial_{y}^{-1}[\partial_{x}^{2}\mathbb{U}](\partial_{x}^{s}\partial_{y}^{2}\mathbb{U}-\rho\partial_{x}^{s}\partial_{y}\mathbb{U})
  +\mathcal{P}(\rho\bar{\mathcal{T}}_{1R}-\partial_{y}\bar{\mathcal{T}}_{1R}),
\end{align*}
with
\begin{align*}
\bar{\mathcal{T}}_{1R}(t,x,y)=
 \sum_{i=3}^{s-1}\binom{s+1}{i}\partial_{x}^{s-i+2}\mathbb{U}\partial_{x}^{i}\mathbb{U}
-\sum_{i=2}^{s-2}\binom{s+1}{i}\partial_{y}^{-1}[\partial_{x}^{i+1}\mathbb{U}]\partial_{x}^{s-i+1}\partial_{y}\mathbb{U}.
\end{align*}
Accordingly,
\begin{align}\label{J6}
J_{6}
=&-\frac{s(s+1)}{2}\left(\partial_{x}^{2}\mathbb{U}\mathbf{U}_{s},\hat{\varphi}^{2}\mathbf{U}_{s+1}\right)
  +(s+1)\left(\mathcal{P}\left(\rho\partial_{x}^{2}\mathbb{U}-\partial_{y}\partial_{x}^{2}\mathbb{U}\right)\partial_{x}^{s}\mathbb{U},\hat{\varphi}^{2}\mathbf{U}_{s+1}\right)\\
 &+\frac{s(s+1)}{2}\left(\mathcal{P}\partial_{y}^{-1}[\partial_{x}^{s}\mathbb{U}](\partial_{x}^{2}\partial_{y}^{2}\mathbb{U}-\rho\partial_{x}^{2}\partial_{y}\mathbb{U}),\hat{\varphi}^{2}\mathbf{U}_{s+1}\right)\nonumber\\
 &+(s+1)\left(\mathcal{P}\partial_{y}^{-1}[\partial_{x}^{2}\mathbb{U}](\rho\partial_{x}^{s}\partial_{y}\mathbb{U}-\partial_{x}^{s}\partial_{y}^{2}\mathbb{U}),\hat{\varphi}^{2}\mathbf{U}_{s+1}\right)\nonumber\\
 &+\left(\mathcal{P}(\rho\bar{\mathcal{T}}_{1R}-\partial_{y}\bar{\mathcal{T}}_{1R}),\hat{\varphi}^{2}\mathbf{U}_{s+1}\right).\nonumber
\end{align}
It is clearly that
\begin{align}\label{J61}
-\frac{s(s+1)}{2}\left(\partial_{x}^{2}\mathbb{U}\mathbf{U}_{s},\hat{\varphi}^{2}\mathbf{U}_{s+1}\right)
\leq& C_{*}\|\mathbf{U}_{s}\|_{L_{\varphi}^{2}(\Omega)}\|\mathbf{U}_{s+1}\|_{L_{\hat{\varphi}}^{2}(\Omega)}\\
\leq& C_{*}\|\mathbf{U}_{s}\|_{L_{\varphi}^{2}(\Omega)}^{2}
     +C_{*}\|\mathbf{U}_{s+1}\|_{L_{\hat{\varphi}}^{2}(\Omega)}^{2}.\nonumber
\end{align}
Using Corollary \ref{ncC2}, one can deduce
\begin{align}\label{J62}
&(s+1)\left(\mathcal{P}\left(\rho\partial_{x}^{2}\mathbb{U}-\partial_{y}\partial_{x}^{2}\mathbb{U}\right)\partial_{x}^{s}\mathbb{U},\hat{\varphi}^{2}\mathbf{U}_{s+1}\right)\\
\leq&~C(\|\partial_{x}^{2}\mathbb{U}\|_{L^{\infty}(\Omega)}+\|\partial_{y}\partial_{x}^{2}\mathbb{U}\|_{L^{\infty}(\Omega)})
      \|\mathcal{P}^{2}\partial_{x}^{s}\mathbb{U}\|_{L_{\varphi}^{2}(\Omega)}
      \|\mathbf{U}_{s+1}\|_{L_{\hat{\varphi}}^{2}(\Omega)}\nonumber\\
\leq&~C_{*}\|\mathcal{P}\mathbf{U}_{s}\|_{L_{\varphi}^{2}(\Omega)}
      \|\mathbf{U}_{s+1}\|_{L_{\hat{\varphi}}^{2}(\Omega)}\nonumber\\
\leq&~C_{*}\|\mathbf{U}_{s+1}\|_{L_{\hat{\varphi}}^{2}(\Omega)}^{2}
     +\|\mathcal{P}\mathbf{U}_{s}\|_{L_{\varphi}^{2}(\Omega)}^{2}.\nonumber
\end{align}
Using Lemma \ref{rhoP} and recall the definition of $u_{0}$, one has
$\left(\|\mathcal{P}\partial_{x}^{2}\partial_{y}^{2}\mathbb{U}\|_{L_{x}^{\infty}L_{y}^{2}(\Omega)}
 +\|\mathcal{P}^{2}\partial_{x}^{2}\partial_{y}\mathbb{U}\|_{L_{x}^{\infty}L_{y}^{2}(\Omega)}\right)\leq C$.
Also by Corollary \ref{ncC2} one can obtain
\begin{align}\label{J63}
&\frac{s(s+1)}{2}\left(\mathcal{P}\partial_{y}^{-1}[\partial_{x}^{s}\mathbb{U}](\partial_{x}^{2}\partial_{y}^{2}\mathbb{U}-\rho\partial_{x}^{2}\partial_{y}\mathbb{U}),\hat{\varphi}^{2}\mathbf{U}_{s+1}\right)\\
\leq&~C_{*}\|\partial_{x}^{s}\mathbb{U}\|_{L_{\varphi}^{2}(\Omega)}
       \left(\|\mathcal{P}\partial_{x}^{2}\partial_{y}^{2}\mathbb{U}\|_{L_{x}^{\infty}L_{y}^{2}(\Omega)}+\|\mathcal{P}^{2}\partial_{x}^{2}\partial_{y}\mathbb{U}\|_{L_{x}^{\infty}L_{y}^{2}(\Omega)}\right)
       \|\mathbf{U}_{s+1}\|_{L_{\hat{\varphi}}^{2}(\Omega)}\nonumber\\
\leq&~C_{*}\|\mathcal{P}\mathbf{U}_{s}\|_{L_{\varphi}^{2}(\Omega)}
       \|\mathbf{U}_{s+1}\|_{L_{\hat{\varphi}}^{2}(\Omega)}\nonumber\\
\leq&~C_{*}\|\mathbf{U}_{s+1}\|_{L_{\hat{\varphi}}^{2}(\Omega)}^{2}
     +\|\mathcal{P}\mathbf{U}_{s}\|_{L_{\varphi}^{2}(\Omega)}^{2}.\nonumber
\end{align}
From Corollary \ref{wdotuC} and Lemma \ref{sdotu} one can deduce that for any $t\in[0,t_{3}]$ and $x\in\mathbb{T}$,
\begin{align*}
|\partial_{y}^{-1}[\partial_{x}^{2}\mathbb{U}]|
\leq \int_{0}^{y}\int_{\sigma}^{+\infty}|\partial_{y}\partial_{x}^{2}\mathbb{U}|d\sigma dy
\leq C\int_{0}^{y}\int_{\sigma}^{+\infty}(1+\y t)e^{-\frac{3\delta\y}{4}}d\sigma dy
\leq C.
\end{align*}
As a consequence,  by utilizing Corollary \ref{ncC2} one has
\begin{align}\label{J64'}
&(s+1)\left(\mathcal{P}\partial_{y}^{-1}[\partial_{x}^{2}\mathbb{U}](\rho\partial_{x}^{s}\partial_{y}\mathbb{U}-\partial_{x}^{s}\partial_{y}^{2}\mathbb{U}),\hat{\varphi}^{2}\mathbf{U}_{s+1}\right)\\
\leq&~C\left(\|\mathcal{P}\partial_{x}^{s}\partial_{y}\mathbb{U}\|_{L_{\varphi}^{2}(\Omega)}
             +\|\partial_{x}^{s}\partial_{y}^{2}\mathbb{U}\|_{L_{\varphi}^{2}(\Omega)}\right)
       \|\mathcal{P}\mathbf{U}_{s+1}\|_{L_{\hat{\varphi}}^{2}(\Omega)}\nonumber\\
\leq&~C_{\lambda}\left(\|\mathcal{P}\mathbf{U}_{s}\|_{L_{\varphi}^{2}(\Omega)}
             +\|\partial_{y}\mathbf{U}_{s}\|_{L_{\varphi}^{2}(\Omega)}
             +\|\mathbf{U}_{s}\|_{L_{\varphi}^{2}(\Omega)}\right)
       \|\mathcal{P}\mathbf{U}_{s+1}\|_{L_{\hat{\varphi}}^{2}(\Omega)}\nonumber\\
\leq&~C_{\lambda}\|\mathcal{P}\mathbf{U}_{s}\|_{L_{\varphi}^{2}(\Omega)}^{2}
     +C_{\lambda}\|\partial_{y}\mathbf{U}_{s}\|_{L_{\varphi}^{2}(\Omega)}^{2}
     +C\|\mathcal{P}\mathbf{U}_{s+1}\|_{L_{\hat{\varphi}}^{2}(\Omega)}^{2}
     +C_{\lambda}\|\mathbf{U}_{s}\|_{L_{\varphi}^{2}(\Omega)}^{2}\nonumber.
\end{align}

Finally, one can use a similar estimate as used in the case $r=s$ above to get
\begin{align}\label{J65}
\left(\mathcal{P}(\rho\bar{\mathcal{T}}_{1R}-\partial_{y}\bar{\mathcal{T}}_{1R}),\hat{\varphi}^{2}\mathbf{U}_{s+1}\right)
\leq  C_{*}\|\mathbf{U}_{s+1}\|_{L_{\hat{\varphi}}^{2}(\Omega)}^{2}
     +\|\mathcal{P}\mathbf{U}_{s+1}\|_{L_{\hat{\varphi}}^{2}(\Omega)}^{2}
     +C_{*}\left(1+\frac{1}{(\bar{\varsigma}+t)^{\frac{1}{2}+\varepsilon_{0}}}\right).
\end{align}

Substituting \eqref{J61},\eqref{J62},\eqref{J63},\eqref{J64'} and \eqref{J65} into \eqref{J6}, one can conclude that when $r=s+1$,
\begin{align}\label{J6'}
J_{6}
\leq& C_{\lambda}\|\mathcal{P}\mathbf{U}_{s}\|_{L_{\varphi}^{2}(\Omega)}^{2}
     +C\|\mathcal{P}\mathbf{U}_{s+1}\|_{L_{\hat{\varphi}}^{2}(\Omega)}^{2}
     +C_{*}\|\mathbf{U}_{s}\|_{L_{\varphi}^{2}(\Omega)}^{2}
     +C_{*}\|\mathbf{U}_{s+1}\|_{L_{\hat{\varphi}}^{2}(\Omega)}^{2}\\
    &+C_{\lambda}\|\partial_{y}\mathbf{U}_{s}\|_{L_{\varphi}^{2}(\Omega)}^{2}
     +C_{*}\left(1+\frac{1}{(\bar{\varsigma}+t)^{\frac{1}{2}+\varepsilon_{0}}}\right).\nonumber
\end{align}


\noindent{\bf \underline{Estimate of~$J_{7}$.}}
Recall
\begin{align}\label{J7}
J_{7}
= 2(\eta\mathcal{P}\partial_{x}\rho\partial_{x}^{r+1}\mathbb{U},\varphi_{r}^{2}\mathbf{U}_{r})
 +(\partial_{y}\eta\mathcal{P}\partial_{x}^{r+2}\mathbb{U},\varphi_{r}^{2}\mathbf{U}_{r})
 -(\mathcal{P}\partial_{y}^{-1}[\partial_{x}^{r+1}\tilde{\mathfrak{U}}](\varsigma'-\rho\varsigma),\varphi_{r}^{2}\mathbf{U}_{r}).
\end{align}
By using integration by parts, one has
\begin{align}\label{J71}
2(\eta\mathcal{P}\partial_{x}\rho\partial_{x}^{r+1}\mathbb{U},\varphi_{r}^{2}\mathbf{U}_{r})
=&-2(\eta\partial_{x}(\mathcal{P}\partial_{x}\rho)\partial_{x}^{r}\mathbb{U},\varphi_{r}^{2}\mathbf{U}_{r})
 -2(\eta\mathcal{P}\partial_{x}\rho\partial_{x}^{r}\mathbb{U},\varphi_{r}^{2}\partial_{x}\mathbf{U}_{r})\\
\leq& \|\eta\partial_{x}(\mathcal{P}\partial_{x}\rho)\partial_{x}^{r}\mathbb{U}\|_{L_{\varphi_{r}}^{2}(\Omega)}^{2}
     +C\|\sqrt{\eta}\mathcal{P}\partial_{x}\rho\partial_{x}^{r}\mathbb{U}\|_{L_{\varphi_{r}}^{2}(\Omega)}^{2}\nonumber\\
    &+\frac{1}{16}\|\sqrt{\eta}\partial_{x}\mathbf{U}_{r}\|_{L_{\varphi_{r}}^{2}(\Omega)}^{2}
     +C\|\mathbf{U}_{r}\|_{L_{\varphi_{r}}^{2}(\Omega)}^{2}.\nonumber
\end{align}
In virtue of Lemma \ref{rhoP}, by noticing that $\bar{\varsigma}=\iota\bar{\eta}^{2}$,
from Lemma \ref{nc} one has
\begin{align}\label{J711}
\|\eta\partial_{x}(\mathcal{P}\partial_{x}\rho)\partial_{x}^{r}\mathbb{U}\|_{L_{\varphi_{r}}^{2}(\Omega)}^{2}
\leq& C_{*}\bar{\eta}^{2}\|\partial_{x}^{r}\mathbb{U}\|_{L_{\hat{\Psi}_{r}}^{2}(\Omega)}^{2}\\
\leq& C_{*}\bar{\eta}^{2}\|\mathcal{P}\mathbf{U}_{r}\|_{L_{\varphi_{r}}^{2}(\Omega)}^{2}\nonumber\\
\leq& C_{*}\iota^{-1}\|\mathbf{U}_{r}\|_{L_{\varphi_{r}}^{2}(\Omega)}^{2},\nonumber
\end{align}
and
\begin{align}\label{J712}
\|\sqrt{\eta}\mathcal{P}\partial_{x}\rho\partial_{x}^{r}\mathbb{U}\|_{L_{\varphi_{r}}^{2}(\Omega)}^{2}
\leq& C_{*}\bar{\eta}\|\partial_{x}^{r}\mathbb{U}\|_{L_{\hat{\Psi}_{r}}^{2}(\Omega)}^{2}\\
\leq& C_{*}\iota^{-\frac{1}{2}}\|\mathbf{U}_{r}\|_{L_{\varphi_{r}}^{2}(\Omega)}\|\mathcal{P}\mathbf{U}_{r}\|_{L_{\varphi_{r}}^{2}(\Omega)}\nonumber\\
\leq& C_{*}\iota^{-1}\|\mathbf{U}_{r}\|_{L_{\varphi_{r}}^{2}(\Omega)}^{2}
     +\frac{1}{2}\|\mathcal{P}\mathbf{U}_{r}\|_{L_{\varphi_{r}}^{2}(\Omega)}^{2}.\nonumber
\end{align}
Substitute \eqref{J711} and \eqref{J712} into \eqref{J71}, one has
\begin{align}\label{J71'}
2(\eta\mathcal{P}\partial_{x}\rho\partial_{x}^{r+1}\mathbb{U},\varphi_{r}^{2}\mathbf{U}_{r})
\leq \frac{1}{16}\|\sqrt{\eta}\partial_{x}\mathbf{U}_{r}\|_{L_{\varphi_{r}}^{2}(\Omega)}^{2}
    +C_{*}\iota^{-1}\|\mathbf{U}_{r}\|_{L_{\varphi_{r}}^{2}(\Omega)}^{2}
    +\frac{1}{2}\|\mathcal{P}\mathbf{U}_{r}\|_{L_{\varphi_{r}}^{2}(\Omega)}^{2}.
\end{align}
Also by integration by parts, noticing that \eqref{theta} and Corollary \ref{ncC3}, one has
\begin{align*}
(\partial_{y}\eta\mathcal{P}\partial_{x}^{r+2}\mathbb{U},\varphi_{r}^{2}\mathbf{U}_{r})
=&-(\partial_{y}\eta\partial_{x}^{r+1}\mathbb{U},\varphi_{r}^{2}\partial_{x}\mathbf{U}_{r})\\
\leq&~C\theta t\|\vartheta_{c}\partial_{x}^{r+1}\mathbb{U}\|_{L_{\varphi_{r}}^{2}(\Omega)}
       \|\vartheta_{c}\partial_{x}\mathbf{U}_{r}\|_{L_{\varphi_{r}}^{2}(\Omega)}\\
\leq&~C_{\lambda}\bar{\eta}\theta\left(\sqrt{t}\|\vartheta_{c}\partial_{x}\mathbf{U}_{r}\|_{L_{\varphi_{r}}^{2}(\Omega)}
                                       +\|\vartheta_{c}\mathbf{U}_{r}\|_{L_{\varphi_{r}}^{2}(\Omega)}\right)
      \|\vartheta_{c}\partial_{x}\mathbf{U}_{r}\|_{L_{\varphi_{r}}^{2}(\Omega)}\\
\leq&~\left(\frac{1}{16}+C_{\lambda}\theta\sqrt{t}\right)\|\sqrt{\eta}\partial_{x}\mathbf{U}_{r}\|_{L_{\varphi_{r}}^{2}(\Omega)}^{2}
     +C_{*}C_{\theta}\|\mathbf{U}_{r}\|_{L_{\varphi_{r}}^{2}(\Omega)}^{2}.
\end{align*}
Noticing ${\bf supp}~\varsigma \subset(\check{y}_{*},\hat{y}_{*})$ and $\mathcal{P}\rho=-2\partial_{y}\mathcal{P}$ on $[\check{y}_{*},\hat{y}_{*}]$, by using the integration by parts one has
\begin{align}\label{J72}
-(\mathcal{P}\partial_{y}^{-1}[\partial_{x}^{r+1}\tilde{\mathfrak{U}}](\varsigma'-\rho\varsigma),\varphi_{r}^{2}\mathbf{U}_{r})
=&-(\mathcal{P}\varsigma'\partial_{y}^{-1}[\partial_{x}^{r+1}\tilde{\mathfrak{U}}],\varphi_{r}^{2}\mathbf{U}_{r})
  -2(\partial_{y}\mathcal{P}\varsigma\partial_{y}^{-1}[\partial_{x}^{r+1}\tilde{\mathfrak{U}}],\varphi_{r}^{2}\mathbf{U}_{r})\\
=& (\mathcal{P}\varsigma'\partial_{y}^{-1}[\partial_{x}^{r+1}\tilde{\mathfrak{U}}],\varphi_{r}^{2}\mathbf{U}_{r})
  +2(\mathcal{P}\varsigma\partial_{x}^{r+1}\tilde{\mathfrak{U}},\varphi_{r}^{2}\mathbf{U}_{r})\nonumber\\
 &+2(\mathcal{P}\varsigma\partial_{y}^{-1}[\partial_{x}^{r+1}\tilde{\mathfrak{U}}],\varphi_{r}^{2}\partial_{y}\mathbf{U}_{r})
  +4(\mathcal{P}\varsigma\partial_{y}^{-1}[\partial_{x}^{r+1}\tilde{\mathfrak{U}}],\varphi_{r}\partial_{y}\varphi_{r}\mathbf{U}_{r})\nonumber\\
=& (\mathcal{P}\varsigma'\partial_{y}^{-1}[\partial_{x}^{r+1}\tilde{\mathbb{U}}],\varphi_{r}^{2}\mathbf{U}_{r})
  +2(\mathcal{P}\varsigma\partial_{x}^{r+1}\tilde{\mathbb{U}},\varphi_{r}^{2}\mathbf{U}_{r})\nonumber\\
 &+2(\mathcal{P}\varsigma\partial_{y}^{-1}[\partial_{x}^{r+1}\tilde{\mathbb{U}}],\varphi_{r}^{2}\partial_{y}\mathbf{U}_{r})
  +4(\mathcal{P}\varsigma\partial_{y}^{-1}[\partial_{x}^{r+1}\tilde{\mathbb{U}}],\varphi_{r}\partial_{y}\varphi_{r}\mathbf{U}_{r})\nonumber\\
 &+(\mathcal{P}\varsigma'\partial_{y}^{-1}[\partial_{x}^{r+1}u_{0}],\varphi_{r}^{2}\mathbf{U}_{r})
  +2(\mathcal{P}\varsigma\partial_{x}^{r+1}u_{0},\varphi_{r}^{2}\mathbf{U}_{r})\nonumber\\
 &+2(\mathcal{P}\varsigma\partial_{y}^{-1}[\partial_{x}^{r+1}u_{0}],\varphi_{r}^{2}\partial_{y}\mathbf{U}_{r})
  +4(\mathcal{P}\varsigma\partial_{y}^{-1}[\partial_{x}^{r+1}u_{0}],\varphi_{r}\partial_{y}\varphi_{r}\mathbf{U}_{r})\nonumber.
\end{align}
Since $\sqrt{\bar{\varsigma}}\mathcal{P}\leq C$, $\sqrt{\tilde{\bar{\varsigma}}}\tilde{\mathcal{P}}\leq C$ and $\bar{\varsigma}=\iota\bar{\eta}^{2}$, it follows from Corollary \ref{ncC3} that
\begin{align}\label{J721}
(\mathcal{P}\varsigma'\partial_{y}^{-1}[\partial_{x}^{r+1}\tilde{\mathbb{U}}],\varphi_{r}^{2}\mathbf{U}_{r})
\leq&~C_{*}\bar{\varsigma}\|\vartheta_{c}\partial_{x}^{r+1}\tilde{\mathbb{U}}\|_{L_{\check{\psi}}^{2}(\Omega)}
                 \|\varphi_{r}\|_{L^{2}[\check{y},\hat{y}]}
                 \|\mathcal{P}\mathbf{U}_{r}\|_{L_{\varphi_{r}}^{2}(\Omega)}\\
\leq&~C_{*}\bar{\varsigma}\|\mathcal{P}\omega_{o}^{-1}\|_{L^{\infty}(\Omega)}
      \left(\|\vartheta_{c}\partial_{x}\tilde{\mathbf{U}}_{r}\|_{L_{\tilde{\varphi}_{r}}^{2}(\Omega)}
            +\|\tilde{\mathcal{P}}\tilde{\mathbf{U}}_{r}\|_{L_{\tilde{\varphi}_{r}}^{2}(\Omega)}\right)
      \|\mathcal{P}\mathbf{U}_{r}\|_{L_{\varphi_{r}}^{2}(\Omega)}\nonumber\\
\leq&~C_{*}\bar{\varsigma}\tilde{\bar{\varsigma}}^{-\frac{1}{2}}
      \left(\|\vartheta_{c}\partial_{x}\tilde{\mathbf{U}}_{r}\|_{L_{\tilde{\varphi}_{r}}^{2}(\Omega)}
            +\|\tilde{\mathcal{P}}\tilde{\mathbf{U}}_{r}\|_{L_{\tilde{\varphi}_{r}}^{2}(\Omega)}\right)
      \|\mathcal{P}\mathbf{U}_{r}\|_{L_{\varphi_{r}}^{2}(\Omega)}\nonumber\\
\leq&~\tau\|\sqrt{\tilde{\eta}}\partial_{x}\tilde{\mathbf{U}}_{r}\|_{L_{\tilde{\varphi}_{r}}^{2}(\Omega)}^{2}
     +C_{*}C_{\tau}\bar{\varsigma}^{2}\tilde{\bar{\varsigma}}^{-1}\tilde{\bar{\eta}}^{-1}\|\mathcal{P}\mathbf{U}_{r}\|_{L_{\varphi_{r}}^{2}(\Omega)}^{2}\nonumber\\
    &+C_{*}\bar{\varsigma}\tilde{\bar{\varsigma}}^{-\frac{1}{2}}\|\tilde{\mathcal{P}}\tilde{\mathbf{U}}_{r}\|_{L_{\tilde{\varphi}_{r}}^{2}(\Omega)}
                                              \|\mathcal{P}\mathbf{U}_{r}\|_{L_{\varphi_{r}}^{2}(\Omega)}\nonumber\\
\leq&~\tau\|\sqrt{\tilde{\eta}}\partial_{x}\tilde{\mathbf{U}}_{r}\|_{L_{\tilde{\varphi}_{r}}^{2}(\Omega)}^{2}
     +C_{*}C_{\tau}\|\mathbf{U}_{r}\|_{L_{\varphi_{r}}^{2}(\Omega)}\|\mathcal{P}\mathbf{U}_{r}\|_{L_{\varphi_{r}}^{2}(\Omega)}\nonumber\\
    &+C_{*}\|\tilde{\mathbf{U}}_{r}\|_{L_{\tilde{\varphi}_{r}}^{2}(\Omega)}\|\mathcal{P}\mathbf{U}_{r}\|_{L_{\varphi_{r}}^{2}(\Omega)}\nonumber\\
\leq&~\tau\|\sqrt{\tilde{\eta}}\partial_{x}\tilde{\mathbf{U}}_{r}\|_{L_{\tilde{\varphi}_{r}}^{2}(\Omega)}^{2}
     +C_{*}\|\mathbf{U}_{r}\|_{L_{\varphi_{r}}^{2}(\Omega)}^{2}
     +C_{*}\|\tilde{\mathbf{U}}_{r}\|_{L_{\tilde{\varphi}_{r}}^{2}(\Omega)}^{2}\nonumber\\
    &+\frac{1}{2}\|\mathcal{P}\mathbf{U}_{r}\|_{L_{\varphi_{r}}^{2}(\Omega)}^{2},\nonumber
\end{align}
Similarly, we have
\begin{align}\label{J722}
2\left|(\mathcal{P}\varsigma\partial_{x}^{r+1}\tilde{\mathbb{U}},\varphi_{r}^{2}\mathbf{U}_{r})\right|
\leq& C\bar{\varsigma}\|\tilde{\mathcal{P}}\vartheta_{c}\partial_{x}^{r+1}\tilde{\mathbb{U}}\|_{L_{\tilde{\varphi}_{r}}^{2}(\Omega)}
             \|\mathbf{U}_{r}\|_{L_{\varphi_{r}}^{2}(\Omega)}\\
\leq& C_{*}\tilde{\bar{\varsigma}}\|\vartheta_{c}\partial_{x}^{r+1}\tilde{\mathbb{U}}\|_{L_{\tilde{\hat{\Psi}}_{r}}^{2}(\Omega)}
                 \|\mathbf{U}_{r}\|_{L_{\varphi_{r}}^{2}(\Omega)}\nonumber\\
\leq& C_{*}\left(\|\sqrt{\tilde{\eta}}\partial_{x}\tilde{\mathbf{U}}_{r}\|_{L_{\tilde{\varphi}_{r}}^{2}(\Omega)}
            +\|\tilde{\mathbf{U}}_{r}\|_{L_{\tilde{\varphi}_{r}}^{2}(\Omega)}\right)
      \|\mathbf{U}_{r}\|_{L_{\varphi_{r}}^{2}(\Omega)}\nonumber\\
\leq& \tau\|\sqrt{\tilde{\eta}}\partial_{x}\tilde{\mathbf{U}}_{r}\|_{L_{\tilde{\varphi}_{r}}^{2}(\Omega)}^{2}
     +C_{*}C_{\tau}\|\mathbf{U}_{r}\|_{L_{\varphi_{r}}^{2}(\Omega)}^{2}
     +C_{*}\|\tilde{\mathbf{U}}_{r}\|_{L_{\tilde{\varphi}_{r}}^{2}(\Omega)}^{2}.\nonumber
\end{align}
Moreover, note that $\|\mathcal{P}\omega_{o}^{-1}\|_{L^{\infty}(\Omega)}\leq C\bar{\varsigma}^{-\frac{1}{2}+\frac{\epsilon_{0}}{2}}$ for any $t\in[0,t_{3}]$, by using Corollary \ref{ncC3} we can obtain
\begin{align}\label{J723}
&2(\mathcal{P}\varsigma\partial_{y}^{-1}[\partial_{x}^{r+1}\tilde{\mathbb{U}}],\varphi_{r}^{2}\partial_{y}\mathbf{U}_{r})\\
\leq&~C_{*}\bar{\varsigma}\|\vartheta_{c}\partial_{x}^{r+1}\tilde{\mathbb{U}}\|_{L_{\check{\psi}}^{2}(\Omega)}
       \|\mathcal{P}\varphi_{r}\|_{L_{x}^{\infty}(\mathbb{T})L_{y}^{2}([\check{y},\hat{y}])}\|\partial_{y}\mathbf{U}_{r}\|_{L_{\varphi_{r}}^{2}(\Omega)}\nonumber\\
\leq&~C_{*}\bar{\varsigma}^{\frac{3}{4}-\frac{\varepsilon_{0}}{2}}\|\vartheta_{c}\partial_{x}^{r+1}\tilde{\mathbb{U}}\|_{L_{\check{\psi}}^{2}(\Omega)}
                               \|\partial_{y}\mathbf{U}_{r}\|_{L_{\varphi_{r}}^{2}(\Omega)}\nonumber\\
\leq&~C_{*}\bar{\varsigma}^{\frac{3}{4}-\frac{\varepsilon_{0}}{2}}\|\mathcal{P}\omega_{o}^{-1}\|_{L^{\infty}(\Omega)}
           \left(\|\vartheta_{c}\partial_{x}\tilde{\mathbf{U}}_{r}\|_{L_{\tilde{\varphi}_{r}}^{2}(\Omega)}
            +\|\tilde{\mathcal{P}}\tilde{\mathbf{U}}_{r}\|_{L_{\tilde{\varphi}_{r}}^{2}(\Omega)}\right)
      \|\partial_{y}\mathbf{U}_{r}\|_{L_{\varphi_{r}}^{2}(\Omega)}\nonumber\\
\leq&~C_{*}\bar{\varsigma}^{\frac{3}{4}}\tilde{\bar{\varsigma}}^{-\frac{1}{2}}
           \left(\|\vartheta_{c}\partial_{x}\tilde{\mathbf{U}}_{r}\|_{L_{\tilde{\varphi}_{r}}^{2}(\Omega)}
            +\|\tilde{\mathcal{P}}\tilde{\mathbf{U}}_{r}\|_{L_{\tilde{\varphi}_{r}}^{2}(\Omega)}\right)
      \|\partial_{y}\mathbf{U}_{r}\|_{L_{\varphi_{r}}^{2}(\Omega)}\nonumber\\
\leq&~C_{*}\bar{\varsigma}^{\frac{3}{2}}\tilde{\bar{\varsigma}}^{-1}
           \|\vartheta_{c}\partial_{x}\tilde{\mathbf{U}}_{r}\|_{L_{\tilde{\varphi}_{r}}^{2}(\Omega)}^{2}
     +C_{*}\bar{\varsigma}^{\frac{3}{2}}\tilde{\bar{\varsigma}}^{-1}\|\tilde{\mathcal{P}}\tilde{\mathbf{U}}_{r}\|_{L_{\tilde{\varphi}_{r}}^{2}(\Omega)}^{2}
     +\frac{1}{16}\|\partial_{y}\mathbf{U}_{r}\|_{L_{\varphi_{r}}^{2}(\Omega)}^{2}\nonumber\\
\leq&~C_{*}\sqrt{\iota}\|\sqrt{\tilde{\eta}}\partial_{x}\tilde{\mathbf{U}}_{r}\|_{L_{\tilde{\varphi}_{r}}^{2}(\Omega)}^{2}
     +\frac{1}{16}\|\partial_{y}\mathbf{U}_{r}\|_{L_{\varphi_{r}}^{2}(\Omega)}^{2}
     +C_{*}\|\tilde{\mathbf{U}}_{r}\|_{L_{\tilde{\varphi}_{r}}^{2}(\Omega)}\|\tilde{\mathcal{P}}\tilde{\mathbf{U}}_{r}\|_{L_{\tilde{\varphi}_{r}}^{2}(\Omega)}\nonumber\\
\leq&~C_{*}\sqrt{\iota}\|\sqrt{\tilde{\eta}}\partial_{x}\tilde{\mathbf{U}}_{r}\|_{L_{\tilde{\varphi}_{r}}^{2}(\Omega)}^{2}
     +\frac{1}{16}\|\partial_{y}\mathbf{U}_{r}\|_{L_{\varphi_{r}}^{2}(\Omega)}^{2}
     +C_{*}C_{\tau}\|\tilde{\mathbf{U}}_{r}\|_{L_{\tilde{\varphi}_{r}}^{2}(\Omega)}^{2}
     +\frac{1}{2}\tau\|\tilde{\mathcal{P}}\tilde{\mathbf{U}}_{r}\|_{L_{\tilde{\varphi}_{r}}^{2}(\Omega)}^{2},\nonumber
\end{align}
With the help of \eqref{varphiy}, we can employ a similar estimate as in \eqref{J723} to obtain
\begin{align}\label{J724}
&4(\mathcal{P}\varsigma\partial_{y}^{-1}[\partial_{x}^{r+1}\tilde{\mathbb{U}}],\varphi_{r}\partial_{y}\varphi_{r}\mathbf{U}_{r})\\
\leq&~C_{*}\bar{\varsigma}^{\frac{3}{4}}\tilde{\bar{\varsigma}}^{-\frac{1}{2}}
      \left(\|\vartheta_{c}\partial_{x}\tilde{\mathbf{U}}_{r}\|_{L_{\tilde{\varphi}_{r}}^{2}(\Omega)}
            +\|\tilde{\mathcal{P}}\tilde{\mathbf{U}}_{r}\|_{L_{\tilde{\varphi}_{r}}^{2}(\Omega)}\right)
      \|\varphi_{r}^{-1}\partial_{y}\varphi_{r}^{2}\mathbf{U}_{r}\|_{L^{2}(\Omega)}\nonumber\\
\leq&~C_{*}\bar{\varsigma}^{\frac{3}{2}}\tilde{\bar{\varsigma}}^{-1}
           \|\vartheta_{c}\partial_{x}\tilde{\mathbf{U}}_{r}\|_{L_{\tilde{\varphi}_{r}}^{2}(\Omega)}^{2}
     +C_{*}\bar{\varsigma}^{\frac{3}{2}}\tilde{\bar{\varsigma}}^{-1}\|\tilde{\mathcal{P}}\tilde{\mathbf{U}}_{r}\|_{L_{\tilde{\varphi}_{r}}^{2}(\Omega)}^{2}
     +\frac{1}{16}\|\sqrt{\omega_{\lambda}}\mathbf{U}_{r}\|_{L_{\varphi_{r}}^{2}(\Omega)}^{2}
     +C\|\mathbf{U}_{r}\|_{L_{\varphi_{r}}^{2}(\Omega)}^{2}\nonumber\\
\leq&~C_{*}\sqrt{\iota}\|\sqrt{\tilde{\eta}}\partial_{x}\tilde{\mathbf{U}}_{r}\|_{L_{\tilde{\varphi}_{r}}^{2}(\Omega)}^{2}
     +\frac{1}{16}\|\sqrt{\omega_{\lambda}}\mathbf{U}_{r}\|_{L_{\varphi_{r}}^{2}(\Omega)}^{2}
     +C_{*}\|\mathbf{U}_{r}\|_{L_{\varphi_{r}}^{2}(\Omega)}^{2}\nonumber\\
    &+C_{*}C_{\tau}\|\tilde{\mathbf{U}}_{r}\|_{L_{\tilde{\varphi}_{r}}^{2}(\Omega)}^{2}
     +\frac{1}{2}\tau\|\tilde{\mathcal{P}}\tilde{\mathbf{U}}_{r}\|_{L_{\tilde{\varphi}_{r}}^{2}(\Omega)}^{2}.\nonumber
\end{align}
It remains to estimate the last four terms on the right hand side of \eqref{J72}. Noticing ${\bf supp}\varsigma \subset(\check{y}_{*},\hat{y}_{*})$ and \eqref{varphiy}, we have
\begin{align}\label{J725}
&\quad(\mathcal{P}\varsigma'\partial_{y}^{-1}[\partial_{x}^{r+1}u_{0}],\varphi_{r}^{2}\mathbf{U}_{r})
+2(\mathcal{P}\varsigma\partial_{x}^{r+1}u_{0},\varphi_{r}^{2}\mathbf{U}_{r})\\
&+2(\mathcal{P}\varsigma\partial_{y}^{-1}[\partial_{x}^{r+1}u_{0}],\varphi_{r}^{2}\partial_{y}\mathbf{U}_{r})
+2(\mathcal{P}\varsigma\partial_{y}^{-1}[\partial_{x}^{r+1}u_{0}],\partial_{y}\varphi_{r}^{2}\mathbf{U}_{r})\nonumber\\
\leq& C_{*}\|\mathbf{U}_{r}\|_{L_{\varphi_{r}}^{2}(\Omega)}^{2}
    +\frac{1}{16}\|\partial_{y}\mathbf{U}_{r}\|_{L_{\varphi_{r}}^{2}(\Omega)}^{2}
    +\frac{1}{16}\|\sqrt{\omega_{\lambda}}\mathbf{U}_{r}\|_{L_{\varphi_{r}}^{2}(\Omega)}^{2}
    +C_{*}.\nonumber
\end{align}

Substituting \eqref{J721},\eqref{J722},\eqref{J723},\eqref{J724} and \eqref{J725} into \eqref{J72}, we have
\begin{align}\label{J72'}
 -(\mathcal{P}\partial_{y}^{-1}[\partial_{x}^{r+1}\tilde{\mathfrak{U}}](\varsigma'-\rho\varsigma),\varphi_{r}^{2}\mathbf{U}_{r})
\leq& (2\tau+C_{*}\sqrt{\iota})\|\sqrt{\tilde{\eta}}\partial_{x}\tilde{\mathbf{U}}_{r}\|_{L_{\tilde{\varphi}_{r}}^{2}(\Omega)}^{2}
     +\frac{1}{8}\|\sqrt{\omega_{\lambda}}\mathbf{U}_{r}\|_{L_{\varphi_{r}}^{2}(\Omega)}^{2}\\
    &+\frac{1}{8}\|\partial_{y}\mathbf{U}_{r}\|_{L_{\varphi_{r}}^{2}(\Omega)}^{2}
     +C_{*}C_{\tau}\|\mathbf{U}_{r}\|_{L_{\varphi_{r}}^{2}(\Omega)}^{2}
     +C_{*}C_{\tau}\|\tilde{\mathbf{U}}_{r}\|_{L_{\tilde{\varphi}_{r}}^{2}(\Omega)}^{2}\nonumber\\
    &+\frac{1}{2}\|\mathcal{P}\mathbf{U}_{r}\|_{L_{\varphi_{r}}^{2}(\Omega)}^{2}
     +\tau\|\tilde{\mathcal{P}}\tilde{\mathbf{U}}_{r}\|_{L_{\tilde{\varphi}_{r}}^{2}(\Omega)}^{2}+C_{*}.\nonumber
\end{align}

Now combining \eqref{J71'} and \eqref{J72'} with \eqref{J7}, we can conclude that
\begin{align*}
J_{7}
\leq& \left(\frac{1}{8}+C_{\lambda}\theta\sqrt{t}\right)\|\sqrt{\eta}\partial_{x}\mathbf{U}_{r}\|_{L_{\varphi_{r}}^{2}(\Omega)}^{2}
     +\frac{1}{8}\|\sqrt{\omega_{\lambda}}\mathbf{U}_{r}\|_{L_{\varphi_{r}}^{2}(\Omega)}^{2}
     +\frac{1}{8}\|\partial_{y}\mathbf{U}_{r}\|_{L_{\varphi_{r}}^{2}(\Omega)}^{2}
     +\|\mathcal{P}\mathbf{U}_{r}\|_{L_{\varphi_{r}}^{2}(\Omega)}^{2}\\
    &+\tau\|\tilde{\mathcal{P}}\tilde{\mathbf{U}}_{r}\|_{L_{\tilde{\varphi}_{r}}^{2}(\Omega)}^{2}
     +C_{*}C_{\tau}C_{\theta}\iota^{-1}\|\mathbf{U}_{r}\|_{L_{\varphi_{r}}^{2}(\Omega)}^{2}
     +(2\tau+C_{*}\sqrt{\iota})\|\sqrt{\tilde{\eta}}\partial_{x}\tilde{\mathbf{U}}_{r}\|_{L_{\varphi_{r}}^{2}(\Omega)}^{2}
     +C_{*}C_{\tau}\|\tilde{\mathbf{U}}_{r}\|_{L_{\tilde{\varphi}_{r}}^{2}(\Omega)}^{2}+C_{*}\nonumber.
\end{align*}

\noindent{\bf \underline{Estimate of~$J_{8}$.}}
Taking notice of the fact that $\partial_{x}\partial_{y}^{2}\mathfrak{U}-\rho\partial_{x}\partial_{y}\mathfrak{U}=\partial_{x}\rho(\dot{\mathfrak{U}}+\varsigma)$,
and $|\partial_{x}\rho|\leq C\mathcal{P}$, we have
\begin{align*}
J_{8}
=&~r\left(\mathcal{P}\left(\partial_{x}\partial_{y}^{2}\mathfrak{U}-\rho\partial_{x}\partial_{y}\mathfrak{U}\right)\partial_{y}^{-1}[\partial_{x}^{r}\mathfrak{U}],\varphi_{r}^{2}\mathbf{U}_{r}\right)\\
=&~r\left((\dot{\mathfrak{U}}+\varsigma)\mathcal{P}\partial_{x}\rho\partial_{y}^{-1}[\partial_{x}^{r}\mathbb{U}],\varphi_{r}^{2}\mathbf{U}_{r}\right)
  +r\left((\dot{\mathfrak{U}}+\varsigma)\mathcal{P}\partial_{x}\rho\partial_{y}^{-1}[\partial_{x}^{r}u_{0}],\varphi_{r}^{2}\mathbf{U}_{r}\right)\\
\leq&~C_{*}\|\varphi_{r}(\dot{\mathfrak{U}}+\varsigma)\mathcal{P}^{2}\|_{L_{x}^{\infty}(\mathbb{T})L_{y}^{2}(\mathbb{R_{+}})}
       \|\partial_{x}^{r}\mathbb{U}\|_{L_{\varphi_{r}}^{2}(\Omega)}\|\mathbf{U}_{r}\|_{L_{\varphi_{r}}^{2}(\Omega)}\\
    &+C_{*}\|\varphi_{r}\sqrt{\y}(\dot{\mathfrak{U}}+\varsigma)\mathcal{P}^{2}\|_{L_{x}^{\infty}(\mathbb{T})L_{y}^{2}(\mathbb{R_{+}})}
           \|\mathbf{U}_{r}\|_{L_{\varphi_{r}}^{2}(\Omega)}\\
\leq&~C_{*}\|\partial_{x}^{r}\mathbb{U}\|_{L_{\varphi_{r}}^{2}(\Omega)}\|\mathbf{U}_{r}\|_{L_{\varphi_{r}}^{2}(\Omega)}
     +C_{*}\left(\|\sqrt{\y}\mathbb{U}\|_{\hat{H}_{\psi}^{s}(\Omega)}+1\right)\|\mathbf{U}_{r}\|_{L_{\varphi_{r}}^{2}(\Omega)}\\
\leq&~C_{*}\|\mathcal{P}\mathbf{U}_{r}\|_{L_{\varphi_{r}}^{2}(\Omega)}\|\mathbf{U}_{r}\|_{L_{\varphi_{r}}^{2}(\Omega)}
     +C_{*}\left(\|\sqrt{\y}\mathbb{U}\|_{\hat{H}_{\psi}^{s}(\Omega)}+1\right)\|\mathbf{U}_{r}\|_{L_{\varphi_{r}}^{2}(\Omega)}\\
\leq&~\|\mathcal{P}\mathbf{U}_{r}\|_{L_{\varphi_{r}}^{2}(\Omega)}^{2}
     +C_{*}C_{\tau}\|\mathbf{U}_{r}\|_{L_{\varphi_{r}}^{2}(\Omega)}^{2}
     +\tau\|\sqrt{\y}\mathbb{U}\|_{\hat{H}_{\psi}^{s}(\Omega)}^{2}
     +C_{*}.
\end{align*}

\noindent{\bf \underline{Estimate of~$J_{9}$.}}
We divide $J_{9}$ into four parts,
\begin{align}\label{J9}
J_{9}
=&\left((\partial_{y}^{2}\hat{F}-\rho\partial_{y}\hat{F})\mathcal{P}\frac{\partial_{x}^{r}\mathbb{U}}{\dot{\mathfrak{U}}+\varsigma},\varphi_{r}^{2}\mathbf{U}_{r}\right)
  -2\left(\eta\mathcal{P}\partial_{x}\partial_{y}\mathbb{U}\partial_{x}\rho\frac{\partial_{x}^{r}\mathbb{U}}{\dot{\mathfrak{U}}+\varsigma},\varphi_{r}^{2}\mathbf{U}_{r}\right)\\
 &+\left(\mathcal{P}\left(\varsigma'\partial_{x}\mathfrak{U}+\varsigma'''+\partial_{y}^{-1}[\partial_{x}\mathfrak{U}]\varsigma''-\varsigma\partial_{x}\partial_{y}\mathfrak{U}+\eta\partial_{x}^{2}\partial_{y}^{2}u_{0}\right)\frac{\partial_{x}^{r}\mathbb{U}}{\dot{\mathfrak{U}}+\varsigma},\varphi_{r}^{2}\mathbf{U}_{r}\right)\nonumber\\
 &-\left(\mathcal{P}\rho\left(\partial_{y}^{-1}[\partial_{x}\mathfrak{U}]\varsigma'+\varsigma''+\eta\partial_{x}^{2}\partial_{y}u_{0}\right)\frac{\partial_{x}^{r}\mathbb{U}}{\dot{\mathfrak{U}}+\varsigma},\varphi_{r}^{2}\mathbf{U}_{r}\right).\nonumber
\end{align}

Since $\|\mathcal{P}(\partial_{y}^{2}\hat{F}-\rho\partial_{y}\hat{F})\|_{L^{\infty}(\Omega_{t_{3}})}\leq C_{*}$,
and for any $\gamma\in\Gamma_{2}, |D^{\gamma}\hat{F}|\leq C_{*}e^{-\frac{4}{3}\delta\y}~\text{in}~\Omega_{T_{*}},$
with Corollary \ref{ncC2} one can deduce that
\begin{align}\label{J91}
\left((\partial_{y}^{2}\hat{F}-\rho\partial_{y}\hat{F})\mathcal{P}\frac{\partial_{x}^{r}\mathbb{U}}{\dot{\mathfrak{U}}+\varsigma},\varphi_{r}^{2}\mathbf{U}_{r}\right)
\leq& C_{*}\|\partial_{x}^{r}\mathbb{U}\|_{L_{\hat{\Psi}_{r}}^{2}(\Omega)}\|\mathbf{U}_{r}\|_{L_{\varphi_{r}}^{2}(\Omega)}\\
\leq& C_{*}\|\mathcal{P}\mathbf{U}_{r}\|_{L_{\varphi_{r}}^{2}(\Omega)}\|\mathbf{U}_{r}\|_{L_{\varphi_{r}}^{2}(\Omega)}\nonumber\\
\leq& C_{*}\|\mathbf{U}_{r}\|_{L_{\varphi_{r}}^{2}(\Omega)}^{2}
     +\frac{1}{3}\|\mathcal{P}\mathbf{U}_{r}\|_{L_{\varphi_{r}}^{2}(\Omega)}^{2}.\nonumber
\end{align}
Due to Corollary \ref{wdotuC'} and Lemma \ref{rhoP}, moreover, $\partial_{x}u_{0}=0$ for $y\in[0,\hat{y}]$, one has $|\vartheta_{c}^{2}\partial_{x}\partial_{y}\mathbb{U}|\leq C\dot{\mathfrak{U}}$ in $\Omega_{t_{3}}$, which leads to
\begin{align}\label{J92}
-2\left(\eta\mathcal{P}\partial_{x}\partial_{y}\mathbb{U}\partial_{x}\rho\frac{\partial_{x}^{r}\mathbb{U}}{\dot{\mathfrak{U}}+\varsigma},\varphi_{r}^{2}\mathbf{U}_{r}\right)
\leq& C_{*}\bar{\eta}\|\partial_{x}^{r}\mathbb{U}\|_{L_{\hat{\Psi}_{r}}^{2}(\Omega)}\|\mathbf{U}_{r}\|_{L_{\varphi_{r}}^{2}(\Omega)}\\
\leq& C_{*}\bar{\eta}\|\mathcal{P}\mathbf{U}_{r}\|_{L_{\varphi_{r}}^{2}(\Omega)}\|\mathbf{U}_{r}\|_{L_{\varphi_{r}}^{2}(\Omega)}\nonumber\\
\leq& C_{*}\|\mathbf{U}_{r}\|_{L_{\varphi_{r}}^{2}(\Omega)}^{2}
     +\frac{1}{3}\|\mathcal{P}\mathbf{U}_{r}\|_{L_{\varphi_{r}}^{2}(\Omega)}^{2},\nonumber
\end{align}
Notice that
$$\varsigma'\partial_{x}\mathfrak{U}+\varsigma'''+\partial_{y}^{-1}[\partial_{x}\mathfrak{U}]\varsigma''-\varsigma\partial_{x}\partial_{y}\mathfrak{U}+\eta\partial_{x}^{2}\partial_{y}^{2}u_{0}
\leq C_{*}\iota^{-\frac{1}{2}}\sqrt{\bar{\varsigma}},$$
and
$$\rho\left(\partial_{y}^{-1}[\partial_{x}\mathfrak{U}]\varsigma'+\varsigma''+\eta\partial_{x}^{2}\partial_{y}u_{0}\right)
\leq C_{*}\iota^{-\frac{1}{2}}\sqrt{\bar{\varsigma}},$$
also by Corollary \ref{ncC2} one gets
\begin{align}\label{J93}
&\left(\mathcal{P}\left(\varsigma'\partial_{x}\mathfrak{U}+\varsigma'''+\partial_{y}^{-1}[\partial_{x}\mathfrak{U}]\varsigma''-\varsigma\partial_{x}\partial_{y}\mathfrak{U}+\eta\partial_{x}^{2}\partial_{y}^{2}u_{0}\right)\frac{\partial_{x}^{r}\mathbb{U}}{\dot{\mathfrak{U}}+\varsigma},\varphi_{r}^{2}\mathbf{U}_{r}\right)\\
&-\left(\mathcal{P}\rho\left(\partial_{y}^{-1}[\partial_{x}\mathfrak{U}]\varsigma'+\varsigma''+\eta\partial_{x}^{2}\partial_{y}u_{0}\right)\frac{\partial_{x}^{r}\mathbb{U}}{\dot{\mathfrak{U}}+\varsigma},\varphi_{r}^{2}\mathbf{U}_{r}\right)\nonumber\\
\leq& C_{*}\iota^{-\frac{1}{2}}\sqrt{\bar{\varsigma}}\|\partial_{x}^{r}\mathbb{U}\|_{L_{\hat{\Psi}_{r}}^{2}(\Omega)}
                               \|\mathcal{P}\mathbf{U}_{r}\|_{L_{\varphi_{r}}^{2}(\Omega)}\nonumber\\
\leq& C_{*}\iota^{-\frac{1}{2}}\sqrt{\bar{\varsigma}}\|\mathcal{P}\mathbf{U}_{r}\|_{L_{\varphi_{r}}^{2}(\Omega)}^{2}\nonumber\\
\leq& C_{*}\iota^{-1}\|\mathbf{U}_{r}\|_{L_{\varphi_{r}}^{2}(\Omega)}^{2}
     +\frac{1}{3}\|\mathcal{P}\mathbf{U}_{r}\|_{L_{\varphi_{r}}^{2}(\Omega)}^{2},\nonumber
\end{align}

A combination of \eqref{J91}, \eqref{J92} and \eqref{J93} with \eqref{J9} gives
\begin{align*}
J_{9}
\leq& C_{*}\iota^{-1}\|\mathbf{U}_{r}\|_{L_{\varphi_{r}}^{2}(\Omega)}^{2}
     +\|\mathcal{P}\mathbf{U}_{r}\|_{L_{\varphi_{r}}^{2}(\Omega)}^{2}.
\end{align*}

\noindent{\bf \underline{Estimate of~$J_{10}$.}}
Recall that
\begin{align}\label{J10}
J_{10}
= (\mathcal{P}\varsigma\partial_{x}^{r+1}u,\varphi_{r}^{2}\mathbf{U}_{r})
 +(\mathcal{P}(\rho\partial_{x}^{r}\mathbb{F}-\partial_{y}\partial_{x}^{r}\mathbb{F}),\varphi_{r}^{2}\mathbf{U}_{r})
\end{align}
In view of $u=\mathbb{U}-\tilde{\mathbb{U}}$, one can divide the first term on the right hand side of \eqref{J10} into two parts,
\begin{align}\label{J101}
(\mathcal{P}\varsigma\partial_{x}^{r+1}u,\varphi_{r}^{2}\mathbf{U}_{r})
=& (\mathcal{P}\varsigma\partial_{x}^{r+1}\mathbb{U},\varphi_{r}^{2}\mathbf{U}_{r})
  -(\mathcal{P}\varsigma\partial_{x}^{r+1}\tilde{\mathbb{U}},\varphi_{r}^{2}\mathbf{U}_{r})
\end{align}
As a result of Corollary \ref{ncC3}, it follows
\begin{align}\label{J1011}
(\mathcal{P}\varsigma\partial_{x}^{r+1}\mathbb{U},\varphi_{r}^{2}\mathbf{U}_{r})
\leq& C\bar{\varsigma}\|\vartheta_{c}\mathcal{P}\partial_{x}^{r+1}\mathbb{U}\|_{L_{\varphi_{r}}^{2}(\Omega)}
       \|\mathbf{U}_{r}\|_{L_{\varphi_{r}}^{2}(\Omega)}\\
\leq& C_{*}\bar{\varsigma}\|\vartheta_{c}\partial_{x}^{r+1}\mathbb{U}\|_{L_{\hat{\Psi}_{r}}^{2}(\Omega)}
       \|\mathbf{U}_{r}\|_{L_{\varphi_{r}}^{2}(\Omega)}\nonumber\\
\leq& C_{*}\sqrt{\bar{\varsigma}}
      \left(\|\vartheta_{c}\partial_{x}\mathbf{U}_{r}\|_{L_{\varphi_{r}}^{2}(\Omega)}
            +\|\mathcal{P}\mathbf{U}_{r}\|_{L_{\varphi_{r}}^{2}(\Omega)}\right)
      \|\mathbf{U}_{r}\|_{L_{\varphi_{r}}^{2}(\Omega)}\nonumber\\
\leq& \frac{1}{16}\|\sqrt{\eta}\partial_{x}\mathbf{U}_{r}\|_{L_{\varphi_{r}}^{2}(\Omega)}^{2}
     +C_{*}\|\mathbf{U}_{r}\|_{L_{\varphi_{r}}^{2}(\Omega)}^{2}.\nonumber
\end{align}
Combining \eqref{J1011} and \eqref{J722} with \eqref{J101}, one has
\begin{align}\label{J101'}
(\mathcal{P}\varsigma\partial_{x}^{r+1}\mathbb{U},\varphi_{r}^{2}\mathbf{U}_{r})
\leq& \frac{1}{16}\|\sqrt{\eta}\partial_{x}\mathbf{U}_{r}\|_{L_{\varphi_{r}}^{2}(\Omega)}^{2}
     +C_{*}C_{\tau}\|\mathbf{U}_{r}\|_{L_{\varphi_{r}}^{2}(\Omega)}^{2}\\
    &+\tau\|\sqrt{\tilde{\eta}}\partial_{x}\tilde{\mathbf{U}}_{r}\|_{L_{\tilde{\varphi}_{r}}^{2}(\Omega)}^{2}
     +C_{*}\|\tilde{\mathbf{U}}_{r}\|_{L_{\tilde{\varphi}_{r}}^{2}(\Omega)}^{2}.\nonumber
\end{align}

Next let's turn to the estimate of the second term on the right hand side of \eqref{J10}. By a direct calculation, we decompose $\rho\partial_{x}^{r}\mathbb{F}-\partial_{y}\partial_{x}^{r}\mathbb{F}$ into four parts according to the order of the derivatives of $u$ with respect to $x$ as follows,
\begin{align*}
\rho\partial_{x}^{r}\mathbb{F}-\partial_{y}\partial_{x}^{r}\mathbb{F}
= \partial_{y}\partial_{x}^{r}\left(u\partial_{x}u-(1-h)\partial_{y}^{-1}[\partial_{x}u]\partial_{y}u\right)
  -\rho\partial_{x}^{r}\left(u\partial_{x}u-(1-h)\partial_{y}^{-1}[\partial_{x}u]\partial_{y}u\right)
=: \sum_{i=1}^{4}\mathbb{J}_{i},
\end{align*}
where
\begin{align*}
\mathbb{J}_{1}
=& u\partial_{x}^{r+1}\partial_{y}u
  -\rho u\partial_{x}^{r+1}u
  +(1-h)\partial_{y}^{-1}[\partial_{x}^{r+1}u](\rho\partial_{y}u-\partial_{y}^{2}u)\\
 &+h'\partial_{y}^{-1}[\partial_{x}^{r+1}u]\partial_{y}u
  +h\partial_{x}^{r+1}u\partial_{y}u,
\end{align*}
\begin{align*}
\mathbb{J}_{2}
=& \partial_{x}^{r}u(\partial_{x}\partial_{y}u-\rho\partial_{x}u)
  +r\partial_{x}u(\partial_{x}^{r}\partial_{y}u-\rho\partial_{x}^{r}u)
  -r(1-h)\partial_{y}^{-1}[\partial_{x}^{r}u](\partial_{x}\partial_{y}^{2}u-\rho\partial_{x}\partial_{y}u)\\
 &-(1-h)\partial_{y}^{-1}[\partial_{x}u](\partial_{x}^{r}\partial_{y}^{2}u-\rho\partial_{x}^{r}\partial_{y}u)
  +h'(r\partial_{y}^{-1}[\partial_{x}^{r}u]\partial_{x}\partial_{y}u+\partial_{y}^{-1}[\partial_{x}u]\partial_{x}^{r}\partial_{y}u)\\
 &+h(r\partial_{x}^{r}u\partial_{x}\partial_{y}u+\partial_{x}u\partial_{x}^{r}\partial_{y}u),
\end{align*}
\begin{align*}
\mathbb{J}_{3}
=& r\partial_{x}^{r-1}u(\partial_{x}^{2}\partial_{y}u-\rho\partial_{x}^{2}u)
  +\frac{r(r-1)}{2}\partial_{x}^{2}u(\partial_{x}^{r-1}\partial_{y}u-\rho\partial_{x}^{r-1}u)\\
 &-\frac{r(r-1)}{2}(1-h)\partial_{y}^{-1}[\partial_{x}^{r-1}u](\partial_{x}^{2}\partial_{y}^{2}u-\rho\partial_{x}^{2}\partial_{y}u)
  -r(1-h)\partial_{y}^{-1}[\partial_{x}^{2}u](\partial_{x}^{r-1}\partial_{y}^{2}u-\rho\partial_{x}^{r-1}\partial_{y}u)\\
 &+rh'\partial_{y}^{-1}[\partial_{x}^{2}u]\partial_{x}^{r-1}\partial_{y}u
  +\frac{r(r-1)}{2}h'\partial_{y}^{-1}[\partial_{x}^{r-1}u]\partial_{x}^{2}\partial_{y}u
  +rh\partial_{x}^{2}u\partial_{x}^{r-1}\partial_{y}u
  +\frac{r(r-1)}{2}h\partial_{x}^{r-1}u\partial_{x}^{2}\partial_{y}u,
\end{align*}
and
\begin{align*}
\mathbb{J}_{4}
=& \sum_{i=3}^{r-2}\binom{r}{i}[\partial_{x}^{i}u\partial_{x}^{r-i+1}\partial_{y}u-\rho\partial_{x}^{i}u\partial_{x}^{r-i+1}u]\\
 &-(1-h)\sum_{i=2}^{r-3}\binom{r}{i}\left[\partial_{y}^{-1}[\partial_{x}^{i+1}u]\partial_{x}^{r-i}\partial_{y}^{2}u-\rho\partial_{y}^{-1}[\partial_{x}^{i+1}u]\partial_{x}^{n-i}\partial_{y}u\right]\\
 &+h\sum_{i=2}^{r-3}\binom{r}{i}\partial_{x}^{i+1}u\partial_{x}^{r-i}\partial_{y}u
  +h'\sum_{i=2}^{r-3}\binom{r}{i}\partial_{y}^{-1}[\partial_{x}^{i+1}u]\partial_{x}^{r-i}\partial_{y}u
\end{align*}
Accordingly, one has
\begin{align*}
(\mathcal{P}(\rho\partial_{x}^{r}\mathbb{F}-\partial_{y}\partial_{x}^{r}\mathbb{F}),\varphi_{r}^{2}\mathbf{U}_{r})
= \sum_{i=1}^{4}(\mathcal{P}\mathbb{J}_{i},\varphi_{r}^{2}\mathbf{U}_{r})
=: \sum_{i=1}^{4}\mathfrak{J}_{i}.
\end{align*}
Let's start with the estimate of $\mathfrak{J}_{1}$. We divide $\mathfrak{J}_{1}$ into three parts,
\begin{align*}
\mathfrak{J}_{1}
&= (\mathcal{P}(u\partial_{x}^{r+1}\partial_{y}u-\rho u\partial_{x}^{r+1}u),\varphi_{r}^{2}\mathbf{U}_{r})
  +(\mathcal{P}(1-h)\partial_{y}^{-1}[\partial_{x}^{r+1}u](\rho\partial_{y}u-\partial_{y}^{2}u),\varphi_{r}^{2}\mathbf{U}_{r})\\
&\quad+(\mathcal{P}(h'\partial_{y}^{-1}[\partial_{x}^{r+1}u]\partial_{y}u+h\partial_{x}^{r+1}u\partial_{y}u),\varphi_{r}^{2}\mathbf{U}_{r})\\
&=: \mathfrak{J}_{11}+\mathfrak{J}_{12}+\mathfrak{J}_{13}.
\end{align*}
By a direct calculation, one has
\begin{align}\label{JJ11}
\mathfrak{J}_{11}
=& (\mathcal{P}u(\partial_{x}^{r+1}\partial_{y}\mathbb{U}-\rho \partial_{x}^{r+1}\mathbb{U}),\varphi_{r}^{2}\mathbf{U}_{r})
  -(\mathcal{P}u(\partial_{x}^{r+1}\partial_{y}\tilde{\mathbb{U}}-\rho \partial_{x}^{r+1}\tilde{\mathbb{U}}),\varphi_{r}^{2}\mathbf{U}_{r})\\
=& (\mathcal{P}u\partial_{x}(\mathcal{P}^{-1}\mathbf{U}_{r}),\varphi_{r}^{2}\mathbf{U}_{r})
  +(\mathcal{P}u\partial_{x}\rho\partial_{x}^{r}\mathbb{U},\varphi_{r}^{2}\mathbf{U}_{r})
  -(\mathcal{P}u(\partial_{x}^{r+1}\partial_{y}\tilde{\mathbb{U}}-\tilde{\rho}\partial_{x}^{r+1}\tilde{\mathbb{U}}),\varphi_{r}^{2}\mathbf{U}_{r})\nonumber\\
 &+(\mathcal{P}u(\rho-\tilde{\rho})\partial_{x}^{r+1}\tilde{\mathbb{U}},\varphi_{r}^{2}\mathbf{U}_{r})\nonumber\\
=& (\mathcal{P}u\partial_{x}(\mathcal{P}^{-1}\mathbf{U}_{r}),\varphi_{r}^{2}\mathbf{U}_{r})
  -(\mathcal{P}u\partial_{x}(\tilde{\mathcal{P}}^{-1}\tilde{\mathbf{U}}_{r}),\varphi_{r}^{2}\mathbf{U}_{r})
  +(\mathcal{P}u\partial_{x}\rho\partial_{x}^{r}\mathbb{U},\varphi_{r}^{2}\mathbf{U}_{r})\nonumber\\
 &-(\mathcal{P}u\partial_{x}\tilde{\rho}\partial_{x}^{r}\tilde{\mathbb{U}},\varphi_{r}^{2}\mathbf{U}_{r})
  +(\mathcal{P}u(\rho-\tilde{\rho})\partial_{x}^{r+1}\tilde{\mathbb{U}},\varphi_{r}^{2}\mathbf{U}_{r}).\nonumber
\end{align}
Utilizing \eqref{us-3}, and recalling that $\bar{\eta}=\epsilon_{0}^{2}$ when $\check{c}\neq0$, one has
\begin{align}\label{JJ111}
 &(\mathcal{P}u\partial_{x}(\mathcal{P}^{-1}\mathbf{U}_{r}),\varphi_{r}^{2}\mathbf{U}_{r})
-(\mathcal{P}u\partial_{x}(\tilde{\mathcal{P}}^{-1}\tilde{\mathbf{U}}_{r}),\varphi_{r}^{2}\mathbf{U}_{r})\\
=& (u\partial_{x}\mathbf{U}_{r},\varphi_{r}^{2}\mathbf{U}_{r})
  -(\mathcal{P}\tilde{\mathcal{P}}^{-1}u\partial_{x}\tilde{\mathbf{U}}_{r},\varphi_{r}^{2}\mathbf{U}_{r})
  +(\mathcal{P}u\partial_{x}\mathcal{P}^{-1}\mathbf{U}_{r},\varphi_{r}^{2}\mathbf{U}_{r})
  -(\mathcal{P}u\partial_{x}\tilde{\mathcal{P}}^{-1}\tilde{\mathbf{U}}_{r},\varphi_{r}^{2}\mathbf{U}_{r})\nonumber\\
\leq& C_{*}(\check{c}+\sqrt{\bar{\eta}}t)\|\vartheta_{c}\partial_{x}\mathbf{U}_{r}\|_{L_{\varphi_{r}}^{2}(\Omega)}
                                             \|\mathbf{U}_{r}\|_{L_{\varphi_{r}}^{2}(\Omega)}
     +C_{*}(\check{c}+\sqrt{\bar{\eta}}t)\|\vartheta_{c}\partial_{x}\tilde{\mathbf{U}}_{r}\|_{L_{\tilde{\varphi}_{r}}^{2}(\Omega)}\|\mathbf{U}_{r}\|_{L_{\varphi_{r}}^{2}(\Omega)}\nonumber\\
    &+C_{*}(\check{c}+\sqrt{\bar{\eta}}t)\|\mathcal{P}\|_{L^{\infty}(\Omega)}\|\mathbf{U}_{r}\|_{L_{\varphi_{r}}^{2}(\Omega)}^{2}
     +C_{*}(\check{c}+\sqrt{\bar{\eta}}t)\|\mathcal{P}\|_{L^{\infty}(\Omega)}\|\mathbf{U}_{r}\|_{L_{\varphi_{r}}^{2}(\Omega)}
                                             \|\tilde{\mathbf{U}}_{r}\|_{L_{\tilde{\varphi}_{r}}^{2}(\Omega)}\nonumber\\
\leq& \frac{1}{16}\|\sqrt{\eta}\partial_{x}\mathbf{U}_{r}\|_{L_{\varphi_{r}}^{2}(\Omega)}^{2}
     +\tau\|\sqrt{\tilde{\eta}}\partial_{x}\tilde{\mathbf{U}}_{r}\|_{L_{\varphi_{r}}^{2}(\Omega)}^{2}
     +C_{*}C_{\tau}\iota^{-1}\|\mathbf{U}_{r}\|_{L_{\varphi_{r}}^{2}(\Omega)}^{2}
     +\|\tilde{\mathbf{U}}_{r}\|_{L_{\tilde{\varphi}_{r}}^{2}(\Omega)}^{2}.\nonumber
\end{align}
By Corollary \ref{ncC2}, one can obtain
\begin{align}\label{JJ112}
&(\mathcal{P}u\partial_{x}\rho\partial_{x}^{r}\mathbb{U},\varphi_{r}^{2}\mathbf{U}_{r})
-(\mathcal{P}u\partial_{x}\tilde{\rho}\partial_{x}^{r}\tilde{\mathbb{U}},\varphi_{r}^{2}\mathbf{U}_{r})\\
\leq& C_{*}(\check{c}+\sqrt{\bar{\eta}}t)\|\partial_{x}^{r}\mathbb{U}\|_{L_{\hat{\Psi}_{r}}^{2}(\Omega)}
                                             \|\mathbf{U}_{r}\|_{L_{\varphi_{r}}^{2}(\Omega)}
     +C_{*}(\check{c}+\sqrt{\bar{\eta}}t)\|\partial_{x}^{r}\tilde{\mathbb{U}}\|_{L_{\tilde{\hat{\Psi}}_{r}}^{2}(\Omega)}
                                             \|\mathbf{U}_{r}\|_{L_{\varphi_{r}}^{2}(\Omega)}\nonumber\\
\leq& C_{*}(\check{c}+\sqrt{\bar{\eta}}t)\|\mathcal{P}\mathbf{U}_{r}\|_{L_{\varphi_{r}}^{2}(\Omega)}
                                             \|\mathbf{U}_{r}\|_{L_{\varphi_{r}}^{2}(\Omega)}
     +C_{*}(\check{c}+\sqrt{\bar{\eta}}t)\|\tilde{\mathcal{P}}\tilde{\mathbf{U}}_{r}\|_{L_{\tilde{\varphi}_{r}}^{2}(\Omega)}
                                             \|\mathbf{U}_{r}\|_{L_{\varphi_{r}}^{2}(\Omega)}\nonumber\\
\leq& C_{*}\iota^{-1}\|\mathbf{U}_{r}\|_{L_{\varphi_{r}}^{2}(\Omega)}^{2}
     +\|\tilde{\mathbf{U}}_{r}\|_{L_{\tilde{\varphi}_{r}}^{2}(\Omega)}^{2}.\nonumber
\end{align}
Using Corollary \ref{ncC3} and the fact $\check{c}\leq C\iota^{-1}\bar{\varsigma}$, one can deduce
\begin{align}\label{JJ113}
&(\mathcal{P}u(\rho-\tilde{\rho})\partial_{x}^{r+1}\tilde{\mathbb{U}},\varphi_{r}^{2}\mathbf{U}_{r})\\
\leq&~C_{*}(\check{c}+\sqrt{\bar{\eta}}t)
           \|\vartheta_{c}\partial_{x}^{r+1}\tilde{\mathbb{U}}\|_{L_{\tilde{\hat{\Psi}}_{r}}^{2}(\Omega)}
           \|\mathbf{U}_{r}\|_{L_{\varphi_{r}}^{2}(\Omega)}\nonumber\\
\leq&~C_{*}(\check{c}+\sqrt{\bar{\eta}}t)\|\tilde{\mathcal{P}}\|_{L^{\infty}(\Omega)}
           \left(\|\vartheta_{c}\partial_{x}\tilde{\mathbf{U}}_{r}\|_{L_{\tilde{\varphi}_{r}}^{2}(\Omega)}
                 +\|\tilde{\mathcal{P}}\tilde{\mathbf{U}}_{r}\|_{L_{\tilde{\varphi}_{r}}^{2}(\Omega)}\right)
      \|\mathbf{U}_{r}\|_{L_{\varphi_{r}}^{2}(\Omega)}\nonumber\\
\leq&~C_{*}(\iota^{-1}\sqrt{\bar{\varsigma}}+\sqrt{\bar{\eta}t})
           \|\vartheta_{c}\partial_{x}\tilde{\mathbf{U}}_{r}\|_{L_{\tilde{\varphi}_{r}}^{2}(\Omega)}\|\mathbf{U}_{r}\|_{L_{\varphi_{r}}^{2}(\Omega)}\nonumber\\
    &+C_{*}(\check{c}+\sqrt{\bar{\eta}}t)\|\tilde{\mathcal{P}}\|_{L^{\infty}(\Omega)}
           \|\tilde{\mathcal{P}}\tilde{\mathbf{U}}_{r}\|_{L_{\tilde{\varphi}_{r}}^{2}(\Omega)}\|\mathbf{U}_{r}\|_{L_{\varphi_{r}}^{2}(\Omega)}\nonumber\\
\leq& \tau\|\sqrt{\tilde{\eta}}\partial_{x}\tilde{\mathbf{U}}_{r}\|_{L_{\tilde{\varphi}_{r}}^{2}(\Omega)}^{2}
     +C_{*}C_{\tau}(\iota^{-2}+t)\|\mathbf{U}_{r}\|_{L_{\varphi_{r}}^{2}(\Omega)}^{2}\nonumber\\
    &+C_{*}\iota^{-1}\|\tilde{\mathbf{U}}_{r}\|_{L_{\tilde{\varphi}_{r}}^{2}(\Omega)}\|\mathbf{U}_{r}\|_{L_{\varphi_{r}}^{2}(\Omega)}\nonumber\\
\leq& \tau\|\sqrt{\tilde{\eta}}\partial_{x}\tilde{\mathbf{U}}_{r}\|_{L_{\tilde{\varphi}_{r}}^{2}(\Omega)}^{2}
     +\|\tilde{\mathbf{U}}_{r}\|_{L_{\tilde{\varphi}_{r}}^{2}(\Omega)}^{2}
     +C_{*}C_{\tau}\iota^{-2}\|\mathbf{U}_{r}\|_{L_{\varphi_{r}}^{2}(\Omega)}^{2}.\nonumber
\end{align}
Substituting \eqref{JJ111}, \eqref{JJ112} and \eqref{JJ113} into \eqref{JJ11}, it follows
\begin{align}\label{JJ11'}
\mathfrak{J}_{11}
\leq \frac{1}{16}\|\sqrt{\eta}\partial_{x}\mathbf{U}_{r}\|_{L_{\varphi_{r}}^{2}(\Omega)}^{2}
    +2\tau\|\sqrt{\tilde{\eta}}\partial_{x}\tilde{\mathbf{U}}_{r}\|_{L_{\varphi_{r}}^{2}(\Omega)}^{2}
    +C_{*}C_{\tau}\iota^{-2}\|\mathbf{U}_{r}\|_{L_{\varphi_{r}}^{2}(\Omega)}^{2}
    +C\|\tilde{\mathbf{U}}_{r}\|_{L_{\tilde{\varphi}_{r}}^{2}(\Omega)}^{2}.
\end{align}
The term $\mathfrak{J}_{12}$ should be divided into two parts,
\begin{align}\label{JJ12}
\mathfrak{J}_{12}
= (\mathcal{P}(1-h)\partial_{y}^{-1}[\partial_{x}^{r+1}\mathbb{U}](\rho\partial_{y}u-\partial_{y}^{2}u),\varphi_{r}^{2}\mathbf{U}_{r})
 -(\mathcal{P}(1-h)\partial_{y}^{-1}[\partial_{x}^{r+1}\tilde{\mathbb{U}}](\rho\partial_{y}u-\partial_{y}^{2}u),\varphi_{r}^{2}\mathbf{U}_{r})
\end{align}
For the first part on the right hand side of \eqref{JJ12}, since
$\|\rho\|_{L^{\infty}(\Omega)}\leq C$ in $\mathbb{R}_{+}\setminus[\check{y},\hat{y}]$,
using \eqref{us-3} and Corollary \ref{ncC3}  we divide it into two part to get its estimate as follows,
\begin{align}\label{JJ121}
&(\mathcal{P}(1-h)\partial_{y}^{-1}[\partial_{x}^{r+1}\mathbb{U}](\rho\partial_{y}u-\partial_{y}^{2}u),\varphi_{r}^{2}\mathbf{U}_{r})\\
=&~\int_{\mathbb{T}}\int_{\check{y}}^{\hat{y}}\varphi_{r}^{2}\mathcal{P}(1-h)\partial_{y}^{-1}[\partial_{x}^{r+1}\mathbb{U}](\rho\partial_{y}u-\partial_{y}^{2}u)\mathbf{U}_{r}dydx\nonumber\\
 &+\int_{\mathbb{T}}\int_{\mathbb{R}_{+}\setminus[\check{y},\hat{y}]}\varphi_{r}^{2}\mathcal{P}(1-h)\partial_{y}^{-1}[\partial_{x}^{r+1}\mathbb{U}](\rho\partial_{y}u-\partial_{y}^{2}u)\mathbf{U}_{r}dydx\nonumber\\
\leq&~C_{*}\left(\|\partial_{y}u\|_{L^{\infty}(\mathbb{T}\times[\check{y},\hat{y}])}
                 +\|\partial_{y}^{2}u\|_{L^{\infty}(\mathbb{T}\times[\check{y},\hat{y}])}\right)
          \|\vartheta_{c}\partial_{x}^{r+1}\mathbb{U}\|_{L_{\hat{\Psi}_{c}}^{2}(\Omega)}
          \|\mathcal{P}\varphi_{r}\|_{L_{x}^{\infty}(\mathbb{T})L_{y}^{2}([\check{y},\hat{y}])}
          \|\mathcal{P}\mathbf{U}_{r}\|_{L_{\varphi_{r}}^{2}(\Omega)}\nonumber\\
    &+C_{*}\left(\|\psi\partial_{y}u\|_{L_{x}^{\infty}L_{y}^{2}(\Omega)}+\|\psi\partial_{y}^{2}u\|_{L_{x}^{\infty}L_{y}^{2}(\Omega)}\right)
          \|\vartheta_{c}\partial_{x}^{r+1}\mathbb{U}\|_{L_{\hat{\Psi}_{r}}^{2}(\Omega)}\|\mathbf{U}_{r}\|_{L_{\varphi_{r}}^{2}(\Omega)}\nonumber\\
\leq&~C_{*}(\check{c}+\sqrt{\bar{\eta}}t)
          \left(\int_{\check{y}_{*}}^{\hat{y}_{*}}\frac{1}{((y-y_{*})^2+\bar{\varsigma}+t)^{1+\varepsilon_{0}}}dy\right)^{\frac{1}{2}}
          \|\vartheta_{c}\partial_{x}^{r+1}\mathbb{U}\|_{L_{\hat{\Psi}_{c}}^{2}(\Omega)}
          \|\mathcal{P}\mathbf{U}_{r}\|_{L_{\varphi_{r}}^{2}(\Omega)}\nonumber\\
    &+C_{*}(\check{c}+\sqrt{\bar{\eta}}t)\|\vartheta_{c}\partial_{x}^{r+1}\mathbb{U}\|_{L_{\hat{\Psi}_{r}}^{2}(\Omega)}
          \|\mathbf{U}_{r}\|_{L_{\varphi_{r}}^{2}(\Omega)}\nonumber\\
\leq&~C_{*}(\check{c}\iota^{-\frac{1}{4}-\frac{\varepsilon_{0}}{2}}+\sqrt{\bar{\eta}}t^{\frac{3}{4}-\frac{\varepsilon_{0}}{2}})
           \|\mathcal{P}\omega_{o}^{-1}\|_{L^{\infty}(\Omega)}
           \left(\|\vartheta_{c}\partial_{x}\mathbf{U}_{r}\|_{L_{\varphi_{r}}^{2}(\Omega)}
                 +\|\mathcal{P}\mathbf{U}_{r}\|_{L_{\varphi_{r}}^{2}(\Omega)}\right)
          \|\mathcal{P}\mathbf{U}_{r}\|_{L_{\varphi_{r}}^{2}(\Omega)}\nonumber\\
    &+C_{*}(\check{c}+\sqrt{\bar{\eta}}t)\|\vartheta_{c}\partial_{x}^{r+1}\mathbb{U}\|_{L_{\hat{\Psi}_{r}}^{2}(\Omega)}
          \|\mathbf{U}_{r}\|_{L_{\varphi_{r}}^{2}(\Omega)}\nonumber\\
\leq&~C_{*}(\check{c}\iota^{-\frac{3}{4}}+\sqrt{\bar{\eta}}t^{\frac{1}{4}})
      \left(\|\vartheta_{c}\partial_{x}\mathbf{U}_{r}\|_{L_{\varphi_{r}}^{2}(\Omega)}
            +\|\mathcal{P}\mathbf{U}_{r}\|_{L_{\varphi_{r}}^{2}(\Omega)}\right)
      \|\mathcal{P}\mathbf{U}_{r}\|_{L_{\varphi_{r}}^{2}(\Omega)}\nonumber\\
    &+C_{*}(\check{c}\iota^{-\frac{1}{2}}+\sqrt{\bar{\eta}t})
      \left(\|\vartheta_{c}\partial_{x}\mathbf{U}_{r}\|_{L_{\varphi_{r}}^{2}(\Omega)}
            +\|\mathcal{P}\mathbf{U}_{r}\|_{L_{\varphi_{r}}^{2}(\Omega)}\right)
      \|\mathbf{U}_{r}\|_{L_{\varphi_{r}}^{2}(\Omega)}\nonumber\\
\leq&~\frac{1}{16}\|\sqrt{\eta}\partial_{x}\mathbf{U}_{r}\|_{L_{\varphi_{r}}^{2}(\Omega)}^{2}
     +C_{*}(\sqrt{t}+\sqrt{\bar{\eta}}t^{\frac{1}{4}})\|\mathcal{P}\mathbf{U}_{r}\|_{L_{\varphi_{r}}^{2}(\Omega)}^{2}
     +C_{*}\iota^{-\frac{5}{2}}\|\mathbf{U}_{r}\|_{L_{\varphi_{r}}^{2}(\Omega)}^{2}.\nonumber
\end{align}
Here $\hat{\Psi}_{c}=\y^{-\frac{1}{2}}\Psi_{c}$, with $\Psi_{c}$ given below Corollary \ref{ncC1}. Analogously, it follows from Corollary \ref{ncC3} that
\begin{align}\label{JJ122}
&-(\mathcal{P}\partial_{y}^{-1}[\partial_{x}^{r+1}\tilde{\mathbb{U}}](\rho\partial_{y}u-\partial_{y}^{2}u),\varphi_{r}^{2}\mathbf{U}_{r})\\
\leq& C_{*}(\check{c}\iota^{-\frac{3}{4}}+\sqrt{\bar{\eta}}t^{\frac{1}{4}})
      \left(\|\vartheta_{c}\partial_{x}\tilde{\mathbf{U}}_{r}\|_{L_{\tilde{\varphi}_{r}}^{2}(\Omega)}
            +\|\tilde{\mathcal{P}}\tilde{\mathbf{U}}_{r}\|_{L_{\tilde{\varphi}_{r}}^{2}(\Omega)}\right)
      \|\mathcal{P}\mathbf{U}_{r}\|_{L_{\varphi_{r}}^{2}(\Omega)}\nonumber\\
    &+C_{*}(\check{c}\iota^{-\frac{1}{2}}+\sqrt{\bar{\eta}t})
      \left(\|\vartheta_{c}\partial_{x}\tilde{\mathbf{U}}_{r}\|_{L_{\tilde{\varphi}_{r}}^{2}(\Omega)}
            +\|\tilde{\mathcal{P}}\tilde{\mathbf{U}}_{r}\|_{L_{\tilde{\varphi}_{r}}^{2}(\Omega)}\right)
      \|\mathbf{U}_{r}\|_{L_{\varphi_{r}}^{2}(\Omega)}\nonumber\\
\leq& \tau\|\sqrt{\tilde{\eta}}\partial_{x}\tilde{\mathbf{U}}_{r}\|_{L_{\tilde{\varphi}_{r}}^{2}(\Omega)}
     +C_{*}C_{\tau}\sqrt{t}\|\mathcal{P}\mathbf{U}_{r}\|_{L_{\varphi_{r}}^{2}(\Omega)}^{2}
     +\tau\|\tilde{\mathcal{P}}\tilde{\mathbf{U}}_{r}\|_{L_{\tilde{\varphi}_{r}}^{2}(\Omega)}^{2}
     +C_{*}C_{\tau}\iota^{-\frac{5}{2}}\|\mathbf{U}_{r}\|_{L_{\varphi_{r}}^{2}(\Omega)}^{2}.\nonumber
\end{align}
Combining \eqref{JJ121} and \eqref{JJ122} with \eqref{JJ12}, we have
\begin{align}\label{JJ12'}
\mathfrak{J}_{12}
\leq& \frac{1}{16}\|\sqrt{\eta}\partial_{x}\mathbf{U}_{r}\|_{L_{\varphi_{r}}^{2}(\Omega)}^{2}
     +C_{*}C_{\tau}(\sqrt{t}+\sqrt{\bar{\eta}}t^{\frac{1}{4}})\|\mathcal{P}\mathbf{U}_{r}\|_{L_{\varphi_{r}}^{2}(\Omega)}^{2}\\
    &+C_{*}C_{\tau}\iota^{-\frac{5}{2}}\|\mathbf{U}_{r}\|_{L_{\varphi_{r}}^{2}(\Omega)}^{2}
     +\tau\|\sqrt{\tilde{\eta}}\partial_{x}\tilde{\mathbf{U}}_{r}\|_{L_{\tilde{\varphi}_{r}}^{2}(\Omega)}^{2}
     +\tau\|\tilde{\mathcal{P}}\tilde{\mathbf{U}}_{r}\|_{L_{\tilde{\varphi}_{r}}^{2}(\Omega)}^{2}.\nonumber
\end{align}

Noticing that ${\bf supp}~h(y)\in(\hat{y},+\infty)$, we have
\begin{align}\label{JJ13}
\mathfrak{J}_{13}
=&(\mathcal{P}(h'\partial_{y}^{-1}[\partial_{x}^{r+1}\mathbb{U}]\partial_{y}u+h\partial_{x}^{r+1}\mathbb{U}\partial_{y}u),\varphi_{r}^{2}\mathbf{U}_{r})
 -(\mathcal{P}(h'\partial_{y}^{-1}[\partial_{x}^{r+1}\tilde{\mathbb{U}}]\partial_{y}u+h\partial_{x}^{r+1}\tilde{\mathbb{U}}\partial_{y}u),\varphi_{r}^{2}\mathbf{U}_{r})\\
\leq& C_{*}\left(\|\partial_{y}u\|_{L^{\infty}(\Omega)}+\|\partial_{y}u\|_{L_{\varphi}^{2}(\Omega)}\right)
           \|\vartheta_{c}\partial_{x}^{r+1}\mathbb{U}\|_{L_{\hat{\Psi}_{r}}^{2}(\Omega)}\|\mathbf{U}_{r}\|_{L_{\varphi_{r}}^{2}(\Omega)}\nonumber\\
    &+C_{*}\left(\|\partial_{y}u\|_{L^{\infty}(\Omega)}+\|\partial_{y}u\|_{L_{\varphi}^{2}(\Omega)}\right)
           \|\vartheta_{c}\partial_{x}^{r+1}\tilde{\mathbb{U}}\|_{L_{\tilde{\hat{\Psi}}_{r}}^{2}(\Omega)}\|\mathbf{U}_{r}\|_{L_{\varphi_{r}}^{2}(\Omega)}\nonumber\\
\leq& C_{*}(\check{c}\iota^{-\frac{1}{2}}+\sqrt{\bar{\eta}t})
           \left(\|\vartheta_{c}\partial_{x}\mathbf{U}_{r}\|_{L_{\varphi_{r}}^{2}(\Omega)}+\|\mathcal{P}\mathbf{U}_{r}\|_{L_{\varphi_{r}}^{2}(\Omega)}\right)
           \|\mathbf{U}_{r}\|_{L_{\varphi_{r}}^{2}(\Omega)}\nonumber\\
    &+C_{*}(\check{c}\iota^{-\frac{1}{2}}+\sqrt{\bar{\eta}t})
           \left(\|\vartheta_{c}\partial_{x}\tilde{\mathbf{U}}_{r}\|_{L_{\tilde{\varphi}_{r}}^{2}(\Omega)}+\|\tilde{\mathcal{P}}\tilde{\mathbf{U}}_{r}\|_{L_{\tilde{\varphi}_{r}}^{2}(\Omega)}\right)
           \|\mathbf{U}_{r}\|_{L_{\varphi_{r}}^{2}(\Omega)}\nonumber\\
\leq& \frac{1}{16}\|\sqrt{\eta}\partial_{x}\mathbf{U}_{r}\|_{L_{\varphi_{r}}^{2}(\Omega)}^{2}
     +C_{*}C_{\tau}\iota^{-1}\|\mathbf{U}_{r}\|_{L_{\varphi_{r}}^{2}(\Omega)}^{2}
     +\tau\|\sqrt{\tilde{\eta}}\partial_{x}\tilde{\mathbf{U}}_{r}\|_{L_{\tilde{\varphi}_{r}}^{2}(\Omega)}^{2}
     +\|\tilde{\mathbf{U}}_{r}\|_{L_{\tilde{\varphi}_{r}}^{2}(\Omega)}^{2}.\nonumber
\end{align}

Form \eqref{JJ11'}, \eqref{JJ12'} and \eqref{JJ13} one knows that
\begin{align}\label{JJ1}
\mathfrak{J}_{1}
\leq& \frac{3}{16}\|\sqrt{\eta}\partial_{x}\mathbf{U}_{r}\|_{L_{\varphi_{r}}^{2}(\Omega)}^{2}
     +C_{*}C_{\tau}(\sqrt{t}+\sqrt{\bar{\eta}}t^{\frac{1}{4}})\|\mathcal{P}\mathbf{U}_{r}\|_{L_{\varphi_{r}}^{2}(\Omega)}^{2}
     +C_{*}C_{\tau}\iota^{-\frac{5}{2}}\|\mathbf{U}_{r}\|_{L_{\varphi_{r}}^{2}(\Omega)}^{2}\\
    &+4\tau\|\sqrt{\tilde{\eta}}\partial_{x}\tilde{\mathbf{U}}_{r}\|_{L_{\tilde{\varphi}_{r}}^{2}(\Omega)}^{2}
     +\tau\|\tilde{\mathcal{P}}\tilde{\mathbf{U}}_{r}\|_{L_{\tilde{\varphi}_{r}}^{2}(\Omega)}^{2}
     +C\|\tilde{\mathbf{U}}_{r}\|_{L_{\tilde{\varphi}_{r}}^{2}(\Omega)}^{2}.\nonumber
\end{align}

Next let's turn to the estimate of $\mathfrak{J}_{2}$. We shall estimate $\mathfrak{J}_{2}$ in different ways for $r=s$ and $r=s+1$ respectively.
When $r=s$, by a direct calculation,  with \eqref{us-3} and Corollary \ref{ncC2} one can obtain
\begin{align}\label{JJ2(1)}
\mathfrak{J}_{2}
=& (\mathcal{P}\partial_{x}^{s}u(\partial_{x}\partial_{y}u-\rho\partial_{x}u),\varphi^{2}\mathbf{U}_{s})
  +s(\mathcal{P}\partial_{x}u\partial_{x}^{s}\partial_{y}u,\varphi^{2}\mathbf{U}_{s})
  -s(\mathcal{P}\rho\partial_{x}u\partial_{x}^{s}u,\varphi^{2}\mathbf{U}_{s})\\
 &-s(\mathcal{P}(1-h)\partial_{y}^{-1}[\partial_{x}^{s}u](\partial_{x}\partial_{y}^{2}u-\rho\partial_{x}\partial_{y}u),\varphi^{2}\mathbf{U}_{s})\nonumber\\
 &-(\mathcal{P}(1-h)\partial_{y}^{-1}[\partial_{x}u](\partial_{x}^{s}\partial_{y}^{2}u-\rho\partial_{x}^{s}\partial_{y}u),\varphi^{2}\mathbf{U}_{s})\nonumber\\
 &+s(\mathcal{P}h'\partial_{y}^{-1}[\partial_{x}^{s}u]\partial_{x}\partial_{y}u+\mathcal{P}h\partial_{x}^{s}u\partial_{x}\partial_{y}u,\varphi^{2}\mathbf{U}_{s})\nonumber\\
 &+(\mathcal{P}h'\partial_{y}^{-1}[\partial_{x}u]\partial_{x}^{s}\partial_{y}u+\mathcal{P}h\partial_{x}u\partial_{x}^{s}\partial_{y}u),\varphi^{2}\mathbf{U}_{s})\nonumber\\
\leq&~C_{*}\left(\|\partial_{x}\partial_{y}u\|_{L^{\infty}(\Omega)}+\|\partial_{x}u\|_{L^{\infty}(\Omega)}\right)
           \|\partial_{x}^{s}u\|_{L_{\hat{\Psi}}^{2}(\Omega)}\|\mathbf{U}_{s}\|_{L_{\varphi}^{2}(\Omega)}\nonumber\\
    &+C_{*}\|\partial_{x}u\|_{L^{\infty}(\Omega)}
           \left(\|\partial_{x}^{s}\partial_{y}u\|_{L_{\hat{\Psi}}^{2}(\Omega)}+\|\partial_{x}^{s}\partial_{y}^{2}u\|_{L_{\hat{\Psi}}^{2}(\Omega)}\right)
           \|\mathbf{U}_{s}\|_{L_{\varphi}^{2}(\Omega)}\nonumber\\
    &+C_{*}(\|\partial_{x}\partial_{y}^{2}u\|_{L^{\infty}(\Omega)}+\|\partial_{x}\partial_{y}u\|_{L^{\infty}(\Omega)})
           \|\mathcal{P}\varphi\|_{L_{x}^{\infty}(\mathbb{T})L_{y}^{2}(0,\hat{y})}\|\partial_{x}^{s}u\|_{L_{\hat{\psi}}^{2}(\Omega)}\|\mathcal{P}\mathbf{U}_{s}\|_{L_{\varphi}^{2}(\Omega)}\nonumber\\
\leq&~C_{*}(\check{c}+\sqrt{\bar{\eta}}t)
           \|\tilde{\mathcal{P}}\mathcal{U}_{s}\|_{L_{\tilde{\varphi}}^{2}(\Omega)}\|\mathbf{U}_{s}\|_{L_{\varphi}^{2}(\Omega)}
     +C_{*}(\check{c}\iota^{-1}+\sqrt{\bar{\eta}})
           \|\tilde{\mathcal{P}}\mathcal{U}_{s}\|_{L_{\tilde{\varphi}}^{2}(\Omega)}\|\mathbf{U}_{s}\|_{L_{\varphi}^{2}(\Omega)}\nonumber\\
    &+C_{*}(\check{c}+\sqrt{\bar{\eta}}t)
           \left(\|\mathcal{U}_{s}\|_{L_{\tilde{\varphi}}^{2}(\Omega)}
                 +\|\tilde{\mathcal{P}}\mathcal{U}_{s}\|_{L_{\tilde{\varphi}}^{2}(\Omega)}
                 +\|\partial_{y}\mathcal{U}_{s}\|_{L_{\tilde{\varphi}}^{2}(\Omega)}\right)
           \|\mathbf{U}_{s}\|_{L_{\varphi}^{2}(\Omega)}\nonumber\\
\leq&~\|\mathcal{U}_{s}\|_{L_{\tilde{\varphi}}^{2}(\Omega)}^{2}
     +\tau\|\tilde{\mathcal{P}}\mathcal{U}_{s}\|_{L_{\tilde{\varphi}}^{2}(\Omega)}
     +\tau\|\partial_{y}\mathcal{U}_{s}\|_{L_{\tilde{\varphi}}^{2}(\Omega)}^{2}
     +C_{*}C_{\tau}\iota^{-2}\|\mathbf{U}_{s}\|_{L_{\varphi}^{2}(\Omega)}^{2}.\nonumber
\end{align}
When $r=s+1$, we divide $\mathfrak{J}_{2}$ into four parts as follows,
\begin{align}\label{JJ2(2)}
\mathfrak{J}_{2}
&= (\mathcal{P}\partial_{x}^{s+1}u(\partial_{x}\partial_{y}u-\rho\partial_{x}u),\hat{\varphi}^{2}\mathbf{U}_{s+1})
  +(s+1)(\mathcal{P}\partial_{x}u(\partial_{x}^{s+1}\partial_{y}u-\rho\partial_{x}^{s+1}u),\hat{\varphi}^{2}\mathbf{U}_{s+1})\\
&\quad-(s+1)(\mathcal{P}(1-h)\partial_{y}^{-1}[\partial_{x}^{s+1}u](\partial_{x}\partial_{y}^{2}u-\rho\partial_{x}\partial_{y}u),\hat{\varphi}^{2}\mathbf{U}_{s+1})\nonumber\\
&\quad-(\mathcal{P}(1-h)\partial_{y}^{-1}[\partial_{x}u](\partial_{x}^{s+1}\partial_{y}^{2}u-\rho\partial_{x}^{s+1}\partial_{y}u),\hat{\varphi}^{2}\mathbf{U}_{s+1})\nonumber\\
&\quad+(s+1)(\mathcal{P}\partial_{x}\partial_{y}u(h'\partial_{y}^{-1}[\partial_{x}^{s+1}u]+h\partial_{x}^{s+1}u),\hat{\varphi}^{2}\mathbf{U}_{s+1})\nonumber\\
&\quad+(\mathcal{P}\partial_{x}^{s+1}\partial_{y}u(h'\partial_{y}^{-1}[\partial_{x}u]+h\partial_{x}u),\hat{\varphi}^{2}\mathbf{U}_{s+1})\nonumber\\
&=: \mathfrak{J}_{21}+\mathfrak{J}_{22}+\mathfrak{J}_{23}+\mathfrak{J}_{24}+\mathfrak{J}_{25}+\mathfrak{J}_{26}.\nonumber
\end{align}
By using \eqref{us-3} and Corollary \ref{ncC3} we have
\begin{align}\label{JJ21}
\mathfrak{J}_{21}
=&~(\mathcal{P}\partial_{x}^{s+1}u(\partial_{x}\partial_{y}u-\rho\partial_{x}u),\hat{\varphi}^{2}\mathbf{U}_{s+1})\\
\leq&~C_{*}\left(\|\vartheta_{c}^{-1}\partial_{x}\partial_{y}u\|_{L^{\infty}(\Omega)}+\|\vartheta_{c}^{-1}\partial_{x}u\|_{L^{\infty}(\Omega)}\right)
      \|\vartheta_{c}\tilde{\mathcal{P}}^{2}\partial_{x}^{s+1}u\|_{L_{\tilde{\varphi}}^{2}(\Omega)}\|\mathbf{U}_{s+1}\|_{L_{\hat{\varphi}}^{2}(\Omega)}\nonumber\\
\leq&~C_{*}(\check{c}+\sqrt{\bar{\eta}}t)
      \|\vartheta_{c}\partial_{x}^{s+1}u\|_{L_{\tilde{\hat{\Psi}}}^{2}(\Omega)}\|\mathbf{U}_{s+1}\|_{L_{\hat{\varphi}}^{2}(\Omega)}\nonumber\\
\leq&~C_{*}(\check{c}+\sqrt{\bar{\eta}}t)\|\tilde{\mathcal{P}}\|_{L^{\infty}(\Omega)}
      \left(\|\tilde{\mathcal{P}}\mathcal{U}_{s}\|_{L_{\tilde{\varphi}}^{2}(\Omega)}+\|\vartheta_{c}\partial_{x}\mathcal{U}_{s}\|_{L_{\tilde{\varphi}}^{2}(\Omega)}\right)
      \|\mathbf{U}_{s+1}\|_{L_{\hat{\varphi}}^{2}(\Omega)}\nonumber\\
\leq&~C_{*}(\check{c}\iota^{-\frac{1}{2}}+\sqrt{\bar{\eta}t})
      \left(\|\tilde{\mathcal{P}}\mathcal{U}_{s}\|_{L_{\tilde{\varphi}}^{2}(\Omega)}+\|\vartheta_{c}\partial_{x}\mathcal{U}_{s}\|_{L_{\tilde{\varphi}}^{2}(\Omega)}\right)
      \|\mathbf{U}_{s+1}\|_{L_{\hat{\varphi}}^{2}(\Omega)}\nonumber\\
\leq&~\tau\|\sqrt{\eta}\partial_{x}\mathcal{U}_{s}\|_{L_{\tilde{\varphi}}^{2}(\Omega)}^{2}
     +C_{*}C_{\tau}\iota^{-2}\|\mathbf{U}_{s+1}\|_{L_{\hat{\varphi}}^{2}(\Omega)}^{2}
     +\|\mathcal{U}_{s}\|_{L_{\tilde{\varphi}}^{2}(\Omega)}^{2}.\nonumber
\end{align}
Since $\partial_{y}\partial_{x}^{s+1}u-\tilde{\rho}\partial_{x}^{s+1}u=\partial_{x}(\tilde{\mathcal{P}}^{-1}\mathcal{U}_{s})+\partial_{x}\tilde{\rho}\partial_{x}^{s}u$,
similar to \eqref{JJ21} we have
\begin{align}\label{JJ22}
\mathfrak{J}_{22}
=&~(s+1)(\mathcal{P}\partial_{x}u(\partial_{x}^{s+1}\partial_{y}u-\rho\partial_{x}^{s+1}u),\hat{\varphi}^{2}\mathbf{U}_{s+1})\\
=&~(s+1)(\mathcal{P}\partial_{x}u\partial_{x}(\tilde{\mathcal{P}}^{-1}\mathcal{U}_{s}),\hat{\varphi}^{2}\mathbf{U}_{s+1})
  +(s+1)(\mathcal{P}\partial_{x}u\partial_{x}\tilde{\rho}\partial_{x}^{s}u,\hat{\varphi}^{2}\mathbf{U}_{s+1})\nonumber\\
 &+(s+1)(\mathcal{P}\partial_{x}u(\tilde{\rho}-\rho)\partial_{x}^{s+1}u,\hat{\varphi}^{2}\mathbf{U}_{s+1})\nonumber\\
\leq&~C\|\vartheta_{c}^{-1}\partial_{x}u\|_{L^{\infty}(\Omega)}
       \|\vartheta_{c}\partial_{x}\mathcal{U}_{s}\|_{L_{\varphi}^{2}(\Omega)}
       \|\mathbf{U}_{s+1}\|_{L_{\hat{\varphi}}^{2}(\Omega)}\nonumber\\
    &+C_{*}\|\partial_{x}u\|_{L^{\infty}(\Omega)}
           \|\tilde{\mathcal{P}}\mathcal{U}_{s}\|_{L_{\tilde{\varphi}}^{2}(\Omega)}
           \|\mathbf{U}_{s+1}\|_{L_{\hat{\varphi}}^{2}(\Omega)}\nonumber\\
    &+C_{*}\|\partial_{x}u\|_{L^{\infty}(\Omega)}
           \|\partial_{x}^{s}u\|_{L_{\hat{\Psi}}^{2}(\Omega)}
           \|\mathbf{U}_{s+1}\|_{L_{\hat{\varphi}}^{2}(\Omega)}\nonumber\\
    &+C_{*}\|\vartheta_{c}^{-1}\partial_{x}u\|_{L^{\infty}(\Omega)}
           \|\vartheta_{c}\tilde{\mathcal{P}}^{2}\partial_{x}^{s+1}u\|_{L_{\tilde{\hat{\varphi}}}^{2}(\Omega)}
           \|\mathbf{U}_{s+1}\|_{L_{\hat{\varphi}}^{2}(\Omega)}\nonumber\\
\leq&~C_{*}(\check{c}+\sqrt{\bar{\eta}}t)
      \left(\|\tilde{\mathcal{P}}\mathcal{U}_{s}\|_{L_{\tilde{\varphi}}^{2}(\Omega)}
            +\|\vartheta_{c}\partial_{x}\mathcal{U}_{s}\|_{L_{\tilde{\varphi}}^{2}(\Omega)}\right)
      \|\mathbf{U}_{s+1}\|_{L_{\hat{\varphi}}^{2}(\Omega)}\nonumber\\
    &+C_{*}(\check{c}+\sqrt{\bar{\eta}}t)
      \|\vartheta_{c}\partial_{x}^{s+1}u\|_{L_{\hat{\Psi}}^{2}(\Omega)}\|\mathbf{U}_{s+1}\|_{L_{\tilde{\hat{\varphi}}}^{2}(\Omega)}\nonumber\\
\leq&~\tau\|\sqrt{\eta}\partial_{x}\mathcal{U}_{s}\|_{L_{\tilde{\varphi}}^{2}(\Omega)}^{2}
     +C_{*}C_{\tau}\iota^{-2}\|\mathbf{U}_{s+1}\|_{L_{\hat{\varphi}}^{2}(\Omega)}^{2}
     +\|\mathcal{U}_{s}\|_{L_{\tilde{\varphi}}^{2}(\Omega)}^{2}.\nonumber
\end{align}
Similar to that in \eqref{JJ121}, one has
\begin{align}\label{JJ23}
\mathfrak{J}_{23}
=&-(s+1)(\mathcal{P}(1-h)\partial_{y}^{-1}[\partial_{x}^{s+1}u](\partial_{x}\partial_{y}^{2}u-\rho\partial_{x}\partial_{y}u),\hat{\varphi}^{2}\mathbf{U}_{s+1})\\
=&~(s+1)\int_{\mathbb{T}}\int_{\check{y}}^{\hat{y}}\hat{\varphi}^{2}\mathcal{P}(1-h)\partial_{y}^{-1}[\partial_{x}^{s+1}u](\rho\partial_{x}\partial_{y}u-\partial_{x}\partial_{y}^{2}u)\mathbf{U}_{s+1}dydx\nonumber\\
 &+(s+1)\int_{\mathbb{T}}\int_{\mathbb{R}_{+}\setminus[\check{y},\hat{y}]}\hat{\varphi}^{2}\mathcal{P}(1-h)\partial_{y}^{-1}[\partial_{x}^{s+1}u](\rho\partial_{x}\partial_{y}u-\partial_{x}\partial_{y}^{2}u)\mathbf{U}_{s+1}dydx\nonumber\\
\leq&~C_{*}\left(\|\partial_{x}\partial_{y}u\|_{L^{\infty}(\mathbb{T}\times[\check{y},\hat{y}])}
                 +\|\partial_{x}\partial_{y}^{2}u\|_{L^{\infty}(\mathbb{T}\times[\check{y},\hat{y}])}\right)
          \|\vartheta_{c}\partial_{x}^{s+1}u\|_{L_{\hat{\psi}}^{2}(\Omega)}
          \|\mathcal{P}\hat{\varphi}\|_{L_{x}^{\infty}(\mathbb{T})L_{y}^{2}([\check{y},\hat{y}])}
          \|\mathcal{P}\mathbf{U}_{s+1}\|_{L_{\hat{\varphi}}^{2}(\Omega)}\nonumber\\
    &+C_{*}\left(\|\varphi\partial_{x}\partial_{y}u\|_{L_{x}^{\infty}L_{y}^{2}(\Omega)}+\|\varphi\partial_{x}\partial_{y}^{2}u\|_{L_{x}^{\infty}L_{y}^{2}(\Omega)}\right)
          \|\vartheta_{c}\partial_{x}^{s+1}u\|_{L_{\tilde{\hat{\Psi}}}^{2}(\Omega)}\|\mathbf{U}_{s+1}\|_{L_{\hat{\varphi}}^{2}(\Omega)}\nonumber\\
\leq&~C_{*}(\check{c}+\sqrt{\bar{\eta}}t)
          \left(\int_{\check{y}_{*}}^{\hat{y}_{*}}\frac{1}{((y-y_{*})^2+\bar{\varsigma}+t)^{1+\varepsilon_{0}}}dy\right)^{\frac{1}{2}}
          \|\vartheta_{c}\partial_{x}^{s+1}u\|_{L_{\hat{\psi}}^{2}(\Omega)}
          \|\mathcal{P}\mathbf{U}_{s+1}\|_{L_{\hat{\varphi}}^{2}(\Omega)}\nonumber\\
    &+C_{*}(\check{c}+\sqrt{\bar{\eta}}t)\|\vartheta_{c}\partial_{x}^{s+1}u\|_{L_{\tilde{\hat{\Psi}}}^{2}(\Omega)}
          \|\mathbf{U}_{s+1}\|_{L_{\hat{\varphi}}^{2}(\Omega)}\nonumber\\
\leq&~C_{*}(\check{c}\iota^{-\frac{3}{4}}+\sqrt{\bar{\eta}}t^{\frac{1}{4}})
      \left(\|\vartheta_{c}\partial_{x}\mathcal{U}_{s}\|_{L_{\tilde{\varphi}}^{2}(\Omega)}
            +\|\tilde{\mathcal{P}}\mathcal{U}_{s}\|_{L_{\tilde{\varphi}}^{2}(\Omega)}\right)
      \|\mathcal{P}\mathbf{U}_{s+1}\|_{L_{\hat{\varphi}}^{2}(\Omega)}\nonumber\\
    &+C_{*}(\check{c}\iota^{-\frac{1}{2}}+\sqrt{\bar{\eta}t})
      \left(\|\vartheta_{c}\partial_{x}\mathcal{U}_{s}\|_{L_{\tilde{\varphi}}^{2}(\Omega)}
            +\|\tilde{\mathcal{P}}\mathcal{U}_{s}\|_{L_{\tilde{\varphi}}^{2}(\Omega)}\right)
      \|\mathbf{U}_{s+1}\|_{L_{\hat{\varphi}}^{2}(\Omega)}\nonumber\\
\leq&~\tau\|\sqrt{\eta}\partial_{x}\mathcal{U}_{s}\|_{L_{\tilde{\varphi}}^{2}(\Omega)}^{2}
     +\tau\|\tilde{\mathcal{P}}\mathcal{U}_{s}\|_{L_{\tilde{\varphi}}^{2}(\Omega)}^{2}
     +C_{*}C_{\tau}\sqrt{t}\|\mathcal{P}\mathbf{U}_{s+1}\|_{L_{\hat{\varphi}}^{2}(\Omega)}^{2}
     +C_{*}C_{\tau}\iota^{-\frac{5}{2}}\|\mathbf{U}_{s+1}\|_{L_{\hat{\varphi}}^{2}(\Omega)}^{2}.\nonumber
\end{align}

By using $u=\mathbb{U}-\tilde{\mathbb{U}}$, we decompose $\mathfrak{J}_{24}$ into two parts as follows,
\begin{align}\label{JJ24}
\mathfrak{J}_{24}
=&-(\mathcal{P}(1-h)\partial_{y}^{-1}[\partial_{x}u](\partial_{x}^{s+1}\partial_{y}^{2}u-\rho\partial_{x}^{s+1}\partial_{y}u),\hat{\varphi}^{2}\mathbf{U}_{s+1})\\
=&~(\mathcal{P}(1-h)\partial_{y}^{-1}[\partial_{x}u](\rho\partial_{x}^{s+1}\partial_{y}\mathbb{U}-\partial_{x}^{s+1}\partial_{y}^{2}\mathbb{U}),\hat{\varphi}^{2}\mathbf{U}_{s+1})\nonumber\\
 &-(\mathcal{P}(1-h)\partial_{y}^{-1}[\partial_{x}u](\rho\partial_{x}^{s+1}\partial_{y}\tilde{\mathbb{U}}-\partial_{x}^{s+1}\partial_{y}^{2}\tilde{\mathbb{U}}),\hat{\varphi}^{2}\mathbf{U}_{s+1}).\nonumber
\end{align}
Since $\partial_{x}^{s+1}\partial_{y}^{2}\mathbb{U}-\rho\partial_{x}^{s+1}\partial_{y}\mathbb{U}=\partial_{y}(\mathcal{P}^{-1}\mathbf{U}_{s+1})+\partial_{y}\rho\partial_{x}^{s+1}\mathbb{U}$,
one has
\begin{align}\label{JJ241}
&(\mathcal{P}(1-h)\partial_{y}^{-1}[\partial_{x}u](\rho\partial_{x}^{s+1}\partial_{y}\mathbb{U}-\partial_{x}^{s+1}\partial_{y}^{2}\mathbb{U}),\hat{\varphi}^{2}\mathbf{U}_{s+1})\\
=&-(\mathcal{P}(1-h)\partial_{y}^{-1}[\partial_{x}u]\partial_{y}\mathcal{P}^{-1}\mathbf{U}_{s+1},\hat{\varphi}^{2}\mathbf{U}_{s+1})
  -((1-h)\partial_{y}^{-1}[\partial_{x}u]\partial_{y}\mathbf{U}_{s+1},\hat{\varphi}^{2}\mathbf{U}_{s+1})\nonumber\\
 &-(\mathcal{P}(1-h)\partial_{y}^{-1}[\partial_{x}u]\partial_{y}\rho\partial_{x}^{s+1}\mathbb{U},\hat{\varphi}^{2}\mathbf{U}_{s+1})\nonumber\\
\leq&~C_{*}\|\psi\partial_{x}u\|_{L_{x}^{\infty}L_{y}^{2}(\Omega)}
           \|\mathcal{P}\mathbf{U}_{s+1}\|_{L_{\hat{\varphi}}^{2}(\Omega)}
           \|\mathbf{U}_{s+1}\|_{L_{\hat{\varphi}}^{2}(\Omega)}
     +C\|\psi\partial_{x}u\|_{L_{x}^{\infty}L_{y}^{2}(\Omega)}
       \|\partial_{y}\mathbf{U}_{s+1}\|_{L_{\hat{\varphi}}^{2}(\Omega)}
       \|\mathbf{U}_{s+1}\|_{L_{\hat{\varphi}}^{2}(\Omega)}\nonumber\\
    &+C_{*}\|\psi\partial_{x}u\|_{L_{x}^{\infty}L_{y}^{2}(\Omega)}
           \|\mathcal{P}^{3}\partial_{x}^{s+1}\mathbb{U}\|_{L_{\hat{\varphi}}^{2}(\Omega)}
           \|\mathbf{U}_{s+1}\|_{L_{\hat{\varphi}}^{2}(\Omega)}\nonumber\\
\leq&~C_{*}\iota^{-\frac{1}{2}}\|\mathbf{U}_{s+1}\|_{L_{\hat{\varphi}}^{2}(\Omega)}^{2}
     +\frac{1}{16}\|\partial_{y}\mathbf{U}_{s+1}\|_{L_{\hat{\varphi}}^{2}(\Omega)}^{2}
     +C_{*}(\check{c}\iota^{-\frac{1}{2}}+\sqrt{\bar{\eta}t})
      \|\partial_{x}^{s+1}\mathbb{U}\|_{L_{\hat{\Psi}_{s+1}}^{2}(\Omega)}\|\mathbf{U}_{s+1}\|_{L_{\hat{\varphi}}^{2}(\Omega)}\nonumber\\
\leq&~C_{*}\iota^{-\frac{1}{2}}\|\mathbf{U}_{s+1}\|_{L_{\hat{\varphi}}^{2}(\Omega)}^{2}
     +\frac{1}{16}\|\partial_{y}\mathbf{U}_{s+1}\|_{L_{\hat{\varphi}}^{2}(\Omega)}^{2}
     +C_{*}(\check{c}\iota^{-\frac{1}{2}}+\sqrt{\bar{\eta}t})
      \|\mathcal{P}\mathbf{U}_{s+1}\|_{L_{\hat{\varphi}}^{2}(\Omega)}\|\mathbf{U}_{s+1}\|_{L_{\hat{\varphi}}^{2}(\Omega)}\nonumber\\
\leq&~C_{*}\iota^{-1}\|\mathbf{U}_{s+1}\|_{L_{\hat{\varphi}}^{2}(\Omega)}^{2}
     +\frac{1}{16}\|\partial_{y}\mathbf{U}_{s+1}\|_{L_{\hat{\varphi}}^{2}(\Omega)}^{2}.\nonumber
\end{align}
Parallel to \eqref{JJ241}, one has
\begin{align}\label{JJ242}
&(\mathcal{P}(1-h)\partial_{y}^{-1}[\partial_{x}u](\rho\partial_{x}^{s+1}\partial_{y}\tilde{\mathbb{U}}-\partial_{x}^{s+1}\partial_{y}^{2}\tilde{\mathbb{U}}),\hat{\varphi}^{2}\mathbf{U}_{s+1})\\
=&-(\mathcal{P}(1-h)\partial_{y}^{-1}[\partial_{x}u]\partial_{y}\mathcal{P}^{-1}\tilde{\mathbf{U}}_{s+1},\hat{\varphi}^{2}\mathbf{U}_{s+1})
  -((1-h)\partial_{y}^{-1}[\partial_{x}u]\partial_{y}\tilde{\mathbf{U}}_{s+1},\hat{\varphi}^{2}\mathbf{U}_{s+1})\nonumber\\
 &-(\mathcal{P}(1-h)\partial_{y}^{-1}[\partial_{x}u]\partial_{y}\rho\partial_{x}^{s+1}\tilde{\mathbb{U}},\hat{\varphi}^{2}\mathbf{U}_{s+1})\nonumber\\
\leq&~C_{*}\|\psi\partial_{x}u\|_{L_{x}^{\infty}L_{y}^{2}(\Omega)}
           \|\mathcal{P}\mathbf{U}_{s+1}\|_{L_{\hat{\varphi}}^{2}(\Omega)}
           \|\tilde{\mathbf{U}}_{s+1}\|_{L_{\tilde{\hat{\varphi}}}^{2}(\Omega)}
     +C\|\psi\partial_{x}u\|_{L_{x}^{\infty}L_{y}^{2}(\Omega)}
       \|\partial_{y}\tilde{\mathbf{U}}_{s+1}\|_{L_{{\tilde{\hat\varphi}}}^{2}(\Omega)}
       \|\mathbf{U}_{s+1}\|_{L_{\hat{\varphi}}^{2}(\Omega)}\nonumber\\
    &+C_{*}\|\psi\partial_{x}u\|_{L_{x}^{\infty}L_{y}^{2}(\Omega)}
           \|\tilde{\mathcal{P}}^{2}\partial_{x}^{s+1}\tilde{\mathbb{U}}\|_{L_{\tilde{\hat{\varphi}}}^{2}(\Omega)}
           \|\mathcal{P}\mathbf{U}_{s+1}\|_{L_{\hat{\varphi}}^{2}(\Omega)}\nonumber\\
\leq&~C_{*}(\check{c}\iota^{-\frac{1}{2}}+\sqrt{\bar{\eta}t})
            \|\mathbf{U}_{s+1}\|_{L_{\hat{\varphi}}^{2}(\Omega)}\|\tilde{\mathbf{U}}_{s+1}\|_{L_{\tilde{\hat{\varphi}}}^{2}(\Omega)}
     +\frac{1}{16}\|\partial_{y}\tilde{\mathbf{U}}_{s+1}\|_{L_{\tilde{\hat{\varphi}}}^{2}(\Omega)}^{2}
     +C_{*}\|\mathbf{U}_{s+1}\|_{L_{\hat{\varphi}}^{2}(\Omega)}^{2}\nonumber\\
    &+C_{*}(\check{c}+\sqrt{\bar{\eta}}t)
      \|\partial_{x}^{s+1}\tilde{\mathbb{U}}\|_{L_{\tilde{\hat{\Psi}}_{s+1}}^{2}(\Omega)}\|\mathcal{P}\mathbf{U}_{s+1}\|_{L_{\hat{\varphi}}^{2}(\Omega)}\nonumber\\
\leq&~C_{*}\iota^{-1}\|\mathbf{U}_{s+1}\|_{L_{\hat{\varphi}}^{2}(\Omega)}^{2}
     +\frac{1}{2}\|\tilde{\mathbf{U}}_{s+1}\|_{L_{\tilde{\hat\varphi}}^{2}(\Omega)}^{2}
     +\frac{1}{16}\|\partial_{y}\tilde{\mathbf{U}}_{s+1}\|_{L_{\tilde{\hat{\varphi}}}^{2}(\Omega)}^{2}\nonumber\\
    &+C_{*}(\check{c}+\sqrt{\bar{\eta}}t)
      \|\tilde{\mathcal{P}}\tilde{\mathbf{U}}_{s+1}\|_{L_{\tilde{\hat{\varphi}}}^{2}(\Omega)}\|\mathcal{P}\mathbf{U}_{s+1}\|_{L_{\hat{\varphi}}^{2}(\Omega)}\nonumber\\
\leq&~C_{*}\iota^{-2}\|\mathbf{U}_{s+1}\|_{L_{\hat{\varphi}}^{2}(\Omega)}^{2}
     +\|\tilde{\mathbf{U}}_{s+1}\|_{L_{\tilde{\hat{\varphi}}}^{2}(\Omega)}^{2}
     +\frac{1}{16}\|\partial_{y}\mathbf{U}_{s+1}\|_{L_{\hat{\varphi}}^{2}(\Omega)}^{2}.\nonumber
\end{align}
Substituting \eqref{JJ241} and \eqref{JJ242} into \eqref{JJ24}, one has
\begin{align}\label{JJ24'}
\mathfrak{J}_{24}
\leq C_{*}\iota^{-2}\|\mathbf{U}_{s+1}\|_{L_{\hat{\varphi}}^{2}(\Omega)}^{2}
    +\|\tilde{\mathbf{U}}_{s+1}\|_{L_{\tilde{\hat{\varphi}}}^{2}(\Omega)}^{2}
    +\frac{1}{8}\|\partial_{y}\mathbf{U}_{s+1}\|_{L_{\hat{\varphi}}^{2}(\Omega)}^{2}.
\end{align}

Similar to \eqref{JJ21}, one has
\begin{align}\label{JJ25}
\mathfrak{J}_{25}
=&~(s+1)(\mathcal{P}\partial_{x}\partial_{y}u(h'\partial_{y}^{-1}[\partial_{x}^{s+1}u]+h\partial_{x}^{s+1}u),\hat{\varphi}^{2}\mathbf{U}_{s+1})\\
\leq&~C_{*}\left(\|\partial_{x}\partial_{y}u\|_{L^{\infty}(\Omega)}+\|\partial_{x}\partial_{y}u\|_{L_{x}^{\infty}L_{y}^{2}(\Omega)}\right)
           \|\vartheta_{c}\partial_{x}^{s+1}u\|_{L_{\tilde{\hat{\Psi}}}^{2}(\Omega)}\|\mathbf{U}_{s+1}\|_{L_{\hat{\varphi}}^{2}(\Omega)}\nonumber\\
\leq&~C_{*}(\check{c}+\sqrt{\bar{\eta}}t)
      \|\vartheta_{c}\partial_{x}^{s+1}u\|_{L_{\tilde{\hat{\Psi}}}^{2}(\Omega)}\|\mathbf{U}_{s+1}\|_{L_{\hat{\varphi}}^{2}(\Omega)}\nonumber\\
\leq&~C_{*}(\check{c}+\sqrt{\bar{\eta}}t)\|\tilde{\mathcal{P}}\|_{L^{\infty}(\Omega)}
      \left(\|\tilde{\mathcal{P}}\mathcal{U}_{s}\|_{L_{\tilde{\varphi}}^{2}(\Omega)}+\|\vartheta_{c}\partial_{x}\mathcal{U}_{s}\|_{L_{\tilde{\varphi}}^{2}(\Omega)}\right)
      \|\mathbf{U}_{s+1}\|_{L_{\hat{\varphi}}^{2}(\Omega)}\nonumber\\
\leq&~C_{*}(\check{c}\iota^{-\frac{1}{2}}+\sqrt{\bar{\eta}t})
      \left(\|\tilde{\mathcal{P}}\mathcal{U}_{s}\|_{L_{\tilde{\varphi}}^{2}(\Omega)}+\|\vartheta_{c}\partial_{x}\mathcal{U}_{s}\|_{L_{\tilde{\varphi}}^{2}(\Omega)}\right)
      \|\mathbf{U}_{s+1}\|_{L_{\hat{\varphi}}^{2}(\Omega)}\nonumber\\
\leq&~\tau\|\sqrt{\eta}\partial_{x}\mathcal{U}_{s}\|_{L_{\tilde{\varphi}}^{2}(\Omega)}^{2}
     +C_{*}C_{\tau}\iota^{-2}\|\mathbf{U}_{s+1}\|_{L_{\hat{\varphi}}^{2}(\Omega)}^{2}
     +\|\mathcal{U}_{s}\|_{L_{\tilde{\varphi}}^{2}(\Omega)}^{2}.\nonumber
\end{align}

Since $\partial_{y}\partial_{x}^{s+1}u-\tilde{\rho}\partial_{x}^{s+1}u=\partial_{x}(\tilde{\mathcal{P}}^{-1}\mathcal{U}_{s})+\partial_{x}\tilde{\rho}\partial_{x}^{s}u$, one has
\begin{align*}
\|\vartheta_{c}\partial_{y}\partial_{x}^{s+1}u\|_{L_{\varphi}^{2}(\Omega)}
\leq& C_{*}\|\vartheta_{c}\partial_{x}^{s+1}u\|_{L_{\tilde{\hat{\Psi}}}^{2}(\Omega)}
    +C_{*}\|\vartheta_{c}\partial_{x}^{s}u\|_{L_{\tilde{\hat{\Psi}}}^{2}(\Omega)}
    +\|\vartheta_{c}\partial_{x}(\mathcal{P}^{-1}\mathcal{U}_{s})\|_{L_{\varphi}^{2}(\Omega)}\\
\leq& C_{*}\|\tilde{\mathcal{P}}\|_{L^{\infty}(\Omega)}
           \left(\|\tilde{\mathcal{P}}\mathcal{U}_{s}\|_{L_{\tilde{\varphi}}^{2}(\Omega)}+\|\vartheta_{c}\partial_{x}\mathcal{U}_{s}\|_{L_{\tilde{\varphi}}^{2}(\Omega)}\right)
     +C_{*}\|\tilde{\mathcal{P}}\mathcal{U}_{s}\|_{L_{\tilde{\varphi}}^{2}(\Omega)}\\
    &+C_{*}\|\mathcal{U}_{s}\|_{L_{\tilde{\varphi}}^{2}(\Omega)}
     +C_{*}\|\vartheta_{c}\partial_{x}\mathcal{U}_{s}\|_{L_{\tilde{\varphi}}^{2}(\Omega)},
\end{align*}
moreover, $\vartheta_{c}^{-1}\leq C\psi$ owing to $t\leq \min\{\frac{1}{12\Lambda},\frac{1}{4\theta}\}$, thus
\begin{align}\label{JJ26}
\mathfrak{J}_{26}
\leq&~C_{*}\left(\|\vartheta_{c}^{-1}\partial_{x}u\|_{L^{\infty}(\Omega)}
                 +\|\vartheta_{c}^{-1}\partial_{x}u\|_{L_{x}^{\infty}L_{y}^{2}(\Omega)}\right)
          \|\vartheta_{c}\partial_{y}\partial_{x}^{s+1}u\|_{L_{\varphi}^{2}(\Omega)}\|\mathbf{U}_{s+1}\|_{L_{\hat{\varphi}}^{2}(\Omega)}\\
\leq&~C_{*}(\check{c}\iota^{-1}+\sqrt{\bar{\eta}})
           \left(\|\mathcal{U}_{s}\|_{L_{\tilde{\varphi}}^{2}(\Omega)}+\|\vartheta_{c}\partial_{x}\mathcal{U}_{s}\|_{L_{\tilde{\varphi}}^{2}(\Omega)}\right)
           \|\mathbf{U}_{s+1}\|_{L_{\hat{\varphi}}^{2}(\Omega)}\nonumber\\
\leq&~\tau\|\sqrt{\eta}\partial_{x}\mathcal{U}_{s}\|_{L_{\varphi}^{2}(\Omega)}^{2}
     +C_{*}C_{\tau}\iota^{-2}\|\mathbf{U}_{s+1}\|_{L_{\hat{\varphi}}^{2}(\Omega)}^{2}
     +\|\mathcal{U}_{s}\|_{L_{\tilde{\varphi}}^{2}(\Omega)}^{2}.\nonumber
\end{align}

Combining \eqref{JJ21},\eqref{JJ22},\eqref{JJ23},\eqref{JJ24'},\eqref{JJ25} and \eqref{JJ26} with \eqref{JJ2(2)}, one can conclude that when $r=s+1$,
\begin{align*}
\mathfrak{J}_{2}
\leq&~\frac{1}{8}\|\partial_{y}\mathbf{U}_{s+1}\|_{L_{\hat{\varphi}}^{2}(\Omega)}^{2}
     +C_{*}C_{\tau}\iota^{-\frac{5}{2}}\|\mathbf{U}_{s+1}\|_{L_{\hat{\varphi}}^{2}(\Omega)}^{2}
     +C_{*}C_{\tau}\sqrt{t}\|\mathcal{P}\mathbf{U}_{s+1}\|_{L_{\hat{\varphi}}^{2}(\Omega)}^{2}\\
    &+\|\tilde{\mathbf{U}}_{s+1}\|_{L_{\tilde{\hat{\varphi}}}^{2}(\Omega)}^{2}
     +C\|\mathcal{U}_{s}\|_{L_{\tilde{\varphi}}^{2}(\Omega)}^{2}
     +5\tau\|\sqrt{\eta}\partial_{x}\mathcal{U}_{s}\|_{L_{\tilde{\varphi}}^{2}(\Omega)}^{2}
     +\tau\|\tilde{\mathcal{P}}\mathcal{U}_{s}\|_{L_{\tilde{\varphi}}^{2}(\Omega)}^{2}.
\end{align*}

Now let's focus on the estimate of $\mathfrak{J}_{3}$. When $r=s+1$, the estimate of $\mathfrak{J}_{3}$ is similar to that in \eqref{JJ2(1)},
\begin{align*}
\mathfrak{J}_{3}
=&~(s+1)(\mathcal{P}\partial_{x}^{s}u(\partial_{x}^{2}\partial_{y}u-\rho\partial_{x}^{2}u),\hat{\varphi}^{2}\mathbf{U}_{s+1})
  +\frac{s(s+1)}{2}(\partial_{x}^{2}u\mathcal{U}_{s},\hat{\varphi}^{2}\mathbf{U}_{s+1})\\
 &-\frac{s(s+1)}{2}(\mathcal{P}(1-h)\partial_{y}^{-1}[\partial_{x}^{s}u](\partial_{x}^{2}\partial_{y}^{2}u-\rho\partial_{x}^{2}\partial_{y}u),\hat{\varphi}^{2}\mathbf{U}_{s+1})\nonumber\\
 &+\frac{s(s+1)}{2}(\mathcal{P}\partial_{x}^{2}\partial_{y}u(h'\partial_{y}^{-1}[\partial_{x}^{s}u]+h\partial_{x}^{s}u),\hat{\varphi}^{2}\mathbf{U}_{s+1})\nonumber\\
 &-(s+1)(\mathcal{P}(1-h)\partial_{y}^{-1}[\partial_{x}^{2}u][\partial_{y}(\mathcal{P}^{-1}\mathcal{U}_{s})+\partial_{y}\rho\partial_{x}^{s}u],\hat{\varphi}^{2}\mathbf{U}_{s+1})\nonumber\\
 &+(s+1)(\mathcal{P}\partial_{x}^{s}\partial_{y}u(h'\partial_{y}^{-1}[\partial_{x}^{2}u]+h\partial_{x}^{2}u),\hat{\varphi}^{2}\mathbf{U}_{s+1})\nonumber\\
\leq&~C_{*}(\|\partial_{x}^{2}\partial_{y}u\|_{L^{\infty}(\Omega)}+\|\partial_{x}^{2}u\|_{L^{\infty}(\Omega)})
           \|\partial_{x}^{s}u\|_{L_{\hat{\Psi}}^{2}(\Omega)}\|\mathbf{U}_{s+1}\|_{L_{\hat{\varphi}}^{2}(\Omega)}
     +C\|\partial_{x}^{2}u\|_{L^{\infty}(\Omega)}\|\mathcal{U}_{s}\|_{L_{\tilde{\varphi}}^{2}(\Omega)}\|\mathbf{U}_{s+1}\|_{L_{\hat{\varphi}}^{2}(\Omega)}\nonumber\\
    &+C_{*}(\|\partial_{x}^{2}\partial_{y}^{2}u\|_{L^{\infty}(\Omega)}+\|\partial_{x}^{2}\partial_{y}u\|_{L^{\infty}(\Omega)})
           \|\mathcal{P}\varphi\|_{L_{x}^{\infty}(\mathbb{T})L_{y}^{2}(0,\hat{y})}
           \|\partial_{x}^{s}u\|_{L_{\hat{\psi}}^{2}(\Omega)}\|\mathcal{P}\mathbf{U}_{s+1}\|_{L_{\hat{\varphi}}^{2}(\Omega)}\nonumber\\
    &+C_{*}\|\psi\partial_{x}^{2}u\|_{L_{x}^{\infty}L_{y}^{2}(\Omega)}\|\partial_{y}\mathcal{U}_{s}\|_{L_{\tilde{\varphi}}^{2}(\Omega)}
           \|\mathbf{U}_{s+1}\|_{L_{\hat{\varphi}}^{2}(\Omega)}
     +C_{*}\|\psi\partial_{x}^{2}u\|_{L_{x}^{\infty}L_{y}^{2}(\Omega)}\|\tilde{\mathcal{P}}\mathcal{U}_{s}\|_{L_{\tilde{\varphi}}^{2}(\Omega)}
           \|\mathbf{U}_{s+1}\|_{L_{\hat{\varphi}}^{2}(\Omega)}\nonumber\\
    &+C_{*}\|\psi\partial_{x}^{2}u\|_{L_{x}^{\infty}L_{y}^{2}(\Omega)}\|\partial_{x}^{s}u\|_{L_{\hat{\Psi}}^{2}(\Omega)}
           \|\mathcal{P}\mathbf{U}_{s+1}\|_{L_{\hat{\varphi}}^{2}(\Omega)}\nonumber\\
    &+C_{*}\left(\|\psi\partial_{x}^{2}u\|_{L_{x}^{\infty}L_{y}^{2}(\Omega)}+\|\partial_{x}^{2}u\|_{L^{\infty}(\Omega)}\right)
           \|\partial_{y}\partial_{x}^{s}u\|_{L_{\varphi}^{2}(\Omega)}\|\mathbf{U}_{s+1}\|_{L_{\check{\varphi}}^{2}(\Omega)}\\
\leq&~C_{*}(\check{c}+\sqrt{\bar{\eta}}t)\|\tilde{\mathcal{P}}\mathcal{U}_{s}\|_{L_{\tilde{\varphi}}^{2}(\Omega)}
                                             \|\mathbf{U}_{s+1}\|_{L_{\hat{\varphi}}^{2}(\Omega)}
     +C_{*}\|\mathcal{U}_{s}\|_{L_{\tilde{\varphi}}^{2}(\Omega)}\|\mathbf{U}_{s+1}\|_{L_{\hat{\varphi}}^{2}(\Omega)}\nonumber\\
    &+C_{*}(\check{c}\iota^{-1}+\sqrt{\bar{\eta}})\|\tilde{\mathcal{P}}\mathcal{U}_{s}\|_{L_{\tilde{\varphi}}^{2}(\Omega)}
                                                  \|\mathbf{U}_{s+1}\|_{L_{\hat{\varphi}}^{2}(\Omega)}
     +C_{*}\|\partial_{y}\mathcal{U}_{s}\|_{L_{\tilde{\varphi}}^{2}(\Omega)}\|\mathbf{U}_{s+1}\|_{L_{\hat{\varphi}}^{2}(\Omega)}\nonumber\\
\leq&~\tau\|\tilde{\mathcal{P}}\mathcal{U}_{s}\|_{L_{\tilde{\varphi}}^{2}(\Omega)}
     +\tau\|\partial_{y}\mathcal{U}_{s}\|_{L_{\tilde{\varphi}}^{2}(\Omega)}^{2}
     +C\|\mathcal{U}_{s}\|_{L_{\tilde{\varphi}}^{2}(\Omega)}^{2}
     +C_{*}C_{\tau}\iota^{-2}\|\mathbf{U}_{s+1}\|_{L_{\hat{\varphi}}^{2}(\Omega)}^{2},\nonumber
\end{align*}
while when $r=s$, since $\bar{\varsigma}=\iota\bar{\eta}^{2},\varepsilon_{0}<\frac{1}{4}$ and $\bar{\varsigma}=\iota\epsilon_{0}^{4}$ when $\check{c}\neq 0$, by noticing \eqref{us-3} one gets,
\begin{align*}
\mathfrak{J}_{3}
=&~s(\mathcal{P}\partial_{x}^{s-1}u(\partial_{x}^{2}\partial_{y}u-\rho\partial_{x}^{2}u),\varphi^{2}\mathbf{U}_{s})
  +\frac{s(s-1)}{2}(\mathcal{P}\partial_{x}^{2}u(\partial_{x}^{s-1}\partial_{y}u-\rho\partial_{x}^{s-1}u),\varphi^{2}\mathbf{U}_{s})\\
 &-\frac{s(s-1)}{2}(\mathcal{P}(1-h)\partial_{y}^{-1}[\partial_{x}^{s-1}u](\partial_{x}^{2}\partial_{y}^{2}u-\rho\partial_{x}^{2}\partial_{y}u),\varphi^{2}\mathbf{U}_{s})\\
 &-s(\mathcal{P}(1-h)\partial_{y}^{-1}[\partial_{x}^{2}u](\partial_{x}^{s-1}\partial_{y}^{2}u-\rho\partial_{x}^{s-1}\partial_{y}u),\varphi^{2}\mathbf{U}_{s})\nonumber\\
 &+\frac{s(s-1)}{2}(\mathcal{P}\partial_{x}^{2}\partial_{y}u(h'\partial_{y}^{-1}[\partial_{x}^{s-1}u]+h\partial_{x}^{s-1}u),\varphi^{2}\mathbf{U}_{s})\\
 &+s(\mathcal{P}\partial_{x}^{s-1}\partial_{y}u(h'\partial_{y}^{-1}[\partial_{x}^{2}u]+h\partial_{x}^{2}u),\varphi^{2}\mathbf{U}_{s})\nonumber\\
\leq&~C_{*}(\|\mathcal{P}\varphi\|_{L^{\infty}(\Omega)}\|\partial_{x}^{2}\partial_{y}u\|_{L^{\infty}(\Omega)}
            +\|\mathcal{P}^{2}\varphi\|_{L^{\infty}(\Omega)}\|\partial_{x}^{2}u\|_{L^{\infty}(\Omega)})
           \|u\|_{\hat{H}_{\psi}^{s}(\Omega)}\|\mathbf{U}_{s}\|_{L_{\varphi}^{2}(\Omega)}\nonumber\\
    &+C_{*}\|\mathcal{P}\varphi\|_{L^{\infty}(\Omega)}(1+\|\mathcal{P}\|_{L^{\infty}(\Omega)})
           \left(\|\partial_{x}^{2}u\|_{L^{\infty}(\Omega)}+\|\partial_{x}^{2}u\|_{L_{x}^{\infty}L_{y}^{2}(\Omega)}\right)
           \|u\|_{\hat{H}_{\psi}^{s}(\Omega)}\|\mathbf{U}_{s}\|_{L_{\varphi}^{2}(\Omega)}\nonumber\\
    &+C_{*}(\|\mathcal{P}\varphi\|_{L^{\infty}(\Omega)}\|\partial_{x}^{2}\partial_{y}^{2}u\|_{L_{x}^{\infty}L_{y}^{2}(\Omega)}
            +\|\mathcal{P}^{2}\varphi\|_{L^{\infty}(\Omega)}\|\partial_{x}^{2}\partial_{y}u\|_{L_{x}^{\infty}L_{y}^{2}(\Omega)})
           \|u\|_{\hat{H}_{\psi}^{s}(\Omega)}\|\mathbf{U}_{s}\|_{L_{\varphi}^{2}(\Omega)}\nonumber\\
    &+C_{*}\|\mathcal{P}\varphi\|_{L^{\infty}(\Omega)}\left(\|\partial_{x}^{2}\partial_{y}u\|_{L_{x}^{\infty}L_{y}^{2}(\Omega)}
                                                     +\|\partial_{x}^{2}\partial_{y}u\|_{L^{\infty}(\Omega)}\right)
           \|u\|_{\hat{H}_{\psi}^{s}(\Omega)}\|\mathbf{U}_{s}\|_{L_{\varphi}^{2}(\Omega)}\nonumber\\
\leq&~C_{*}\left(\|\mathcal{P}\varphi\|_{L^{\infty}(\Omega)}+\|\mathcal{P}^{2}\varphi\|_{L^{\infty}(\Omega)}\right)
           \|u\|_{H^{s-3}(\Omega)}\|u\|_{\hat{H}_{\psi}^{s}(\Omega)}\|\mathbf{U}_{s}\|_{L_{\varphi}^{2}(\Omega)}\\
\leq&~\frac{C_{*}(\check{c}+\sqrt{\bar{\eta}}t)}{(\bar{\varsigma}+t)^{1+\frac{\varepsilon_{0}}{2}}}
      \|u\|_{\hat{H}_{\psi}^{s}(\Omega)}\|\mathbf{U}_{s}\|_{L_{\varphi}^{2}(\Omega)}\\
\leq&~\|u\|_{\hat{H}_{\psi}^{s}(\Omega)}^{2}
     +C_{*}\iota^{-\frac{5}{2}}\|\mathbf{U}_{s}\|_{L_{\varphi}^{2}(\Omega)}^{2}.
\end{align*}

Finally, using the method for estimating the term $\mathfrak{J}_{3}$ in the case $r=s$, one can obtain,
\begin{align*}
\mathfrak{J}_{4}
=&~\sum_{i=3}^{r-2}\binom{r}{i}(\mathcal{P}\partial_{x}^{i}u\partial_{x}^{r-i+1}\partial_{y}u,\varphi_{r}^{2}\mathbf{U}_{r})
  -\sum_{i=3}^{r-2}\binom{r}{i}(\mathcal{P}\rho\partial_{x}^{i}u\partial_{x}^{r-i+1}u,\varphi_{r}^{2}\mathbf{U}_{r})\\
 &-\sum_{i=2}^{r-3}\binom{r}{i}(\mathcal{P}(1-h)\partial_{y}^{-1}[\partial_{x}^{i+1}u]\partial_{x}^{r-i}\partial_{y}^{2}u,\varphi_{r}^{2}\mathbf{U}_{r})
  +\sum_{i=2}^{r-3}\binom{r}{i}(\mathcal{P}(1-h)\rho\partial_{y}^{-1}[\partial_{x}^{i+1}u]\partial_{x}^{r-i}\partial_{y}u,\varphi_{r}^{2}\mathbf{U}_{r})\\
 &+\sum_{i=2}^{r-3}\binom{r}{i}(\mathcal{P}h\partial_{x}^{i+1}u\partial_{x}^{r-i}\partial_{y}u,\varphi_{r}^{2}\mathbf{U}_{r})
  +\sum_{i=2}^{r-3}\binom{r}{i}(\mathcal{P}h'\partial_{y}^{-1}[\partial_{x}^{i+1}u]\partial_{x}^{r-i}\partial_{y}u,\varphi_{r}^{2}\mathbf{U}_{r})\\
\leq&~C_{*}\left(\|\mathcal{P}\varphi_{r}\|_{L^{\infty}(\Omega)}+\|\mathcal{P}^{2}\varphi_{r}\|_{L^{\infty}(\Omega)}\right)
           \|u\|_{H^{s-3}(\Omega)}\|u\|_{\hat{H}_{\psi}^{s}(\Omega)}\|\mathbf{U}_{r}\|_{L_{\varphi_{r}}^{2}(\Omega)}\\
\leq&~\frac{C_{*}(\check{c}+\sqrt{\bar{\eta}}t)}{(\bar{\varsigma}+t)^{1+\frac{\varepsilon_{0}}{2}}}
      \|u\|_{\hat{H}_{\psi}^{s}(\Omega)}\|\mathbf{U}_{r}\|_{L_{\varphi_{r}}^{2}(\Omega)}\\
\leq&~\|u\|_{\hat{H}_{\psi}^{s}(\Omega)}^{2}
     +C_{*}\iota^{-\frac{5}{2}}\|\mathbf{U}_{r}\|_{L_{\varphi_{r}}^{2}(\Omega)}^{2}.
\end{align*}

Collecting all the estimate of $\mathfrak{J}_{1}, \mathfrak{J}_{2}, \mathfrak{J}_{3}$ and $\mathfrak{J}_{4}$ above, together with \eqref{J10} and \eqref{J101'} one can draw a conclusion that when $r=s$,
\begin{align*}
J_{10}
\leq& \frac{1}{4}\|\sqrt{\eta}\partial_{x}\mathbf{U}_{s}\|_{L_{\varphi}^{2}(\Omega)}^{2}
     +C_{*}C_{\tau}\iota^{-\frac{5}{2}}\|\mathbf{U}_{s}\|_{L_{\varphi}^{2}(\Omega)}^{2}
     +C_{*}C_{\tau}(\sqrt{t}+\sqrt{\bar{\eta}}t^{\frac{1}{4}})\|\mathcal{P}\mathbf{U}_{s}\|_{L_{\varphi}^{2}(\Omega)}^{2}\\
    &+5\tau\|\sqrt{\tilde{\eta}}\partial_{x}\tilde{\mathbf{U}}_{s}\|_{L_{\tilde{\varphi}}^{2}(\Omega)}^{2}
     +2\tau\|\tilde{\mathcal{P}}\tilde{\mathbf{U}}_{s}\|_{L_{\tilde{\varphi}}^{2}(\Omega)}^{2}
     +\tau\|\tilde{\mathcal{P}}\mathcal{U}_{s}\|_{L_{\tilde{\varphi}}^{2}(\Omega)}
     +\tau\|\partial_{y}\mathcal{U}_{s}\|_{L_{\tilde{\varphi}}^{2}(\Omega)}^{2}\\
    &+C\|\mathcal{U}_{s}\|_{L_{\tilde{\varphi}}^{2}(\Omega)}^{2}
     +C_{*}\|\tilde{\mathbf{U}}_{s}\|_{L_{\tilde{\varphi}}^{2}(\Omega)}^{2}
     +C\|u\|_{\hat{H}_{\psi}^{s}(\Omega)}^{2},
\end{align*}
and when $r=s+1$,
\begin{align*}
J_{10}
\leq&~\frac{1}{4}\|\sqrt{\eta}\partial_{x}\mathbf{U}_{s+1}\|_{L_{\hat{\varphi}}^{2}(\Omega)}^{2}
     +\frac{1}{8}\|\partial_{y}\mathbf{U}_{s+1}\|_{L_{\hat{\varphi}}^{2}(\Omega)}^{2}
     +C_{*}C_{\tau}(\sqrt{t}+\sqrt{\bar{\eta}}t^{\frac{1}{4}})\|\mathcal{P}\mathbf{U}_{s+1}\|_{L_{\hat{\varphi}}^{2}(\Omega)}^{2}\\
    &+C_{*}C_{\tau}\iota^{-\frac{5}{2}}\|\mathbf{U}_{s+1}\|_{L_{\hat{\varphi}}^{2}(\Omega)}^{2}
     +5\tau\|\sqrt{\tilde{\eta}}\partial_{x}\tilde{\mathbf{U}}_{s+1}\|_{L_{\tilde{\hat{\varphi}}}^{2}(\Omega)}^{2}
     +\tau\|\tilde{\mathcal{P}}\tilde{\mathbf{U}}_{s+1}\|_{L_{\tilde{\hat{\varphi}}}^{2}(\Omega)}^{2}\\
    &+C\|\tilde{\mathbf{U}}_{s+1}\|_{L_{\tilde{\hat{\varphi}}}^{2}(\Omega)}^{2}
     +\tau\|\partial_{y}\mathcal{U}_{s}\|_{L_{\tilde{\varphi}}^{2}(\Omega)}^{2}
     +5\tau\|\sqrt{\eta}\partial_{x}\mathcal{U}_{s}\|_{L_{\tilde{\varphi}}^{2}(\Omega)}^{2}
     +C\|\mathcal{U}_{s}\|_{L_{\tilde{\varphi}}^{2}(\Omega)}^{2}\\
    &+2\tau\|\tilde{\mathcal{P}}\mathcal{U}_{s}\|_{L_{\tilde{\varphi}}^{2}(\Omega)}^{2}
     +C\|u\|_{\hat{H}_{\psi}^{s}(\Omega)}^{2}.
\end{align*}

\noindent{\bf \underline{Estimate of~$J_{11}$.}}
We decompose $J_{11}$ into four parts,
\begin{align}\label{J11}
J_{11}
=\sum_{i=1}^{4}(\mathcal{P}\bar{\mathcal{T}}_{8i},\varphi_{r}^{2}\mathbf{U}_{r}),
\end{align}
where
\begin{align*}
\bar{\mathcal{T}}_{81}
=& \rho\sum_{i=2}^{r}\binom{r}{i}\partial_{x}^{r-i+1}\mathbb{U}\partial_{x}^{i}u_{0}
  +\rho\sum_{i=0}^{r-1}\binom{r}{i}\partial_{x}^{r-i+1}u_{0}\partial_{x}^{i}\mathbb{U}\\
 &-\sum_{i=0}^{r-1}\binom{r}{i}\partial_{y}\partial_{x}^{r-i+1}u_{0}\partial_{x}^{i}\mathbb{U}
  -\sum_{i=2}^{r}\binom{r}{i}\partial_{y}\partial_{x}^{r-i+1}\mathbb{U}\partial_{x}^{i}u_{0}\\
 &+\sum_{i=0}^{r-2}\binom{r}{i}\partial_{y}^{-1}[\partial_{x}^{i+1}\mathbb{U}]\partial_{x}^{r-i}\partial_{y}^{2}u_{0}
  -\rho\sum_{i=0}^{r-2}\binom{r}{i}\partial_{y}^{-1}[\partial_{x}^{i+1}\mathbb{U}]\partial_{x}^{r-i}\partial_{y}u_{0},
\end{align*}
\begin{align*}
\bar{\mathcal{T}}_{82}
= \sum_{i=1}^{r-2}\binom{r}{i}\partial_{y}^{-1}[\partial_{x}^{i+1}u_{0}]\partial_{x}^{r-i}\partial_{y}^{2}\mathbb{U}
  -\rho\sum_{i=1}^{r-2}\binom{r}{i}\partial_{y}^{-1}[\partial_{x}^{i+1}u_{0}]\partial_{x}^{r-i}\partial_{y}\mathbb{U},
\end{align*}
\begin{align*}
\bar{\mathcal{T}}_{83}
=\rho\partial_{x}^{r}(\partial_{t}u_{0}+u_{0}\partial_{x}u_{0}+\partial_{x}P),
\end{align*}
and
\begin{align*}
\bar{\mathcal{T}}_{84}
=&-\sum_{i=0}^{r}\binom{r}{i}\partial_{y}\partial_{x}^{r-i+1}u_{0}\partial_{x}^{i}u_{0}
  +\sum_{i=0}^{r-2}\binom{r}{i}\partial_{y}^{-1}[\partial_{x}^{i+1}u_{0}]\partial_{x}^{r-i}\partial_{y}^{2}u_{0}\\
 &-\rho\sum_{i=0}^{r-2}\binom{r}{i}\partial_{y}^{-1}[\partial_{x}^{i+1}u_{0}]\partial_{x}^{r-i}\partial_{y}u_{0}
  +\partial_{y}^{3}\partial_{x}^{r}u_{0}
  -\rho\partial_{x}^{r}\partial_{y}^{2}u_{0}
  -\partial_{t}\partial_{y}\partial_{x}^{r}u_{0}.
\end{align*}
Recall the definition of $u_{0}$, by using Corollary \ref{ncC2} one has
\begin{align}\label{J111}
(\mathcal{P}\bar{\mathcal{T}}_{81},\varphi_{r}^{2}\mathbf{U}_{r})
\leq&~C_{*}\left(\|\mathbb{U}\|_{\hat{H}_{\psi}^{s}(\Omega)}+\|\partial_{x}^{s}\mathbb{U}\|_{L_{\varphi}^{2}(\Omega)}
              +\|\partial_{y}\partial_{x}^{s}\mathbb{U}\|_{L_{\varphi}^{2}(\Omega)}\right)
        \|\mathbf{U}_{r}\|_{L_{\varphi_{r}}^{2}(\Omega)}\\
\leq&~C_{*}\left(\|\mathbb{U}\|_{\hat{H}_{\psi}^{s}(\Omega)}+\|\mathcal{P}\mathbf{U}_{s}\|_{L_{\varphi}^{2}(\Omega)}
              +\|\mathbf{U}_{s}\|_{L_{\varphi}^{2}(\Omega)}\right)
        \|\mathbf{U}_{r}\|_{L_{\varphi_{r}}^{2}(\Omega)}\nonumber\\
\leq&~C_{*}\|\mathbb{U}\|_{\hat{H}_{\psi}^{s}(\Omega)}^{2}
     +\|\mathcal{P}\mathbf{U}_{s}\|_{L_{\varphi}^{2}(\Omega)}^{2}
     +C_{*}\|\mathbf{U}_{r}\|_{L_{\varphi_{r}}^{2}(\Omega)}^{2}
     +C_{*}\|\mathbf{U}_{s}\|_{L_{\varphi}^{2}(\Omega)}.\nonumber
\end{align}

The term $(\mathcal{P}\bar{\mathcal{T}}_{82},\varphi_{r}^{2}\mathbf{U}_{r})$ is the most difficult term to estimate owing to the nonlocal term. Thanks to the factor $\y^{-\frac{1}{2}}$ in the weight function $\varphi$, if $r=s$,
\begin{align}\label{J112-1}
(\mathcal{P}\bar{\mathcal{T}}_{82},\varphi_{r}^{2}\mathbf{U}_{r})
\leq&~C_{*}\left(\|\mathbb{U}\|_{\hat{H}_{\psi}^{s}(\Omega)}+\|\sqrt{\y}\partial_{x}^{s-1}\partial_{y}^{2}\mathbb{U}\|_{L_{\hat{\psi}}^{2}(\Omega)}\right)
        \|\sqrt{\y}\mathbf{U}_{r}\|_{L_{\varphi_{r}}^{2}(\Omega)}\\
\leq&~C_{*}\left(\|\mathbb{U}\|_{\hat{H}_{\psi}^{s}(\Omega)}+\|\mathbb{U}\|_{\hat{H}_{\psi,\y}^{s}(\Omega)}\right)
        \|\sqrt{\y}\mathbf{U}_{r}\|_{L_{\varphi_{r}}^{2}(\Omega)}\nonumber\\
\leq&~C_{*}\|\mathbb{U}\|_{\hat{H}_{\psi}^{s}(\Omega)}^{2}
     +\tau\|\mathbb{U}\|_{\hat{H}_{\psi,\y}^{s}(\Omega)}^{2}
     +C_{*}C_{\tau}\|\sqrt{\y}\mathbf{U}_{r}\|_{L_{\varphi_{r}}^{2}(\Omega)}^{2}.\nonumber
\end{align}
and if $r=s+1$,
\begin{align}\label{J112-2}
(\mathcal{P}\bar{\mathcal{T}}_{82},\varphi_{r}^{2}\mathbf{U}_{r})
\leq&~C_{*}\left(\|\mathbb{U}\|_{\hat{H}_{\psi}^{s}(\Omega)}+\|\partial_{x}^{s-1}\partial_{y}^{2}\mathbb{U}\|_{L_{\hat{\psi}}^{2}(\Omega)}
                 +\|\partial_{x}^{s}\partial_{y}^{2}\mathbb{U}\|_{L_{\varphi}^{2}(\Omega)}\right)
        \|\sqrt{\y}\mathbf{U}_{r}\|_{L_{\varphi_{r}}^{2}(\Omega)}\\
\leq&~C_{*}\left(\|\mathbb{U}\|_{\hat{H}_{\psi}^{s}(\Omega)}+\|\mathbf{U}_{s}\|_{L_{\varphi}^{2}(\Omega)}
                 +\|\mathcal{P}\mathbf{U}_{s}\|_{L_{\varphi}^{2}(\Omega)}+\|\partial_{y}\mathbf{U}_{s}\|_{L_{\varphi}^{2}(\Omega)}\right)
        \|\sqrt{\y}\mathbf{U}_{r}\|_{L_{\varphi_{r}}^{2}(\Omega)}\nonumber\\
\leq&~C_{*}\|\mathbb{U}\|_{\hat{H}_{\psi}^{s}(\Omega)}^{2}
     +C_{*}\|\mathbf{U}_{s}\|_{L_{\varphi}^{2}(\Omega)}^{2}
     +\|\mathcal{P}\mathbf{U}_{s}\|_{L_{\varphi}^{2}(\Omega)}^{2}\nonumber\\
    &+\frac{1}{16}\|\partial_{y}\mathbf{U}_{s}\|_{L_{\varphi}^{2}(\Omega)}^{2}
     +C_{*}\|\sqrt{\y}\mathbf{U}_{r}\|_{L_{\varphi_{r}}^{2}(\Omega)}^{2}.\nonumber
\end{align}
A combination of \eqref{J112-1} with \eqref{J112-2} gives
\begin{align}\label{J112}
(\mathcal{P}\bar{\mathcal{T}}_{82},\varphi_{r}^{2}\mathbf{U}_{r})
\leq&~C_{*}\|\mathbb{U}\|_{\hat{H}_{\psi}^{s}(\Omega)}^{2}
     +C_{*}\|\mathbf{U}_{s}\|_{L_{\varphi}^{2}(\Omega)}^{2}
     +\tau\|\mathbb{U}\|_{\hat{H}_{\psi,\y}^{s}(\Omega)}^{2}\\
    &+\|\mathcal{P}\mathbf{U}_{s}\|_{L_{\varphi}^{2}(\Omega)}^{2}
     +\frac{1}{16}\|\partial_{y}\mathbf{U}_{s}\|_{L_{\varphi}^{2}(\Omega)}^{2}
     +C_{*}C_{\tau}\|\sqrt{\y}\mathbf{U}_{r}\|_{L_{\varphi_{r}}^{2}(\Omega)}^{2}.\nonumber
\end{align}

From \eqref{B} one can deduce
$$\partial_{t}u_{0}+u_{0}\partial_{x}u_{0}+\partial_{x}P=\partial_{t}(u_{0}-U)+(u_{0}-U)\partial_{x}u_{0}+U\partial_{x}(u_{0}-U),$$
moreover, from the definition of $u_{0}$ one can get
$$\partial_{x}^{r}(\partial_{t}u_{0}+u_{0}\partial_{x}u_{0})=0,~\text{in}~\hat{\Omega}_{t_{3}},$$
thus one has
\begin{align}\label{J113}
(\mathcal{P}\bar{\mathcal{T}}_{83},\varphi_{r}^{2}\mathbf{U}_{r})
\leq&~C_{*}\|\mathcal{P}\partial_{x}^{r+1}P\|_{L_{\varphi_{r}}^{2}(\hat{\Omega})}\|\mathcal{P}\mathbf{U}_{r}\|_{L_{\varphi_{r}}^{2}(\Omega)}\\
    &+C_{*}\|\partial_{x}^{r}(\partial_{t}u_{0}+u_{0}\partial_{x}u_{0}+\partial_{x}P)\|_{L_{\psi}^{2}(\Omega)}\|\mathbf{U}_{r}\|_{L_{\varphi_{r}}^{2}(\Omega)}\nonumber\\
\leq&~\frac{C_{*}}{(\bar{\varsigma}+t)^{\frac{1}{2}+\varepsilon_{0}}}\|\partial_{x}^{r+1}P\|_{L^{2}(\mathbb{T})}^{2}
     +\|\mathcal{P}\mathbf{U}_{r}\|_{L_{\varphi_{r}}^{2}(\Omega)}^{2}
     +\|\mathbf{U}_{r}\|_{L_{\varphi_{r}}^{2}(\Omega)}^{2}
     +C_{*}.\nonumber
\end{align}

Finally, it is easy to see that
\begin{align}\label{J114}
(\mathcal{P}\bar{\mathcal{T}}_{84},\varphi_{r}^{2}\mathbf{U}_{r})
\leq \|\mathbf{U}_{r}\|_{L_{\varphi_{r}}^{2}(\Omega)}^{2}
    +C_{*}.
\end{align}

Combining \eqref{J111}, \eqref{J112}, \eqref{J113} and \eqref{J114} with \eqref{J11}, one has
\begin{align*}
J_{11}
\leq&~C_{*}\|\mathbb{U}\|_{\hat{H}_{\psi}^{s}(\Omega)}^{2}
     +C_{*}\|\mathbf{U}_{s}\|_{L_{\varphi}^{2}(\Omega)}^{2}
     +C_{*}\|\mathbf{U}_{r}\|_{L_{\varphi_{r}}^{2}(\Omega)}^{2}
     +2\|\mathcal{P}\mathbf{U}_{s}\|_{L_{\varphi}^{2}(\Omega)}^{2}
     +\|\mathcal{P}\mathbf{U}_{r}\|_{L_{\varphi_{r}}^{2}(\Omega)}^{2}\nonumber\\
    &+\frac{1}{16}\|\partial_{y}\mathbf{U}_{s}\|_{L_{\varphi}^{2}(\Omega)}^{2}
     +C_{*}C_{\tau}\|\sqrt{\y}\mathbf{U}_{r}\|_{L_{\varphi_{r}}^{2}(\Omega)}^{2}
     +\frac{C_{*}}{(\bar{\varsigma}+t)^{\frac{1}{2}+\varepsilon_{0}}}\|\partial_{x}^{r+1}P\|_{L^{2}(\mathbb{T})}^{2}
     +C_{*}.
\end{align*}


Now we are in a position to establish the estimate of $\mathbf{U}_{s}$ and $\mathbf{U}_{s+1}$. Combining the estimates of $J_{0}$,$J_{2}$,...,$J_{11}$ above with \eqref{UnE}, for any $t\in[0,t_{3}]$ we have
\begin{align*}
&\frac{1}{2}\frac{d}{dt}\|\mathbf{U}_{s}\|_{L_{\varphi}^{2}(\Omega)}^{2}
+\frac{3}{4}\Lambda\delta\|\sqrt{\y}\mathbf{U}_{s}\|_{L_{\varphi}^{2}(\Omega)}^{2}
+\frac{\lambda}{64}\|\sqrt{\omega_{\lambda}}\mathbf{U}_{s}\|_{L_{\varphi}^{2}(\Omega)}^{2}
+\|\partial_{y}\mathbf{U}_{s}\|_{L_{\varphi}^{2}(\Omega)}^{2}
+\|\sqrt{\eta}\partial_{x}\mathbf{U}_{s}\|_{L_{\varphi}^{2}(\Omega)}^{2}\\
\leq&~C_{*}C_{\tau}\iota^{-\frac{5}{2}}\|\mathbf{U}_{s}\|_{L_{\varphi}^{2}(\Omega)}^{2}
     +\frac{3}{8}\|\partial_{y}\mathbf{U}_{s}\|_{L_{\varphi}^{2}(\Omega)}^{2}
     +\left[C_{*}C_{\tau}(\sqrt{t}+\sqrt{\bar{\eta}}t^{\frac{1}{4}})+C\right]\|\mathcal{P}\mathbf{U}_{s}\|_{L_{\varphi}^{2}(\Omega)}^{2}\\
    &+\left(\frac{3}{8}+C_{\lambda}\theta\sqrt{t}\right)\|\sqrt{\eta}\partial_{x}\mathbf{U}_{s}\|_{L_{\varphi}^{2}(\Omega)}^{2}
     +\frac{1}{4}\|\sqrt{\omega_{\lambda}}\mathbf{U}_{s}\|_{L_{\varphi}^{2}(\Omega)}^{2}
     +(7\tau+C_{*}\sqrt{\iota})\|\sqrt{\tilde{\eta}}\partial_{x}\tilde{\mathbf{U}}_{s}\|_{L_{\tilde{\varphi}}^{2}(\Omega)}^{2}\\
    &+3\tau\|\tilde{\mathcal{P}}\tilde{\mathbf{U}}_{s}\|_{L_{\tilde{\varphi}}^{2}(\Omega)}^{2}
     +C_{*}C_{\tau}\|\tilde{\mathbf{U}}_{s}\|_{L_{\tilde{\varphi}}^{2}(\Omega)}^{2}
     +\tau\|\partial_{y}\mathcal{U}_{s}\|_{L_{\tilde{\varphi}}^{2}(\Omega)}^{2}
     +\tau\|\tilde{\mathcal{P}}\mathcal{U}_{s}\|_{L_{\tilde{\varphi}}^{2}(\Omega)}^{2}
     +C_{*}C_{\tau}\|\sqrt{\y}\mathbf{U}_{s}\|_{L_{\varphi}^{2}(\Omega)}^{2}\\
    &+\tau\|\sqrt{\y}\mathbb{U}\|_{\hat{H}_{\psi}^{s}(\Omega)}^{2}
     +C\|u\|_{\mathcal{H}_{\psi,\tilde{\varphi}}^{s}(\Omega)}^{2}
     +\frac{C_{*}}{(\bar{\varsigma}+t)^{\frac{1}{2}+\varepsilon_{0}}}\|\partial_{x}^{s+1}P\|_{L^{2}(\mathbb{T})}^{2}
     +C_{*}\left(1+\frac{1}{(\bar{\varsigma}+t)^{\frac{1}{2}+\varepsilon_{0}}}\right),
\end{align*}
and
\begin{align*}
&\frac{1}{2}\frac{d}{dt}\|\mathbf{U}_{s+1}\|_{L_{\hat{\varphi}}^{2}(\Omega)}^{2}
+\frac{3}{4}\Lambda\delta\|\sqrt{\y}\mathbf{U}_{s+1}\|_{L_{\hat{\varphi}}^{2}(\Omega)}^{2}
+\frac{\lambda}{64}\|\sqrt{\omega_{\lambda}}\mathbf{U}_{s+1}\|_{L_{\hat{\varphi}}^{2}(\Omega)}^{2}
+\|\partial_{y}\mathbf{U}_{s+1}\|_{L_{\hat{\varphi}}^{2}(\Omega)}^{2}
+\|\sqrt{\eta}\partial_{x}\mathbf{U}_{s+1}\|_{L_{\hat{\varphi}}^{2}(\Omega)}^{2}\\
\leq&~\frac{7}{16}\|\partial_{y}\mathbf{U}_{s+1}\|_{L_{\hat{\varphi}}^{2}(\Omega)}^{2}
     +\left(\frac{3}{8}+C_{\lambda}\theta\sqrt{t}\right)\|\sqrt{\eta}\partial_{x}\mathbf{U}_{s+1}\|_{L_{\hat{\varphi}}^{2}(\Omega)}^{2}
     +\frac{1}{4}\|\sqrt{\omega_{\lambda}}\mathbf{U}_{s+1}\|_{L_{\hat{\varphi}}^{2}(\Omega)}^{2}
     +C_{*}C_{\tau}\iota^{-\frac{5}{2}}\|\mathbf{U}_{s+1}\|_{L_{\hat{\varphi}}^{2}(\Omega)}^{2}\\
    &+\left[C_{*}C_{\tau}(\sqrt{t}+\sqrt{\bar{\eta}}t^{\frac{1}{4}})+C\right]\|\mathcal{P}\mathbf{U}_{s+1}\|_{L_{\hat{\varphi}}^{2}(\Omega)}^{2}
     +C_{*}C_{\tau}\|\sqrt{\y}\mathbf{U}_{s+1}\|_{L_{\hat{\varphi}}^{2}(\Omega)}^{2}
     +C_{\lambda}\|\mathcal{P}\mathbf{U}_{s}\|_{L_{\varphi}^{2}(\Omega)}^{2}
     +C_{\lambda}\|\partial_{y}\mathbf{U}_{s}\|_{L_{\varphi}^{2}(\Omega)}^{2}\\
    &+C_{*}\|\mathbf{U}_{s}\|_{L_{\varphi}^{2}(\Omega)}^{2}
     +(7\tau+C_{*}\sqrt{\iota})\|\sqrt{\tilde{\eta}}\partial_{x}\tilde{\mathbf{U}}_{s+1}\|_{L_{\tilde{\hat{\varphi}}}^{2}(\Omega)}^{2}
     +3\tau\|\tilde{\mathcal{P}}\tilde{\mathbf{U}}_{s+1}\|_{L_{\tilde{\hat{\varphi}}}^{2}(\Omega)}^{2}
     +C_{*}C_{\tau}\|\tilde{\mathbf{U}}_{s+1}\|_{L_{\tilde{\hat{\varphi}}}^{2}(\Omega)}^{2}\\
    &+\tau\|\partial_{y}\mathcal{U}_{s}\|_{L_{\tilde{\varphi}}^{2}(\Omega)}^{2}
     +5\tau\|\sqrt{\eta}\partial_{x}\mathcal{U}_{s}\|_{L_{\tilde{\varphi}}^{2}(\Omega)}^{2}
     +2\tau\|\tilde{\mathcal{P}}\mathcal{U}_{s}\|_{L_{\tilde{\varphi}}^{2}(\Omega)}^{2}
     +C\|u\|_{\mathcal{H}_{\psi,\tilde{\varphi}}^{s}(\Omega)}^{2}
     +\tau\|\sqrt{\y}\mathbb{U}\|_{\hat{H}_{\psi}^{s}(\Omega)}^{2}\\
    &+\frac{C_{*}}{(\bar{\varsigma}+t)^{\frac{1}{2}+\varepsilon_{0}}}\|\partial_{x}^{s+2}P\|_{L_{x}^{2}(\mathbb{T})}^{2}
     +C_{*}\left(1+\frac{1}{(\bar{\varsigma}+t)^{\frac{1}{2}+\varepsilon_{0}}}\right).
\end{align*}
By taking $\lambda$ sufficiently large in the above inequalities, we obtain the following proposition.
\begin{prop}
Let $\lambda>\lambda_{*}, \Lambda>\Lambda_{*}, \theta>\theta_{*}, \iota\in(0,\iota_{*}), \mathcal{Z}>\mathcal{Z}_{*}$, where $\lambda_{*}, \Lambda_{*},\iota_{*}, \mathcal{Z}_{*}$ are given in Theorem \ref{wpae}, and $\theta_{*}$ is given in Lemma \ref{wdotu}. Under the assumption of Theorem \ref{Phe}, if $\lambda$ is large enough, then the following inequality for any $t\in[0,t_{3}]$ and $\bar{\tau}\in(0,1)$,
\begin{align}\label{phe}
&\frac{1}{2}\frac{d}{dt}\left(\|\mathbf{U}_{s}\|_{L_{\varphi}^{2}(\Omega)}^{2}+\bar{\tau}\|\mathbf{U}_{s+1}\|_{L_{\hat{\varphi}}^{2}(\Omega)}^{2}\right)
 +\frac{3}{4}\Lambda\delta\left(\|\sqrt{\y}\mathbf{U}_{s}\|_{L_{\varphi}^{2}(\Omega)}^{2}+\bar{\tau}\|\sqrt{\y}\mathbf{U}_{s+1}\|_{L_{\hat{\varphi}}^{2}(\Omega)}^{2}\right)\\
&+\frac{\lambda}{128}\|\sqrt{\omega_{\lambda}}\mathbf{U}_{s}\|_{L_{\varphi}^{2}(\Omega)}^{2}
 +\frac{\lambda\bar{\tau}}{128}\|\sqrt{\omega_{\lambda}}\mathbf{U}_{s+1}\|_{L_{\hat{\varphi}}^{2}(\Omega)}^{2}
 +\frac{1}{2}\|\partial_{y}\mathbf{U}_{s}\|_{L_{\varphi}^{2}(\Omega)}^{2}
 +\frac{\bar{\tau}}{2}\|\partial_{y}\mathbf{U}_{s+1}\|_{L_{\hat{\varphi}}^{2}(\Omega)}^{2}\nonumber\\
&+\frac{5}{8}\|\sqrt{\eta}\partial_{x}\mathbf{U}_{s}\|_{L_{\varphi}^{2}(\Omega)}^{2}
 +\frac{5\bar{\tau}}{8}\|\sqrt{\eta}\partial_{x}\mathbf{U}_{s+1}\|_{L_{\hat{\varphi}}^{2}(\Omega)}^{2}\nonumber\\
\leq&~C_{*}C_{\tau}\iota^{-\frac{5}{2}}
      \left(\|\mathbf{U}_{s}\|_{L_{\varphi}^{2}(\Omega)}^{2}+\bar{\tau}\|\mathbf{U}_{s+1}\|_{L_{\hat{\varphi}}^{2}(\Omega)}^{2}\right)
      +C_{*}C_{\tau}\left(\|\sqrt{\y}\mathbf{U}_{s}\|_{L_{\varphi}^{2}(\Omega)}^{2}+\bar{\tau}\|\sqrt{\y}\mathbf{U}_{s+1}\|_{L_{\hat{\varphi}}^{2}(\Omega)}^{2}\right)\nonumber\\
    &+\left[C_{*}C_{\tau}(\sqrt{t}+\sqrt{\bar{\eta}}t^{\frac{1}{4}})+\check{C}\right]
      \left(\|\mathcal{P}\mathbf{U}_{s}\|_{L_{\varphi}^{2}(\Omega)}^{2}+\bar{\tau}\|\mathcal{P}\mathbf{U}_{s+1}\|_{L_{\hat{\varphi}}^{2}(\Omega)}^{2}\right)
     +\check{\mathcal{C}}_{\lambda}\bar{\tau}\|\mathcal{P}\mathbf{U}_{s}\|_{L_{\varphi}^{2}(\Omega)}^{2}\nonumber\\
    &+\check{\mathcal{C}}_{\lambda}\bar{\tau}\|\partial_{y}\mathbf{U}_{s}\|_{L_{\varphi}^{2}(\Omega)}^{2}
     +(7\tau+C_{*}\sqrt{\iota})\left(\|\sqrt{\tilde{\eta}}\partial_{x}\tilde{\mathbf{U}}_{s}\|_{L_{\tilde{\varphi}}^{2}(\Omega)}^{2}
                                     +\bar{\tau}\|\sqrt{\tilde{\eta}}\partial_{x}\tilde{\mathbf{U}}_{s+1}\|_{L_{\tilde{\hat{\varphi}}}^{2}(\Omega)}^{2}\right)\nonumber\\
    &+C_{\lambda}\theta\sqrt{t}\left(\|\sqrt{\eta}\partial_{x}\mathbf{U}_{s}\|_{L_{\varphi}^{2}(\Omega)}^{2}
                             +\bar{\tau}\|\sqrt{\eta}\partial_{x}\mathbf{U}_{s+1}\|_{L_{\hat{\varphi}}^{2}(\Omega)}^{2}\right)\nonumber\\
    &+3\tau\|\tilde{\mathcal{P}}\tilde{\mathbf{U}}_{s}\|_{L_{\tilde{\varphi}}^{2}(\Omega)}^{2}
     +3\tau\bar{\tau}\|\tilde{\mathcal{P}}\tilde{\mathbf{U}}_{s+1}\|_{L_{\tilde{\hat{\varphi}}}^{2}(\Omega)}^{2}
     +C_{*}C_{\tau}\|\tilde{\mathbf{U}}_{s}\|_{L_{\tilde{\varphi}}^{2}(\Omega)}^{2}
     +C_{*}C_{\tau}\bar{\tau}\|\tilde{\mathbf{U}}_{s+1}\|_{L_{\tilde{\hat{\varphi}}}^{2}(\Omega)}^{2}\nonumber\\
    &+2\tau\|\sqrt{\y}\mathbb{U}\|_{\hat{H}_{\psi}^{s}(\Omega)}^{2}
     +C\|u\|_{\mathcal{H}_{\psi,\tilde{\varphi}}^{s}(\Omega)}^{2}
     +\frac{C_{*}}{(\bar{\varsigma}+t)^{\frac{1}{2}+\varepsilon_{0}}}\|\partial_{x}^{s+1}P\|_{L_{x}^{2}(\mathbb{T})}^{2}
     +\frac{C_{*}}{(\bar{\varsigma}+t)^{\frac{1}{2}+\varepsilon_{0}}}\|\partial_{x}^{s+2}P\|_{L_{x}^{2}(\mathbb{T})}^{2}\nonumber\\
    &+2\tau\|\partial_{y}\mathcal{U}_{s}\|_{L_{\tilde{\varphi}}^{2}(\Omega)}^{2}
     +5\tau\|\sqrt{\eta}\partial_{x}\mathcal{U}_{s}\|_{L_{\tilde{\varphi}}^{2}(\Omega)}^{2}
     +3\tau\|\tilde{\mathcal{P}}\mathcal{U}_{s}\|_{L_{\tilde{\varphi}}^{2}(\Omega)}^{2}
     +C_{*}\left(1+\frac{1}{(\bar{\varsigma}+t)^{\frac{1}{2}+\varepsilon_{0}}}\right),\nonumber
\end{align}
where $\check{\mathcal{C}}_{\lambda}$ is a positive constant depending only on $\lambda$, and $\check{C}$ is a positive constant independent of $\lambda$ and $\mathcal{Z}$.
\end{prop}

To prove Theorem \ref{Phe}, let $\theta>\theta^{*}$ be fixed, and write
$\mathbb{S}(t)=\|\mathbf{U}_{s}\|_{L_{\varphi}^{2}(\Omega)}^{2}+\bar{\tau}\|\mathbf{U}_{s+1}\|_{L_{\hat{\varphi}}^{2}(\Omega)}^{2}$,
\begin{align*}
\mathbb{W}(t)
=&\Lambda\delta\|\sqrt{\y}\mathbf{U}_{s}\|_{L_{\varphi}^{2}(\Omega)}^{2}
 +\bar{\tau}\Lambda\delta\|\sqrt{\y}\mathbf{U}_{s+1}\|_{L_{\hat{\varphi}}^{2}(\Omega)}^{2}
 +\frac{\lambda}{128}\|\sqrt{\omega_{\lambda}}\mathbf{U}_{s}\|_{L_{\varphi}^{2}(\Omega)}^{2}
 +\frac{\lambda\bar{\tau}}{128}\|\sqrt{\omega_{\lambda}}\mathbf{U}_{s+1}\|_{L_{\hat{\varphi}}^{2}(\Omega)}^{2}\\
&+\|\partial_{y}\mathbf{U}_{s}\|_{L_{\varphi}^{2}(\Omega)}^{2}
 +\bar{\tau}\|\partial_{y}\mathbf{U}_{s+1}\|_{L_{\hat{\varphi}}^{2}(\Omega)}^{2}
 +\|\sqrt{\eta}\partial_{x}\mathbf{U}_{s}\|_{L_{\varphi}^{2}(\Omega)}^{2}
 +\bar{\tau}\|\sqrt{\eta}\partial_{x}\mathbf{U}_{s+1}\|_{L_{\hat{\varphi}}^{2}(\Omega)}^{2},
\end{align*}
and
\begin{align*}
\mathbb{G}(t)
=&~(7\tau+C_{*}\sqrt{\iota})\left(\|\sqrt{\tilde{\eta}}\partial_{x}\tilde{\mathbf{U}}_{s}\|_{L_{\tilde{\varphi}}^{2}(\Omega)}^{2}
                                     +\bar{\tau}\|\sqrt{\tilde{\eta}}\partial_{x}\tilde{\mathbf{U}}_{s+1}\|_{L_{\tilde{\hat{\varphi}}}^{2}(\Omega)}^{2}\right)\\
 &+3\tau\|\tilde{\mathcal{P}}\tilde{\mathbf{U}}_{s}\|_{L_{\tilde{\varphi}}^{2}(\Omega)}^{2}
  +3\tau\bar{\tau}\|\tilde{\mathcal{P}}\tilde{\mathbf{U}}_{s+1}\|_{L_{\tilde{\hat{\varphi}}}^{2}(\Omega)}^{2}
  +C_{*}C_{\tau}\|\tilde{\mathbf{U}}_{s}\|_{L_{\tilde{\varphi}}^{2}(\Omega)}^{2}
  +C_{*}C_{\tau}\bar{\tau}\|\tilde{\mathbf{U}}_{s+1}\|_{L_{\tilde{\hat{\varphi}}}^{2}(\Omega)}^{2}\nonumber\\
 &+2\tau\|\sqrt{\y}\mathbb{U}\|_{\hat{H}_{\psi}^{s}(\Omega)}^{2}
  +C\|u\|_{\mathcal{H}_{\psi,\tilde{\varphi}}^{s}(\Omega)}^{2}
  +\frac{C_{*}}{(\bar{\varsigma}+t)^{\frac{1}{2}+2\varepsilon_{0}}}\|\partial_{x}^{s+1}P\|_{L_{x}^{2}(\mathbb{T})}^{2}
  +\frac{C_{*}}{(\bar{\varsigma}+t)^{\frac{1}{2}+2\varepsilon_{0}}}\|\partial_{x}^{s+2}P\|_{L_{x}^{2}(\mathbb{T})}^{2}\nonumber\\
 &+2\tau\|\partial_{y}\mathcal{U}_{s}\|_{L_{\tilde{\varphi}}^{2}(\Omega)}^{2}
  +5\tau\|\sqrt{\eta}\partial_{x}\mathcal{U}_{s}\|_{L_{\tilde{\varphi}}^{2}(\Omega)}^{2}
  +3\tau\|\tilde{\mathcal{P}}\mathcal{U}_{s}\|_{L_{\tilde{\varphi}}^{2}(\Omega)}^{2}
  +C_{*}\left(1+\frac{1}{(\bar{\varsigma}+t)^{\frac{1}{2}+\varepsilon_{0}}}\right).\nonumber
\end{align*}

By taking $\lambda$ large enough, one has
\begin{align*}
(\check{C}+2)\|\mathcal{P}\mathbf{U}_{s}\|_{L_{\varphi}^{2}(\Omega)}^{2}
\leq \frac{\lambda}{256}\|\sqrt{\omega_{\lambda}}\mathbf{U}_{s}\|_{L_{\varphi}^{2}(\Omega)}^{2}
    +C\|\mathbf{U}_{s}\|_{L_{\varphi}^{2}(\Omega)}^{2},~\forall~t\in[0,t_{3}],
\end{align*}
and
\begin{align}\label{pus}
\|\mathcal{P}\mathbf{U}_{s}\|_{L_{\varphi}^{2}(\Omega)}^{2}
\leq \frac{\lambda}{256}\|\sqrt{\omega_{\lambda}}\mathbf{U}_{s}\|_{L_{\varphi}^{2}(\Omega)}^{2}
    +\|\mathbf{U}_{s}\|_{L_{\varphi}^{2}(\Omega)}^{2},   ~\forall~t\in[0,t_{3}],
\end{align}
moreover,
\begin{align*}
(\check{C}+2)\|\mathcal{P}\mathbf{U}_{s+1}\|_{L_{\hat{\varphi}}^{2}(\Omega)}^{2}
\leq \frac{\lambda}{256}\|\sqrt{\omega_{\lambda}}\mathbf{U}_{s+1}\|_{L_{\hat{\varphi}}^{2}(\Omega)}^{2}
    +C\|\mathbf{U}_{s+1}\|_{L_{\hat{\varphi}}^{2}(\Omega)}^{2},~\forall~t\in[0,t_{3}].
\end{align*}
and
\begin{align}\label{pus+1}
\|\mathcal{P}\mathbf{U}_{s+1}\|_{L_{\hat{\varphi}}^{2}(\Omega)}^{2}
\leq \frac{\lambda}{256}\|\sqrt{\omega_{\lambda}}\mathbf{U}_{s+1}\|_{L_{\hat{\varphi}}^{2}(\Omega)}^{2}
    +\|\mathbf{U}_{s+1}\|_{L_{\hat{\varphi}}^{2}(\Omega)}^{2},~\forall~t\in[0,t_{3}].
\end{align}
With $\lambda$ fixed, we take $\bar{\tau}$ small enough such that $\check{\mathcal{C}}_{\lambda}\bar{\tau}\leq1$, then under the assumption of Theorem \ref{Phe}, since $|\rho(0,x,y)|\leq C|\rho_{in}(x,y)|\leq C\mathcal{P}_{in}$ in $\Omega$, with such fixed $\lambda$ and $\bar{\tau}$ we can take $\mathcal{Z}$ large enough to obtain
\begin{align}\label{ine}
\mathbb{S}(0)
\leq&~C_{\lambda}\|\mathbf{U}_{s}(0,\cdot)\|_{L_{\hat{\psi}_{in}}^{2}(\Omega)}^{2}
     +C_{\lambda}\bar{\tau}\|\mathbf{U}_{s+1}(0,\cdot)\|_{L_{\check{\psi}_{in}}^{2}(\Omega)}^{2}\\
\leq&~C_{\lambda}\|\mathcal{P}_{in}\partial_{y}\partial_{x}^{s}\tilde{u}_{in}\|_{L_{\hat{\psi}_{in}}^{2}(\Omega)}^{2}
     +C_{\lambda}\|\mathcal{P}_{in}\rho(0,x,y)\partial_{x}^{s}\tilde{u}_{in}\|_{L_{\hat{\psi}_{in}}^{2}(\Omega)}^{2}\nonumber\\
    &+C_{\lambda}\bar{\tau}\|\mathcal{P}_{in}\partial_{y}\partial_{x}^{s+1}\tilde{u}_{in}\|_{L_{\check{\psi}_{in}}^{2}(\Omega)}^{2}
     +C_{\lambda}\bar{\tau}\|\mathcal{P}_{in}\rho(0,x,y)\partial_{x}^{s+1}\tilde{u}_{in}\|_{L_{\check{\psi}_{in}}^{2}(\Omega)}^{2}\nonumber\\
\leq&~C_{\lambda}\|\mathcal{P}_{in}\partial_{y}\partial_{x}^{s}\tilde{u}_{in}\|_{L_{\hat{\psi}_{in}}^{2}(\Omega)}^{2}
     +C_{\lambda}\|\mathcal{P}_{in}^{2}\partial_{x}^{s}\tilde{u}_{in}\|_{L_{\hat{\psi}_{in}}^{2}(\Omega)}^{2}\nonumber\\
    &+C_{\lambda}\bar{\tau}\|\mathcal{P}_{in}\partial_{y}\partial_{x}^{s+1}\tilde{u}_{in}\|_{L_{\check{\psi}_{in}}^{2}(\Omega)}^{2}
     +C_{\lambda}\bar{\tau}\|\mathcal{P}_{in}^{2}\partial_{x}^{s+1}\tilde{u}_{in}\|_{L_{\check{\psi}_{in}}^{2}(\Omega)}^{2}\nonumber\\
\leq&~\frac{\bar{\tau}\mathcal{Z}^{2}}{64}.\nonumber
\end{align}
Moreover, we can take $\tilde{t}\in(0,t_{3}]$ and $\tau, \iota$ small enough such that
\begin{align}\label{G}
\int_{0}^{\tilde{t}}\mathbb{G}(\sigma)d\sigma\leq \frac{\bar{\tau}\mathcal{Z}^{2}}{128}.
\end{align}
For such fixed $\lambda,\tau$ and $\mathcal{Z}$, we take $\bar{t}_{*}\in[0,\tilde{t}]$ small enough to obtain
$$C_{*}C_{\tau}(\sqrt{t}+\sqrt{\bar{\eta}}t^{\frac{1}{4}})\leq 1,~\forall~t\in[0,\bar{t}_{*}],$$
With $\lambda,\tau,\bar{\tau},\theta,\iota$ and $\mathcal{Z}$ fixed, by taking $\bar{t}^{*}\in [0,\bar{t}_{*}]$ small enough such that $C_{\lambda}\theta\sqrt{t}\leq \frac{1}{8}$ for any $t\in[0,\bar{t}^{*}]$ and  $\Lambda$ large enough such that $C_{*}C_{\tau}\leq \frac{\Lambda\delta}{4}$ in \eqref{phe}, one can deduce that for any $t\in[0,\bar{t}^{*}]$,
\begin{align}\label{phe0}
&\frac{1}{2}\frac{d}{dt}\left(\|\mathbf{U}_{s}\|_{L_{\varphi}^{2}(\Omega)}^{2}+\bar{\tau}\|\mathbf{U}_{s+1}\|_{L_{\hat{\varphi}}^{2}(\Omega)}^{2}\right)
 +\frac{\Lambda\delta}{2}\left(\|\sqrt{\y}\mathbf{U}_{s}\|_{L_{\varphi}^{2}(\Omega)}^{2}+\bar{\tau}\|\sqrt{\y}\mathbf{U}_{s+1}\|_{L_{\hat{\varphi}}^{2}(\Omega)}^{2}\right)\\
&+\frac{\lambda}{256}\|\sqrt{\omega_{\lambda}}\mathbf{U}_{s}\|_{L_{\varphi}^{2}(\Omega)}^{2}
 +\frac{\lambda\bar{\tau}}{256}\|\sqrt{\omega_{\lambda}}\mathbf{U}_{s+1}\|_{L_{\hat{\varphi}}^{2}(\Omega)}^{2}
 +\frac{1}{2}\|\partial_{y}\mathbf{U}_{s}\|_{L_{\varphi}^{2}(\Omega)}^{2}
 +\frac{\bar{\tau}}{2}\|\partial_{y}\mathbf{U}_{s+1}\|_{L_{\hat{\varphi}}^{2}(\Omega)}^{2}\nonumber\\
&+\frac{1}{2}\|\sqrt{\eta}\partial_{x}\mathbf{U}_{s}\|_{L_{\varphi}^{2}(\Omega)}^{2}
 +\frac{\bar{\tau}}{2}\|\sqrt{\eta}\partial_{x}\mathbf{U}_{s+1}\|_{L_{\hat{\varphi}}^{2}(\Omega)}^{2}\nonumber\\
\leq&~C_{*}C_{\tau}\iota^{-\frac{5}{2}}
      \left(\|\mathbf{U}_{s}\|_{L_{\varphi}^{2}(\Omega)}^{2}+\bar{\tau}\|\mathbf{U}_{s+1}\|_{L_{\hat{\varphi}}^{2}(\Omega)}^{2}\right)
    +\mathbb{G}(t),\nonumber
\end{align}
which gives
\begin{align}\label{phe1}
\frac{d}{dt}\mathbb{S}(t)+\mathbb{W}(t)
\leq C_{*}C_{\tau}\iota^{-\frac{5}{2}}\mathbb{S}(t)+2\mathbb{G}(t).
\end{align}
Utilizing the Gronwall inequality, from \eqref{phe1} one can deduce that for any $t\in[0,\bar{t}^{*}]$,
\begin{align}\label{phe3}
\mathbb{S}(t)+\int_{0}^{t}\mathbb{W}(\sigma)d\sigma
\leq& C_{*}C_{\tau}\iota^{-\frac{5}{2}}\int_{0}^{t}\mathbb{S}(\sigma)d\sigma
     +\int_{0}^{t}2\mathbb{G}(\sigma)d\sigma
     +\mathbb{S}(0)\\
\leq& C_{*}C_{\tau}\iota^{-\frac{5}{2}}\int_{0}^{t}e^{C_{*}C_{\tau}\iota^{-\frac{5}{2}}\sigma}d\sigma
           \left(\int_{0}^{t}2\mathbb{G}(\sigma)d\sigma+\mathbb{S}(0)\right)\nonumber\\
    &+\int_{0}^{t}2\mathbb{G}(\sigma)d\sigma
     +\mathbb{S}(0).\nonumber
\end{align}
With \eqref{G} and \eqref{ine}, from \eqref{phe3} one can choose $\hat{t}\in (0,\bar{t}^{*})$ small enough such that for any $t\in[0,\hat{t}]$, there holds
\begin{align}\label{phe4}
\mathbb{S}(t)+\int_{0}^{t}\mathbb{W}(\sigma)d\sigma\leq \frac{\bar{\tau}\mathcal{Z}^{2}}{16},
\end{align}
Moreover, from \eqref{pus} and \eqref{pus+1} one can choose $t_{*}\in(0,\hat{t}]$ such that
\begin{align}\label{Pus}
\|\mathcal{P}\mathbf{U}_{s}\|_{L_{t}^{2}L_{\varphi}^{2}(\Omega_{t_{*}})}^{2}
\leq \frac{\lambda}{256}\|\sqrt{\omega_{\lambda}}\mathbf{U}_{s}\|_{L_{t}^{2}L_{\varphi}^{2}(\Omega_{t_{*}})}^{2}
    +\|\mathbf{U}_{s}\|_{L_{t}^{\infty}L_{\varphi}^{2}(\Omega_{t_{*}})}^{2},
\end{align}
and
\begin{align}\label{Pus+1}
\|\mathcal{P}\mathbf{U}_{s+1}\|_{L_{t}^{2}L_{\hat{\varphi}}^{2}(\Omega_{t_{*}})}^{2}
\leq \frac{\lambda}{256}\|\sqrt{\omega_{\lambda}}\mathbf{U}_{s+1}\|_{L_{t}^{2}L_{\hat{\varphi}}^{2}(\Omega_{t_{*}})}^{2}
    +\|\mathbf{U}_{s+1}\|_{L_{t}^{\infty}L_{\hat{\varphi}}^{2}(\Omega_{t_{*}})}^{2}.
\end{align}
Recall the definition of $\mathbb{S}$ and $\mathbb{W}$, from \eqref{Pus},\eqref{Pus+1} and \eqref{phe4} we get \eqref{Hoe1} and \eqref{Hoe2}.
\end{proof}

\end{document}
